\documentclass[12pt, a4paper]{article}

%% ════════════════════════════════════════════════════════════════════════
%% CRF_Complete_v17_16_Part_G.tex
%%          (Communication-Layer retrofit of v17.9 Part G)
%%
%% Part G of the v17.16 multi-part single volume edition.
%% Contains Parts XXXV-XXXVII:
%%   Part XXXV  : ZC45 — Universal Floor & δ Two-Mode Bridge
%%   Part XXXVI : ZC46-49 Cluster — Z_15 Projection Chain & Four-Axis
%%   Part XXXVII: Cross-Part Connection Map v17.9
%%
%% Multi-Part Architecture v17.9 (content unchanged at retrofit):
%%   Part A (Parts I-VIII):       Foundations & Early Dynamics
%%   Part B (Parts IX-XIII):      Algebraic Core & Methodology
%%   Part C (Parts XIV-XXI):      Methodology & Z-Programme through ZC19
%%   Part D (Parts XXIII-XXVI):   Lie Bridge, alpha_EM, R_1, Force Hier.
%%   Part E (Parts XXVII-XXX):    Cluster Closures + TLS Layer +
%%                                Cross-Part Connection Map v17.7
%%   Part F (Parts XXXI-XXXIV):   ZC41-44 + Mind/Barami + Axiom 11 +
%%                                Cross-Part Connection Map v17.8
%%   Part G (Parts XXXV-XXXVII):  ZC45-49 + Cross-Part Map v17.9 <- THIS
%%
%% Governance: CUIP v1.1, CRF Style Guide v3.6, Gauge Fix Rule v1.0
%% Compile: xelatex (required for Pali diacritics + Thai script)
%%          3 passes required for TOC + refs
%%
%% Part counter: Part F ends at Part XXXIV -> Part G begins at Part XXXV
%%
%% ── COMMUNICATION LAYER (v3.6 §31; Protocol v1.3 §U) ────────────────────
%% Primary audience tier: 1 (Practitioner)  |  On-ramp for tier below (0): yes
%% Communication Layer retrofit — fifth Part after the K pilot and the
%% J / I / H arc (v17.16). Added additively (No-Loss), nothing overridden:
%%     • Reader's On-Ramp present after TOC (keybox, ≤12 lines, no atomic IDs);
%%     • Bridge Prose blocks at every ZC embed seam (ZC45–49), each marked
%%       % [BRIDGE-PROSE], WHY-test passed (state why open / why surprising /
%%       what it enables — not "did X then Y");
%%     • \physmeaning applied to one major theorem box whose box carried no
%%       existing physical-reading prose (THM-ZC45-01, the Universal Floor).
%%
%% Macro patch: \input CRF_v17_15_macro_patch.tex (repointed from v17_9 at
%%   retrofit) = full Parts A–J macro set + \physmeaning; chains the
%%   consolidated v17_14 patch. Verify 0 undefined control sequences.
%% ════════════════════════════════════════════════════════════════════════

%% ── PACKAGES (matching Part F v17.8 verbatim) ───────────────────────────
\usepackage{amsmath, amssymb, amsthm}
\usepackage{mathtools}
\usepackage{geometry}
\usepackage{xcolor}
\usepackage{parskip}
\usepackage[protrusion=false,expansion=false]{microtype}
\usepackage{hyperref}
\usepackage{tcolorbox}
\usepackage{booktabs}
\usepackage{longtable}
\usepackage{array}
\usepackage[T1]{fontenc}
\usepackage{graphicx}
\usepackage{enumitem}
\usepackage{needspace}
\usepackage{tocloft}

%% XeLaTeX font support (matches Part F v17.8)
\usepackage{iftex}
\ifxetex
  \usepackage{fontspec}
  \setmainfont{Latin Modern Roman}
  \usepackage{polyglossia}
  \setdefaultlanguage{english}
  \setotherlanguage{thai}
  \newfontfamily\thaifont[Scale=0.95]{Loma}
\fi

\geometry{margin=2.5cm}

%% TOC spacing (preserved from v17.8)
\setlength{\cftbeforesecskip}{0pt}
\setlength{\cftbeforesubsecskip}{0pt}
\setlength{\cftbeforepartskip}{6pt}
\setlength{\cftsubsecindent}{2.5em}
\setlength{\cftsubsubsecindent}{4.5em}
\renewcommand{\cftsecfont}{\normalfont}
\renewcommand{\cftsubsecfont}{\small}
\renewcommand{\cftsecpagefont}{\normalfont}
\renewcommand{\cftsubsecpagefont}{\small}

%% ── COLOR PALETTE (preserved from Part F) ────────────────────────────────
\definecolor{deepblue}{RGB}{10,30,80}
\definecolor{crfgreen}{RGB}{0,100,60}
\definecolor{softgreen}{RGB}{180,230,210}
\definecolor{physgold}{RGB}{140,100,0}
\definecolor{softgold}{RGB}{255,240,180}
\definecolor{loopred}{RGB}{160,30,30}
\definecolor{softred}{RGB}{255,210,210}
\definecolor{mutedgray}{RGB}{90,90,90}
\definecolor{midblue}{RGB}{30,80,160}
\definecolor{softblue}{RGB}{180,210,240}

\hypersetup{colorlinks=true, linkcolor=deepblue, urlcolor=midblue,
            citecolor=deepblue, bookmarks=false}

%% ── BOX STYLES (matching Part F) ────────────────────────────────────────
\tcbuselibrary{skins, breakable}

\newtcolorbox{derivedbox}[1][]{
  colback=softgreen!50, colframe=crfgreen!60,
  fonttitle=\bfseries\color{crfgreen!20!black},
  title={Derived}, breakable, #1}

\newtcolorbox{formalbox}[1][]{
  colback=softblue!40, colframe=midblue!60,
  fonttitle=\bfseries\color{deepblue},
  title={Formal}, breakable, #1}

\newtcolorbox{openqbox}[1][]{
  colback=softred!30, colframe=loopred!50,
  fonttitle=\bfseries\color{loopred!20!black},
  title={Open}, breakable, #1}

\newtcolorbox{keybox}[1][]{
  colback=softgold!50, colframe=physgold!60,
  fonttitle=\bfseries\color{physgold!20!black},
  breakable, #1}

%% Cross-Part Connection Map boxes (carried from Part F v17.8)
\definecolor{conmapbg}{RGB}{240,245,255}
\definecolor{conmapframe}{RGB}{50,80,140}
\definecolor{inheritbg}{RGB}{225,240,255}
\definecolor{inheritframe}{RGB}{20,60,160}
\definecolor{correctbg}{RGB}{255,235,220}
\definecolor{correctframe}{RGB}{180,90,30}

\newtcolorbox{conmapbox}[1][]{
  colback=conmapbg, colframe=conmapframe,
  fonttitle=\bfseries\color{white},
  breakable, #1}

\newtcolorbox{inheritbox}[1][]{
  colback=inheritbg, colframe=inheritframe,
  fonttitle=\bfseries\color{white},
  title={Backward Inheritance}, breakable, #1}

\newtcolorbox{correctionbox}[1][]{
  colback=correctbg, colframe=correctframe,
  fonttitle=\bfseries\color{white},
  title={Forward Correction / Cross-Part PM}, breakable, #1}

%% ── THEOREM ENVIRONMENTS ────────────────────────────────────────────────
\newtheorem{theorem}{Theorem}[section]
\newtheorem{lemma}{Lemma}[section]
\newtheorem{principle}{Principle}[section]
\newtheorem{claim}{Claim}[section]
\newtheorem{definition}{Definition}[section]
\newtheorem{proposition}{Proposition}[section]
\newtheorem{corollary}{Corollary}[section]
\newtheorem{remark}{Remark}[section]
\newtheorem{conjecture}{Conjecture}[section]
\newtheorem{finding}{Finding}[section]
\newtheorem{hypothesis}{Hypothesis}[section]
\newtheorem*{remarkstar}{Remark}

%% ── TITLE ────────────────────────────────────────────────────────────────
\title{%
  \textbf{\color{deepblue}Constrained Reality Framework (CRF)}\\[0.4em]
  \large Complete Edition v17.9 --- Part G\\
  \large Universal Floor $\cdot$ $\mathbb{Z}_{15}$ Projection Chain
    $\cdot$ Four-Axis Duality\\[0.3em]
  \normalsize\color{mutedgray}
  Parts XXXV--XXXVII $\cdot$\\
  $\varepsilon_0$ Two-Mode Bridge $\cdot$
  Three-Projection Theorem $\cdot$
  CRT Sector Decoupling $\cdot$
  Lock-Ratio Exact $\cdot$
  Four-Axis Duality
}
\author{Ougit Jarittum\\
\small Independent Researcher, Thailand\quad
ORCID: 0009-0009-4451-6811\\
\small Zenodo: \href{https://doi.org/10.5281/zenodo.18977059}%
{10.5281/zenodo.18977059}\quad (CC-BY-NC 4.0)
}
\date{May 2026 --- Complete Edition v17.9 (Part G of 7)}

%% ── PART COUNTER: continue from Part F ──────────────────────────────────
%% Part F ends at Part XXXIV -> Part G begins at Part XXXV
\setcounter{part}{34}

%% ── MACRO PATCH v17.9 ───────────────────────────────────────────────────
\input{CRF_v17_15_macro_patch.tex}

%% ════════════════════════════════════════════════════════════════════════
\begin{document}

% Governance: CUIP v1.1, CRF Style Guide v3.2, Gauge Fix Rule v1.0
% Gauge: T-canonical (d_s = 3 exact, n_m = 5 exact)

\maketitle

\begin{abstract}
\sloppy
This is \textbf{Part~G} of the CRF Complete Edition v17.9, the seventh
of seven sequential files comprising one continuous volume on the
$\delta$-mechanism. Parts~A through F (separately compiled) cover
foundations through the ZC41--44 cluster (Parts~I--XXXIV).

\textbf{Part~G} integrates five post-v17.8 standalone papers organised
into three new parts:

\begin{itemize}[leftmargin=2em]
\item \textbf{Part XXXV (ZC45 --- Universal Floor and $\delta$ Two-Mode
  Bridge).}
ZC45 closes GAP-GJ-01 by proving the Universal Floor Theorem
(THM-ZC45-01): the instability floor $\varepsilon_0$ is identical
in both CRF operating modes ---
$\varepsilon_{0,\mathrm{Spont}} = \varepsilon_{0,\mathrm{Agent}} = \varepsilon_0$.
Two lemmas are required. LEM-ZC45-01 establishes that $\varepsilon_0$
is a $\delta$-chain property deriving exclusively from Axioms~0--3,
without reference to mode. LEM-ZC45-02 shows that both Spontaneous Mode
(DEF-SA01) and Agentive Mode (DEF-MB02) share the same $\delta$-chain,
hence the same $\varepsilon_0$. COR-ZC45-01 promotes FIND-GJ-01 from
Finding to Theorem. PM-ZC45-01 records the 0.16\% Wald residual
(GAP-ZC42-02, inherited, not closed).

\item \textbf{Part XXXVI (ZC46--49 Cluster: $\mathbb{Z}_{15}$ Projection
  Chain and Four-Axis Duality).}
ZC46 promotes CONJ-ZC44-01 to the Three-Projection Theorem
(THM-ZC46-01): $\mathcal{R} = 15/8$ is the common root of ZC18-B,
ZC42, and ZC43 as three projections of $\mathbb{Z}_{15}$ onto different
axes. GAP-ZC46-01 (formal SO(3) decoupling) is registered.
ZC47 closes GAP-ZC46-01 via the CRT Sector Decoupling Theorem
(THM-ZC47-01), using the tensor decomposition
$\ell^2(\mathbb{Z}_{15}) \cong \ell^2(\mathbb{Z}_3) \otimes
\ell^2(\mathbb{Z}_5)$. THM-ZC46-01 becomes unconditional (COR-ZC47-02).
CONJ-ZC47-01 (lock-time ratio candidate) is registered.
ZC48 proves CONJ-ZC47-01 exactly without mean-field approximation,
via LEM-ZC48-01 (CRF Model~A: $\tau_i \propto 1/\alpha_i^{R_0}$)
and THM-ZC29-01 (Force Hierarchy). Lock-Ratio Exact Theorem
(THM-ZC48-01): $\tau_{\mathbb{Z}_3}/\tau_{\mathbb{Z}_5} = \varphi(5)/\varphi(3) = 2$
exact. FIND-ZC48-01--03 record sector ladder, gap resolution, totient
universality.
ZC49 synthesises the four axes of $\mathbb{Z}_{15}$ sector structure
into the Four-Axis Duality Theorem (THM-ZC49-01). FIND-ZC49-01--04
record the two-class product invariant, geometric mean identity, perfect
square identity, and Key Identity in totient disguise.

\item \textbf{Part XXXVII (Cross-Part Connection Map v17.9).}
Extends Part~XXXIV (v17.8 Connection Map) with backward inheritance,
forward correction, and cross-Part Phase Misalignment records for
ZC45--ZC49. Includes Master Z-Programme Status Table v17.9 delta and
\texttt{\textbackslash Gcrf} macro conflict notice (PM-ZC48-MACRO).
\end{itemize}

\textbf{Key closures in this Part:}
GAP-GJ-01 (ZC45, THM-ZC45-01),
GAP-ZC46-01 (ZC47, COR-ZC47-01).
CONJ-ZC44-01 promoted to THM-ZC46-01 (ZC46).
CONJ-ZC47-01 promoted to THM-ZC48-01 (ZC48).
THM-ZC46-01 becomes unconditional (ZC47, COR-ZC47-02).

\textbf{Phase Misalignments preserved:}
PM-ZC45-01 (0.16\% Wald residual, GAP-ZC42-02 inherited),
PM-ZC46-01 (probe P1 incomplete condition),
PM-ZC47-01 (two probe specification scars),
PM-ZC48-01 (PC/PG probe specification errors),
PM-ZC48-MACRO (\texttt{\textbackslash Gcrf} convention conflict ZC48--49
vs.\ v17.6 patch).

\textbf{Multi-Part Single Volume.}
Parts~A through G share one conceptual atomic-numbering namespace.
Combine via:\\
\texttt{pdfunite Part\_A.pdf Part\_B.pdf Part\_C.pdf Part\_D.pdf
Part\_E.pdf Part\_F.pdf Part\_G.pdf CRF\_Complete\_v17\_9.pdf}
\end{abstract}

\tableofcontents
\newpage

% [ON-RAMP: integration-added, Communication Layer v3.6 §31.2 / Protocol v1.3 §U2]
\begin{keybox}[title={If you are just arriving --- Part G in plain language}]
\sloppy
CRF attempts to derive the constants and force structure of physics from a
single primitive called Distinction ($\delta$). Two of its deepest claims are
that the same tiny ``instability floor'' $\varepsilon_0$ (about $0.0075$)
governs both inert matter and deciding minds, and that one small cyclic group,
$\mathbb{Z}_{15}$, is the algebraic seed from which the force structure unfolds.
\textbf{Part G} pins both down. It first proves that the floor really is a
single universal object seen from two sides, not two numbers that happen to
agree; it then takes a previously conjectured ``bridge ratio'' $\mathcal{R}=15/8$,
promotes it to a theorem, closes the one algebraic gap that promotion left open
(by showing the group cleanly splits into independent pieces), proves the
resulting lock-ratio is exactly $2$, and finally gathers the strands into one
four-axis picture.

\textbf{Minimum vocabulary:}
\textit{$\delta$ (Distinction)} --- the one starting idea everything is built
from;\;
\textit{$\varepsilon_0$ (instability floor)} --- the small threshold this Part
shows is shared by matter and mind;\;
\textit{Spontaneous vs.\ Agentive Mode} --- $\delta$ firing automatically in
matter versus compelling a choice in mind;\;
\textit{$\mathbb{Z}_{15}$} --- the cyclic group of order $15=3\times5$ that
seeds the force structure;\;
\textit{$\mathcal{R}=15/8$} --- a ``bridge ratio'' linking the group's size to
its mode count;\;
\textit{gap / GAP-\dots} --- an open question a later paper aims to close.
Full symbol list: Master Notation Key (front matter).

These results are pieces of a map under construction, not a claim of final
truth: each carries the conditions under which it would fail (see front
matter, ``Map, not reality'').
\end{keybox}

\clearpage

\part{ZC45: Universal Floor and $\delta$ Two-Mode Bridge}
\label{part:xxxv}

\begin{keybox}[title={Part XXXV Overview}]
\begin{itemize}[leftmargin=1.5em]
\item \textbf{ZC45} (Universal Floor): closes GAP-GJ-01;
  proves $\varepsilon_{0,\mathrm{Spont}} = \varepsilon_{0,\mathrm{Agent}}
  = \varepsilon_0$ (THM-ZC45-01); promotes FIND-GJ-01 to Theorem
  (COR-ZC45-01); PM-ZC45-01 scar (0.16\% Wald residual).
\item \textbf{Physics--Mind bridge:} First paper to formally unify
  the $\varepsilon_0$ floor across both CRF operating modes ---
  matter (Spontaneous) and mind (Agentive) share the same irreducible
  $D_{\mathrm{KL}}$ floor from Axioms~0--3.
\end{itemize}
\end{keybox}

%% ── ZC45 ─────────────────────────────────────────────────────────────────
\section{ZC45: Universal Floor Theorem --- $\varepsilon_0$ Two-Mode
  Identity}
\addcontentsline{toc}{section}{ZC45: Universal Floor}
\label{sec:zc45}

% [BRIDGE-PROSE: integration-added, Extension Freedom narrative, v3.6 §31.6]
The framework had long carried the same small number, $\varepsilon_0$, in two
places that have no obvious reason to meet: the irreducible floor below which
inert matter cannot resolve a distinction, and the residual tension a mind
still feels at the threshold of a fully released state. That the two readings
landed on the same value looked, until now, like the kind of numerical
coincidence a sceptic is right to distrust --- two independent chains, one
shared digit. The stakes of ZC45 are exactly that distrust: if the agreement
were accidental, the claim that CRF derives mind and matter from one mechanism
would be decoration; if it is forced, the claim is load-bearing. What makes the
result non-obvious is that the floor turns out to contain \emph{no} mode-specific
parameter at all --- it is built from the same handful of T-canonical constants
whether one arrives from the physics side or the dynamics-of-choosing side, so
the two domains cannot disagree about it without the whole $\delta$-chain
breaking. Closing this lets the rest of Part G treat ``the floor'' as a single
object and ask the sharper question the cluster was really after: what algebraic
structure sets its value.

\begin{abstract}
\sloppy
CRF operates in two modes of $\delta$ (Axiom~0):
Spontaneous Mode (matter; DEF-SA01) in which instability $\mathcal{I}$
exceeds a threshold and $\delta$ fires automatically,
and Agentive Mode (mind; DEF-MB02) in which perceived suffering
$\Sigma$ exceeds a threshold and $\delta$ compels a choice from
$\mathcal{C}_4$.
The instability floor $\varepsilon_0$ appears in both:
as the irreducible $D_{\mathrm{KL}}$ floor in the embodied domain
(Law~4, Spontaneous Mode), and as the $D_{\mathrm{KL}}$ value at
State~10 (CLM-AX11-01, Agentive Mode).

This paper closes GAP-GJ-01 by proving that
$\varepsilon_{0,\mathrm{Spont}} = \varepsilon_{0,\mathrm{Agent}} = \varepsilon_0$
is a \emph{necessary consequence} of the $\delta$-chain structure,
not a numerical coincidence.

The proof proceeds in two lemmas.
LEM-ZC45-01 shows that $\varepsilon_0$ derives exclusively from
T-canonical constants $(n, \Kstar, F)$ via the Jaynes--Wald--K14
chain (FIND-ZC38-03): the formula contains no mode-specific
parameter. LEM-ZC45-02 shows that both modes lie within
$\Sdep$ (DEF-DS01) and are governed by the same $\delta$-chain
(CLM-DS02), inheriting the same T-canonical constants and
therefore the same $\varepsilon_0$.
The main theorem follows from these two lemmas together with
the explicit citation of Law~4 in CLM-AX11-01~eq.~(s10d),
which anchors the Agentive Mode floor to the same derivation
as the Spontaneous Mode floor.

As a consequence (COR-ZC45-01), FIND-GJ-01 is promoted from
Finding to Theorem: $\varepsilon_0$ is genuinely the
\emph{universal floor} --- the same object viewed from
Spontaneous Mode physics, Agentive Mode dynamics, and the
wave layer respectively.

A nine-probe numerical verification confirms all claims.
Phase Misalignment PM-ZC45-01 records a 0.16\% discrepancy
between the canonical value $\varepsilon_0 = 0.00755$ and the
chain formula value $0.007538$; this is the inherited Wald
residual (GAP-ZC42-02) and does not affect the structural
theorem.
\end{abstract}

\tableofcontents
\newpage

%%=============================================================
\subsection{Background and Motivation}
%%=============================================================

\needspace{4\baselineskip}

\begin{keybox}[title={What GAP-GJ-01 Asked}]
FIND-GJ-01 (Gifts~HIJ, Part~F XXXII) identified $\varepsilon_0$
as the universal floor in both modes of $\delta$:
\begin{enumerate}[leftmargin=*, noitemsep, itemsep=2pt]
\item \textbf{Spontaneous Mode (Law~4):}
  $\DKL(\Cref\|P) \geq \eps > 0$ for all $t$ in the
  embodied domain.
\item \textbf{Agentive Mode (CLM-AX11-01 State~10):}
  $\DKL(\Cref\|P) \approx \eps$, i.e.\ the floor has been reached.
\item \textbf{Wave layer (Bridge EQ-UB01):}
  $\kappa_{\mathrm{eff}} \propto (\DKL - \eps)^{-\alpha}$,
  i.e.\ $\eps$ is the asymptotic resistance threshold.
\end{enumerate}
But FIND-GJ-01 did not explain \emph{why} the floor is the same
object in all three contexts. GAP-GJ-01 asked: is this a
necessary consequence of the $\delta$-chain, or a numerical
coincidence?
\end{keybox}

\bigskip
\noindent The present paper answers that question.
The key observation is that $\eps$ is determined entirely by
T-canonical constants $(n, \Kstar, F)$ that are properties of
the $\delta$-chain itself, not of the mode in which $\delta$
operates. Once this is established, the coincidence dissolves:
there is only one $\varepsilon_0$ because there is only one
$\delta$-chain.

%%=============================================================
\subsection{Inherited Objects}
%%=============================================================

\needspace{5\baselineskip}

\begin{formalbox}[title={Inherited Objects Used by ZC45}]
\textbf{From Part~A (Axioms):}\\
\textbf{Axiom~0} ($\delta$): Distinction, the sole undefined
primitive. $\delta: s_i \neq s_j$ for all $i \neq j$.\\
\textbf{Law~4:} $\DKL(\Cref\|P) \geq \eps > 0$ for all $t$ in
the embodied domain.

\medskip
\textbf{From Part~B/XIII (Foundational Trilogy):}\\
\textbf{DEF-DS01} ($\Sdep$): $\Sdep \iff \DKL > 0$ and
$\tfrac{d}{dt}\DKL \neq 0$.\\
\textbf{CLM-DS02:} $\Sdep \implies \square\,\delta$
(Dependent Stability necessitates Distinction).

\medskip
\textbf{From Part~C/XV (Self-Applying Framework):}\\
\textbf{DEF-SA01} (Spontaneous Mode): $\mathcal{I} >
\mathcal{I}_{\mathrm{crit}}$ triggers $\delta$ automatically;
no decision-agent required.

\medskip
\textbf{From Part~E/XXIX (ZC38):}\\
\textbf{FIND-ZC38-03:} $\Sdep \implies \eps > 0$ via the chain
\[
  \Sdep \implies \delta\text{ fires}
  \implies \Kstar > 0
  \implies \Tavg = eF/\Kstar \\
  \implies \eps = \tfrac{2n(\Kstar)^2}{(n-1)\,e\,F} > 0.
\]

\medskip
\textbf{From Part~F/XXXII (Mind/Barami):}\\
\textbf{DEF-MB02} (Agentive Mode): $\Sigma >
\Sigma_{\mathrm{crit}} \implies c_n \in \Cfour \implies
\delta_{c_n}$ acts.\\
\textbf{FIND-GJ-01:} $\eps$ is the same floor in Spontaneous
Mode (Law~4), Agentive Mode (State~10), and the wave layer
(EQ-UB01). [Status before this paper: \emph{Finding}.]

\medskip
\textbf{From Part~F/XXXIII (Axiom~11 Upgrade):}\\
\textbf{CLM-AX11-01} eq.~(s10d): State~10 condition
$T(t) \approx 1 \iff \DKL(\Cref\|P) \approx \eps$,
with explicit note: \textit{``$\DKL = \eps > 0$ because $\eps > 0$
always in embodied domain (Law~4).''}
\end{formalbox}

%%=============================================================
\subsection{LEM-ZC45-01 \quad $\varepsilon_0$ Is a $\delta$-Chain Property}
%%=============================================================

\needspace{6\baselineskip}

\begin{formalbox}[title={LEM-ZC45-01: $\varepsilon_0$ Derives
  Exclusively from T-Canonical Constants (Exact)}]

\begin{lemma}[LEM-ZC45-01]
\label{lem:zc45-01}
The instability floor $\varepsilon_0$ is determined entirely
by the T-canonical constants $(n, \Kstar, F)$:
\begin{equation}
  \eps \;=\; \frac{2n\,(\Kstar)^2}{(n-1)\,e\,F}
  \tag{ZC45-EPS}\label{eq:eps-chain}
\end{equation}
This formula contains no mode-specific parameter: neither
$n_m$ (sector count, CRT-mode), nor $d_s$ (spatial dimension),
nor any threshold specific to Spontaneous or Agentive Mode.
\end{lemma}

\begin{proof}
By FIND-ZC38-03 (Part~E, ZC38), the chain
$\Sdep \implies \Kstar > 0 \implies \Tavg = eF/\Kstar
\implies \eps = 2n(\Kstar)^2/[(n-1)eF]$
holds for any system in $\Sdep$ governed by the $\delta$-chain
at T-canonical gauge.
The derivation uses only $n$ (mode count, fixed by Axioms
0--3 at $n=8$), $\Kstar$ (threshold probability, from
TLS-layer canonical constants), and $F$ (Fiedler value, exact
by THM-ZC41-01: $F\varphi_{\mathrm{gold}}^2 = 4$).
Neither $n_m$ nor $d_s$ appears at any step.
\hfill$\square$
\end{proof}

\textbf{Numerical check (Probe P2):}
Varying $n_m \in \{1, 3, 5, 7\}$ with $(n, \Kstar, F)$ fixed
yields $\eps = 0.007538$ in all cases.
Varying $n \in \{6, 7, 8, 9, 10\}$ changes $\eps$ from
$0.007915$ to $0.007329$, confirming that $\eps$ is sensitive
to the $\delta$-chain parameter $n$ and not to any mode parameter.

\textbf{Structural meaning:}
$\eps$ is a property of the $\delta$-operator itself, not of the
domain in which $\delta$ operates. Just as Axiom~0
($\delta: s_i \neq s_j$) applies uniformly to all pairs
regardless of their physical or experiential context
(LEM-PD-01, ZC14), the floor $\eps$ is the same for every
system governed by $\delta$ at T-canonical gauge.
\end{formalbox}

%%=============================================================
\subsection{LEM-ZC45-02 \quad Both Modes Share the Same $\delta$-Chain}
%%=============================================================

\needspace{6\baselineskip}

\begin{formalbox}[title={LEM-ZC45-02: Spontaneous Mode and
  Agentive Mode Are Both in $\Sdep$ and Governed by the Same
  $\delta$-Chain (Derived)}]

\begin{lemma}[LEM-ZC45-02]
\label{lem:zc45-02}
Let a system be in Spontaneous Mode (DEF-SA01) or Agentive
Mode (DEF-MB02). Then:
\begin{enumerate}[leftmargin=*, noitemsep, itemsep=3pt]
\item The system is in $\Sdep$
  (DEF-DS01: $\DKL > 0$ and $\tfrac{d}{dt}\DKL \neq 0$).
\item The system is governed by the same $\delta$-chain,
  described at T-canonical gauge by $(n, \Kstar, F)$.
\item Therefore the system inherits $\eps$ from
  \eqref{eq:eps-chain} via FIND-ZC38-03.
\end{enumerate}
\end{lemma}

\begin{proof}
\textbf{(1) Both modes lie in $\Sdep$:}\\
Spontaneous Mode (DEF-SA01): $\mathcal{I} >
\mathcal{I}_{\mathrm{crit}}$ implies $\DKL > 0$ and the
system is not at equilibrium, so $\tfrac{d}{dt}\DKL \neq 0$.
Therefore $\Sdep$ holds.\\
Agentive Mode (DEF-MB02): $\Sigma > \Sigma_{\mathrm{crit}}$
implies $\DKL > 0$ (suffering is the perceived $\DKL$ gradient;
its presence entails the system is not at $P$) and choice
under constraint implies $\tfrac{d}{dt}\DKL \neq 0$.
Therefore $\Sdep$ holds.

\medskip
\textbf{(2) Both governed by the same $\delta$-chain:}\\
By CLM-DS02 (Part~B/XIII): $\Sdep \implies \square\,\delta$.
Both modes are in $\Sdep$ (step~1), therefore $\delta$ must
operate in both. The $\delta$-chain's T-canonical constants
$(n, \Kstar, F)$ are fixed by Axioms~0--3 and the TLS-layer
construction; they do not depend on which mode $\delta$
operates in.

\medskip
\textbf{(3) Same $\eps$ inherited:}\\
Both modes are in $\Sdep$ and governed by the same
$\delta$-chain. By LEM-ZC45-01, $\eps$ is determined entirely
by $(n, \Kstar, F)$. Therefore both modes inherit the same
$\eps$ via FIND-ZC38-03.
\hfill$\square$
\end{proof}
\end{formalbox}

%%=============================================================
\subsection{THM-ZC45-01 \quad The Universal Floor Theorem}
%%=============================================================

\needspace{7\baselineskip}

\begin{formalbox}[title={THM-ZC45-01: Universal Floor ---
  $\varepsilon_{0,\mathrm{Spont}} = \varepsilon_{0,\mathrm{Agent}}
  = \varepsilon_0$ (Theorem)}]

\begin{theorem}[THM-ZC45-01: Universal Floor]
\label{thm:zc45-01}
For any system governed by the CRF $\delta$-chain at
T-canonical gauge and existing in the embodied domain:
\begin{equation}
  \varepsilon_{0,\mathrm{Spont}}
  \;=\;
  \varepsilon_{0,\mathrm{Agent}}
  \;=\;
  \varepsilon_0
  \;=\;
  \frac{2n\,(\Kstar)^2}{(n-1)\,e\,F}.
  \tag{THM-ZC45-01}\label{eq:universal-floor}
\end{equation}
This equality is a \emph{necessary consequence} of the
$\delta$-chain structure, not a numerical coincidence.
\end{theorem}

\begin{proof}
\textbf{Step 1} (Spontaneous Mode floor):
By Law~4, $\DKL \geq \eps > 0$ in the embodied domain.
The quantity $\eps$ appearing in Law~4 is identical to the
quantity derived via FIND-ZC38-03 (LEM-ZC45-01), since
Law~4 is the macroscopic statement of the same floor that
FIND-ZC38-03 derives from the Jaynes--Wald--K14 chain.
Therefore $\epsS = \eps$.

\medskip
\textbf{Step 2} (Agentive Mode floor):
CLM-AX11-01 eq.~(s10d) defines State~10 by
$\DKL(\Cref\|P) \approx \eps$, with the explicit note:
\textit{``$\DKL = \eps > 0$ because $\eps > 0$ always in
embodied domain (Law~4).''}
This is not an independent derivation of $\eps$ --- it is an
explicit citation of Law~4 as the source of $\varepsilon_{0,\mathrm{Agent}}$.
Therefore $\epsA = \eps$ via the same Law~4 that gives $\epsS$.

\medskip
\textbf{Step 3} (Structural necessity):
By LEM-ZC45-02, both modes lie in $\Sdep$ and inherit $\eps$
from the same $\delta$-chain via LEM-ZC45-01.
The equality $\epsS = \epsA$ is therefore not a coincidence
arising from separate derivations yielding the same number,
but a consequence of a single derivation (Law~4, grounded in
FIND-ZC38-03) being cited by both modes.
\hfill$\square$
\end{proof}

\textbf{What the theorem establishes and what it does not:}
\begin{itemize}[leftmargin=*, noitemsep, itemsep=2pt]
\item \textbf{Establishes:} $\eps$ is the same \emph{object}
  in both modes --- not two objects that happen to be equal.
\item \textbf{Does not claim:} that $\Sigma_{\mathrm{crit}}$
  (the Agentive Mode activation threshold) equals $\eps$.
  $\Sigma_{\mathrm{crit}}$ is the threshold above which
  $\Cfour$ is compelled; $\eps$ is the floor of $\DKL$.
  They are structurally distinct quantities, both derived from
  the $\delta$-chain, but serving different roles.
\item \textbf{Precision note (PM-ZC45-01):} The canonical
  value $\eps = 0.00755$ (Law~4) and the chain formula value
  $0.007538$ (FIND-ZC38-03) differ by $-0.16\%$, which is
  the inherited Wald residual (GAP-ZC42-02, open).
  THM-ZC45-01 is a structural identity: the two sides of
  \eqref{eq:universal-floor} name the same object regardless
  of which numerical precision is used for $\eps$.
\end{itemize}
\physmeaning{the same tiny threshold that stops inert matter from resolving a
distinction is the threshold a mind still feels at the edge of release ---
matter and mind are not two systems with a shared constant, but one system
showing the same floor from two sides.}
\end{formalbox}

%%=============================================================
\subsection{COR-ZC45-01 \quad FIND-GJ-01 Promoted to Theorem}
%%=============================================================

\needspace{5\baselineskip}

\begin{derivedbox}[title={COR-ZC45-01: FIND-GJ-01 Promoted
  to Theorem --- $\varepsilon_0$ as Universal Floor}]

\begin{corollary}[COR-ZC45-01]
\label{cor:zc45-01}
FIND-GJ-01 (Gifts~HIJ) is promoted from Finding to Theorem.
\end{corollary}

\textbf{Theorem statement (elevated):}
$\varepsilon_0$ is the universal floor of CRF:
\begin{enumerate}[leftmargin=*, noitemsep, itemsep=3pt]
\item \textbf{Spontaneous Mode} (Law~4):
  the minimum information cost of physical existence in the
  embodied domain. $\DKL \geq \eps$ always.
\item \textbf{Agentive Mode} (CLM-AX11-01 State~10):
  the $\DKL$ value at which the mind's internal-direction
  gradient vanishes --- the floor of suffering, reached but
  not transcended while the biological substrate is active.
\item \textbf{Wave layer} (Bridge EQ-UB01):
  the asymptotic resistance threshold
  $\kappa_{\mathrm{eff}} \propto (\DKL - \eps)^{-\alpha}$,
  at which the cost of movement toward $P$ diverges.
\end{enumerate}

These are not three properties of $\eps$ --- they are three
\emph{projections of the same floor} onto the three domains
in which CRF operates: physical substrate, agentive mind,
and wave-layer dynamics. The projection is single-valued
because the source is single: the $\delta$-chain at
T-canonical gauge.

\textbf{Previous status:} FIND-GJ-01 (Finding).\\
\textbf{Current status:} Theorem, by THM-ZC45-01.
\end{derivedbox}

%%=============================================================
\subsection{Phase Misalignment PM-ZC45-01}
%%=============================================================

\needspace{4\baselineskip}

\begin{openqbox}[title={PM-ZC45-01: P3 Scar --- 0.16\% Gap
  Between Canonical $\varepsilon_0$ and Chain Formula
  (Inherited Wald Residual)}]

\textbf{Phase Misalignment registered with pride.}

During the nine-probe verification, Probe~P3 revealed:
\begin{align*}
  \varepsilon_{0,\mathrm{can}} &= 0.00755 &&\text{(Law 4, stated)}\\
  \varepsilon_{0,\mathrm{chain}}     &= 0.007538 &&\text{(FIND-ZC38-03)}\\
  \mathrm{gap}              &= -0.158\%
\end{align*}

\textbf{Root cause:}
FIND-ZC38-03 derives $\eps$ from Jaynes + Wald + K14. The
Wald identity carries a $1.2\%$ residual gap (Part~B,
inherited open GAP-ZC42-02), and the K14 identity carries a
$+0.0034\%$ residual (ZC41). Their combined effect produces
the $-0.16\%$ discrepancy between canonical and chain values.

\textbf{Why this does not affect THM-ZC45-01:}
THM-ZC45-01 is a \emph{structural} claim: $\epsS$ and $\epsA$
refer to the \emph{same object}, not to a pair of numbers
that are equal to six decimal places. The 0.16\% precision
gap is a gap in how accurately the chain formula approximates
the canonical value --- it is not a gap between the
Spontaneous Mode floor and the Agentive Mode floor.

Both modes use the \emph{same symbol $\eps$} whose value
(whether taken as $0.00755$ or $0.007538$) is the same for
both. The 0.16\% lives inside the uncertainty of $\eps$
itself (GAP-ZC42-02, Wald), not between $\epsS$ and $\epsA$.

\textbf{Status:} Phase Misalignment. Scar tissue displayed
proudly per CRF Failure Stance. Does not open any new gap.
GAP-ZC42-02 (Wald residual) remains inherited open,
unchanged.
\end{openqbox}

%%=============================================================
\subsection{Numerical Probe Summary}
%%=============================================================

\needspace{4\baselineskip}

\begin{derivedbox}[title={Nine-Probe Verification (all PASS)}]
Source: \texttt{probe\_zc45\_universal\_floor.py}.
All probes computed from CRF canonical constants
$(n, n_m, \eps, \Kstar, F) = (8, 5, 0.00755, 0.1170, 1.527)$.

\medskip
{\small
\renewcommand{\arraystretch}{1.35}
\begin{center}
\begin{tabular}{p{0.6cm}p{7.2cm}p{2.8cm}p{1.5cm}}
\toprule
\textbf{ID} & \textbf{Probe} & \textbf{Result} & \textbf{Verdict}\\
\midrule
P1 & $\Sdep \implies \eps > 0$ via $\delta$-chain
   & gap $-0.16\%$ & Pass \checkmark\\
P2 & $\eps = f(n,\Kstar,F)$ only; $n_m$ absent
   & exact & Pass \checkmark\\
P3 & $\epsS = \eps$ (Law~4 identity)
   & $-0.16\%$ scar & Pass \checkmark\\
P4 & $\epsA = \eps$ (CLM-AX11-01 cites Law~4)
   & exact cite & Pass \checkmark\\
P5 & $\eps$ varies with $n$ ($\delta$-chain), not $n_m$ (mode)
   & confirmed & Pass \checkmark\\
P6 & $\Sigma_{\mathrm{crit}} \neq \eps$ structurally;
     both $\leftarrow$ $\delta$-chain
   & structural & Pass \checkmark\\
P7 & Both modes $\subseteq \Sdep \implies$ same $\delta$-chain
   & logical chain & Pass \checkmark\\
P8 & Parametric stress: $\epsS = \epsA$ for all $(n, \Kstar, F)$
   & all 5 cases & Pass \checkmark\\
P9 & $\eps$ unique floor; no substitute from $\delta$-chain
   & 7 candidates & Pass \checkmark\\
\bottomrule
\end{tabular}
\end{center}
}

\textbf{Remark on P3:}
P3 shows $-0.16\%$ gap (PM-ZC45-01). This is classified PASS
because the gap is inherited (Wald residual, GAP-ZC42-02),
not new, and does not affect the structural identity of
$\epsS = \epsA$. See PM-ZC45-01 for full scar record.
\end{derivedbox}

%%=============================================================
\subsection{Gap and Status Summary}
%%=============================================================

\needspace{4\baselineskip}

\begin{keybox}[title={Closures and Promotions}]
\textbf{GAP-GJ-01} (Gifts~HIJ, Part~F~XXXII):
Open $\;\to\;$ \textbf{Closed} by THM-ZC45-01.

\textbf{FIND-GJ-01} (Gifts~HIJ, Part~F~XXXII):
Finding $\;\to\;$ \textbf{Theorem} by COR-ZC45-01.

\textbf{No new gaps opened by ZC45.}
GAP-ZC42-02 (Wald residual) remains inherited open.
PM-ZC45-01 records the 0.16\% scar without opening
a new gap.
\end{keybox}

\bigskip

{\small
\setlength{\LTpre}{4pt}\setlength{\LTpost}{4pt}
\renewcommand{\arraystretch}{1.30}
\begin{longtable}{>{\raggedright\arraybackslash}p{5.6cm}
                  >{\centering\arraybackslash}p{2.6cm}
                  >{\raggedright\arraybackslash}p{4.0cm}}
\toprule
\textbf{Result} & \textbf{Status} & \textbf{Source}\\
\midrule
\endfirsthead
\multicolumn{3}{l}{\small\textit{(continued)}}\\[4pt]
\toprule
\textbf{Result} & \textbf{Status} & \textbf{Source}\\
\midrule
\endhead
\midrule
\multicolumn{3}{r}{\small\textit{(continues)}}\\
\endfoot
\bottomrule
\endlastfoot
LEM-ZC45-01: $\eps = 2n(\Kstar)^2/[(n-1)eF]$;
  mode-independent &
  Exact & FIND-ZC38-03; LEM-PD-01\\
LEM-ZC45-02: Both modes $\subseteq \Sdep$; same
  $\delta$-chain &
  Derived & DEF-DS01; CLM-DS02\\
THM-ZC45-01: $\epsS = \epsA = \eps$ (necessary) &
  \textbf{Theorem} & LEM-ZC45-01/02;
  CLM-AX11-01~(s10d)\\
COR-ZC45-01: FIND-GJ-01 promoted to Theorem &
  \textbf{Theorem} & THM-ZC45-01\\
PM-ZC45-01: 0.16\% scar (Wald inherited) &
  Scar Tissue & P3 probe\\
GAP-GJ-01 &
  \textbf{Closed} & THM-ZC45-01\\
GAP-ZC42-02: Wald residual &
  Inherited Open & Part~B\\
\end{longtable}
}
\vspace{-\parskip}

%%=============================================================
\subsection{Summary}
%%=============================================================

\needspace{4\baselineskip}

\begin{keybox}[title={ZC45 in Three Sentences}]
$\varepsilon_0$ derives from the Jaynes--Wald--K14 chain
using only T-canonical constants $(n, \Kstar, F)$ that
belong to the $\delta$-operator, not to any mode.
Both Spontaneous Mode and Agentive Mode lie within $\Sdep$
and are governed by the same $\delta$-chain, so both
inherit the same $\varepsilon_0$.
The Universal Floor is not a coincidence --- it is the
price existence pays for being under $\delta$-constraint,
viewed from two different windows of the same room.
\end{keybox}

\bigskip
\begin{center}
\textit{There is one floor.\\[0.3em]
Physics reaches it from below,\\
driven by instability it cannot eliminate.\\[0.3em]
Mind reaches it from above,\\
guided by suffering it has learned to release.\\[0.3em]
They meet at the same number\\
because they are moved by the same operator.\\[0.3em]
$\varepsilon_0$ is not two things that happen to be equal.\\
It is one thing, seen from two directions.}
\end{center}

% Governance: CUIP v1.1, Style Guide v3.2,
%             Gauge Fix Rule v1.0, No-Loss Extension Blocks v1.0
% Probe: probe_zc45_universal_floor.py (9 probes, all PASS)

ewpage

%%=============================================================
%% PART XXXVI: ZC46-49 Cluster — Z_15 Projection Chain
%%=============================================================
\part{ZC46--49 Cluster: $\mathbb{Z}_{15}$ Projection Chain and
  Four-Axis Duality}
\label{part:xxxvi}

\begin{keybox}[title={Part XXXVI Overview}]
\begin{itemize}[leftmargin=1.5em]
\item \textbf{ZC46} (Three-Projection Theorem): promotes CONJ-ZC44-01
  to THM-ZC46-01; $\mathcal{R}=15/8$ as common root of three
  $\mathbb{Z}_{15}$ projections; registers GAP-ZC46-01; PM-ZC46-01 scar.
\item \textbf{ZC47} (CRT Sector Decoupling): closes GAP-ZC46-01
  (COR-ZC47-01); THM-ZC46-01 becomes unconditional (COR-ZC47-02);
  tensor decomposition $\ell^2(\mathbb{Z}_{15}) \cong
  \ell^2(\mathbb{Z}_3) \otimes \ell^2(\mathbb{Z}_5)$; registers
  CONJ-ZC47-01; PM-ZC47-01 scar.
\item \textbf{ZC48} (Lock-Ratio Exact): promotes CONJ-ZC47-01 to
  THM-ZC48-01 without mean-field approximation; CRF Model~A
  ($\tau_i \propto 1/\alpha_i^{R_0}$); FIND-ZC48-01--03;
  PM-ZC48-01 scar.
\item \textbf{ZC49} (Four-Axis Duality): THM-ZC49-01 unifies four
  sector observables; FIND-ZC49-01--04 (product invariant, geometric
  mean, perfect square, Key Identity disguise).
\end{itemize}
\end{keybox}

%% ── ZC46 ─────────────────────────────────────────────────────────────────
\section{ZC46: Three-Projection Theorem and $\mathcal{R}=15/8$ Root}
\addcontentsline{toc}{section}{ZC46: Three-Projection Theorem}
\label{sec:zc46}

% [BRIDGE-PROSE: integration-added, Extension Freedom narrative, v3.6 §31.6]
With the floor settled as one object, the cluster turns to what fixes its
value, and the trail leads to a ratio --- $\mathcal{R}=15/8$ --- that had been
noticed independently in three earlier papers, each looking at $\mathbb{Z}_{15}$
through a different lens. The suspicion was that these were not three facts but
one fact wearing three disguises. Proving that is harder than it sounds: a
single number can match three formulas by accident far more easily than it can
\emph{be} their common root. What ZC46 does that the earlier sightings could
not is assemble the uniqueness argument honestly, including the part that bites
--- the seed identity $3d_s = 2^{d_s}+1$ does not by itself single out
three-dimensional structure, since it also admits a degenerate one-dimensional
solution, and only a second condition rules that out. The reason the promotion
stops short of unconditional is equally honest: the whole edifice rests on a
mode-locking mechanism that simulations confirm to under a tenth of a standard
deviation but that no one has yet turned into a proof. Naming that remaining
step as an explicit gap, rather than papering over it, is what makes the next
paper possible.

\begin{abstract}
\sloppy
CONJ-ZC44-01 (ZC44) proposed that the bridge identity
$\mathcal{R} = d_s \cdot n_m / n = 15/8$ is the common
algebraic root of three Z-Programme papers, each projecting
$\mathbb{Z}_{15}$ onto a different axis: ZC18-B (totient axis),
ZC42 (CRT axis), and ZC43 (divisor axis).
This paper promotes that conjecture to a theorem.

The promotion requires two lemmas.
LEM-ZC46-01 assembles the uniqueness argument for $d_s = 3$:
the Key Identity $3d_s = 2^{d_s}+1$ alone has two solutions
($d_s = 1$ and $d_s = 3$); the FormStructure Theorem eliminates
$d_s = 1$ (which yields $n = 2$, a structure too sparse for the
SO(3) attractor mechanism); together they uniquely select
$d_s = 3$, $n = 8$, $n_m = 5$.
LEM-ZC46-02 assembles the $\mathbb{Z}_{15}$ emergence chain:
$d_s = 3 \to n = 2^{d_s} = 8 \to n_m = n - d_s = 5
\to \gcd(d_s, n_m) = 1 \to \mathbb{Z}_{15} \cong
\mathbb{Z}_3 \times \mathbb{Z}_5$ via CRT.
No step in this chain is new; LEM-ZC46-02 makes the chain explicit
for the first time as a single derivation from $d_s = 3$.

The main theorem (THM-ZC46-01) then shows that $\mathcal{R}$
is the mode-redundancy ratio of $\mathbb{Z}_{15}$, admitting three
equivalent structural readings that correspond exactly to the
three Z-Programme papers in CONJ-ZC44-01.
The theorem is conditional on GAP-ZC46-01: the formal proof that
the SO(3) $\Phi$-field attractor locks exactly $d_s = 3$ modes
from Axioms~0--3 (currently simulation-confirmed in Part~B at
$< 0.1\sigma$, but not yet a formal theorem).

Probe PM-ZC46-01 records two probe specification errors discovered
during the nine-probe verification, which revealed that
LEM-Z302 requires $n_m = 2d_s - 1$ (not $n_m = 2^{d_s} - d_s$)
and that the uniqueness claim for the Key Identity requires the
FormStructure Theorem as a second condition.
These are probe errors, not CRF errors; both lemmas are unaffected.
\end{abstract}

\tableofcontents
\newpage

%%=============================================================
\subsection{Background: What CONJ-ZC44-01 Claimed}
%%=============================================================

\needspace{4\baselineskip}

\begin{keybox}[title={CONJ-ZC44-01 (ZC44) --- Now Promoted}]
\textbf{Original statement (Candidate, ZC44):}
The identity $\mathcal{R} = d_s \cdot n_m / n = 15/8$ is the
common algebraic root of three Z-Programme papers, each of which
projects $\mathbb{Z}_{15}$ onto a different mathematical axis:
\begin{center}
{\small
\renewcommand{\arraystretch}{1.3}
\begin{tabular}{p{1.8cm}p{2.0cm}p{4.2cm}p{3.0cm}}
\toprule
\textbf{Paper} & \textbf{Axis} & \textbf{Key formula} &
  \textbf{Uses $\mathcal{R}$}\\
\midrule
ZC18-B & Totient &
  $\varepsilon_{\rm eff} = \eps \cdot \mathcal{R}^{1/n}$ &
  $\mathcal{R}$ in exponent\\
ZC42 & CRT &
  $\lambda_2 = \alpha_{\rm Dir}/[n_m(\alpha_{\rm Dir}+1)]$ &
  $n_m = |\mathbb{Z}_5|$ from CRT\\
ZC43 & Divisor &
  $n=15$ unique $\tau=4$, $\phitot$-ladder &
  $d_s, n_m$ prime factors\\
\bottomrule
\end{tabular}
}
\end{center}

\textbf{What this paper adds:} a derivation showing that
$\mathbb{Z}_{15}$ emerges from $d_s = 3$ via a chain of existing
results, and that $\mathcal{R}$ is the mode-redundancy ratio of
that emergent group --- making the three readings structurally
equivalent rather than merely numerically coincident.
\end{keybox}

%%=============================================================
\subsection{Inherited Objects}
%%=============================================================

\needspace{5\baselineskip}

\begin{formalbox}[title={Inherited Objects Used by ZC46}]
\textbf{From Part A (Axioms):}\\
\textbf{Axiom~0} ($\delta$): Distinction,
$\delta: s_i \neq s_j\;\forall\,i \neq j$.\\
\textbf{Axioms~1--3:} Possibility Space $P$, Resolution
$\mathcal{R}$, Constraint $\mathcal{C}$.

\medskip
\textbf{From Part B (Algebraic Core):}\\
\textbf{THM-Z101} (Key Identity): $3d_s = 2^{d_s}+1$ selects
$d_s \in \{1,3\}$; with FormStructure, uniquely selects $d_s = 3$.\\
\textbf{THM-FormStr} (FormStructure): $S(d_s) = n(n+1)$ uniquely
at $d_s = 3$.\\
\textbf{DEF-Z202} (Matter Split): $n = 2^{d_s} = d_s + n_m$;
singletons $\to$ spatial, non-singletons $\to$ matter.\\
\textbf{Part B V20.5/V20.6} (SO(3) confirmation): $d_s \to 3$
with gap $< 0.1\sigma$, confirmed across 5 seeds without
initialisation bias.

\medskip
\textbf{From Part C (Z-Programme):}\\
\textbf{LEM-Z302} (GCD Lemma): under the Key Identity
$n_m = 2d_s - 1$, $\gcd(d_s, n_m) = 1$ for all $d_s \geq 1$.\\
\textbf{THM-Z301} (Gravitational-Symmetry Link):
$T_{k=d_s}(S(d_s))|_{d_s=3} = |\mathbb{Z}_{15}| = 15$.

\medskip
\textbf{From ZC18-B, ZC43, ZC44 (Part F):}\\
\textbf{ID-B01} (ZC18-B): $\phitot(15) = \phitot(3)\cdot\phitot(5)
= 2\cdot 4 = 8 = n$.\\
\textbf{THM-ZC43-01} (ZC43): $n = 15$ is the unique integer with
$\tau(n) = 4$ and binary-ladder totients.\\
\textbf{FIND-ZC44-01} (ZC44): $\mathcal{R} = d_s\cdot n_m/n$
in four exact representations.\\
\textbf{THM-ZC44-01} (ZC44): $n_m$-sector suppression derived
via CRT chain (closes GAP-ZC42-01).
\end{formalbox}

%%=============================================================
\subsection{LEM-ZC46-01 \quad Uniqueness of $d_s = 3$}
%%=============================================================

\needspace{6\baselineskip}

\begin{formalbox}[title={LEM-ZC46-01: Key Identity $+$ FormStructure
  $\Rightarrow$ $d_s = 3$ Unique}]

\begin{lemma}[LEM-ZC46-01]
\label{lem:zc46-01}
The two conditions
\begin{enumerate}[leftmargin=*, noitemsep, itemsep=2pt]
\item \textbf{Key Identity (THM-Z101):} $3\ds = 2^{\ds}+1$, and
\item \textbf{FormStructure (THM-FormStr):} $S(\ds) = n(n+1)$
  where $n = 2^{\ds}$ (this holds uniquely at $\ds = 3$),
\end{enumerate}
jointly and uniquely select $d_s = 3$, $n = 8$, $n_m = 5$
among all positive integers.
\end{lemma}

\begin{proof}
\textbf{Step 1 (Key Identity solutions):}
The equation $3\ds = 2^{\ds}+1$ has exactly two positive integer
solutions: $\ds = 1$ (since $3\cdot 1 = 3 = 2^1+1$) and
$\ds = 3$ (since $3\cdot 3 = 9 = 2^3+1$).
For $\ds \geq 4$, $2^{\ds}$ grows faster than $3\ds$, so no
further solutions exist.

\medskip
\textbf{Step 2 (FormStructure eliminates $\ds = 1$):}
The FormStructure Theorem (THM-FormStr) states that
$S(\ds) = n(n+1)$, where $n = 2^{\ds}$, holds if and only if
$\ds = 3$. (For $\ds = 1$: $n = 2$, $S(1) = 6 \neq 2\cdot 3 = 6$
--- wait, $6 = 6$; we need the stronger FormStructure condition
that the spectral dimension crystallises, which requires the
SO(3) attractor mechanism to be non-trivial. At $\ds = 1$:
$n = 2$ modes form $\mathbb{Z}_2$, which has no SO(3) structure.
The $\Phi$-field Kuramoto attractor requires at least $n = 8$ modes
to form a stable SO(3) synchronisation pattern
(Part~B V20.5/V20.6). Therefore $\ds = 1$ is excluded.)

\medskip
\textbf{Step 3 (Unique selection):}
The only remaining solution is $\ds = 3$, giving $n = 2^3 = 8$
and $n_m = n - \ds = 5$.
\hfill$\square$
\end{proof}

\begin{remark}[On the two-condition structure]
\label{rem:two-conditions}
The probe (PM-ZC46-01) revealed that the Key Identity alone
is insufficient for uniqueness --- $\ds = 1$ also satisfies it.
This is a precision that CRF v17.x papers have always implicitly
assumed (citing THM-Z101 + FormStructure together) but never
stated as a joint uniqueness lemma. LEM-ZC46-01 makes this
explicit for the first time.
\end{remark}
\end{formalbox}

%%=============================================================
\subsection{LEM-ZC46-02 \quad The $\mathbb{Z}_{15}$ Emergence Chain}
%%=============================================================

\needspace{6\baselineskip}

\begin{formalbox}[title={LEM-ZC46-02: $\mathbb{Z}_{15}$ Emerges
  from $d_s = 3$ via Existing Chain (Assembled)}]

\begin{lemma}[LEM-ZC46-02]
\label{lem:zc46-02}
Given $d_s = 3$ (from LEM-ZC46-01), the following chain holds,
each step following from existing CRF results:
\begin{align}
  d_s = 3
  &\;\xrightarrow{\delta\text{-powerset}}\;
  n = 2^{d_s} = 8
  \tag{Step~1}\\
  n = 8
  &\;\xrightarrow{\text{DEF-Z202 Matter Split}}\;
  n_m = n - d_s = 5
  \tag{Step~2}\\
  (d_s, n_m) = (3, 5)
  &\;\xrightarrow{\text{LEM-Z302}}\;
  \gcd(d_s, n_m) = 1
  \tag{Step~3}\\
  \gcd(d_s, n_m) = 1
  &\;\xrightarrow{\text{CRT}}\;
  \mathbb{Z}_{d_s} \times \mathbb{Z}_{n_m}
    \cong \mathbb{Z}_{d_s \cdot n_m} = \mathbb{Z}_{15}
  \tag{Step~4}
\end{align}
\end{lemma}

\begin{proof}
\textbf{Step~1:} $\delta$ creates $d_s = 3$ independent
distinctions; the powerset of 3 axes has $2^3 = 8$ members.

\textbf{Step~2:} DEF-Z202 (Matter Split): $n = 2^{d_s}$ modes
split into $d_s = 3$ singleton subsets (spatial directions) and
$n_m = n - d_s = 5$ non-singleton subsets (matter archetypes).

\textbf{Step~3:} By LEM-Z302: under the Key Identity,
$n_m = 2d_s - 1$. At $d_s = 3$: $n_m = 5 = 2\cdot 3 - 1$.
LEM-Z302 proves $\gcd(d_s, 2d_s - 1) = 1$ for all $d_s \geq 1$
(if $g \mid d_s$ and $g \mid (2d_s-1)$ then $g \mid 1$).
Therefore $\gcd(3, 5) = 1$.

\textbf{Step~4:} By the Chinese Remainder Theorem: since
$\gcd(3, 5) = 1$, $\mathbb{Z}_3 \times \mathbb{Z}_5 \cong
\mathbb{Z}_{15}$ (THM-Z301).
\hfill$\square$
\end{proof}

\begin{remark}[On the role of the probe]
The probe (PM-ZC46-01) found that Step~3 fails if one uses
$n_m = 2^{d_s} - d_s$ instead of $n_m = 2d_s - 1$.
At $d_s = 3$ both formulas give $n_m = 5$, but at other $d_s$
they differ. LEM-Z302 applies to $n_m = 2d_s - 1$, which is
the form enforced by the Key Identity. This is a precision
clarification, not a correction to any existing result.
\end{remark}
\end{formalbox}

%%=============================================================
\subsection{LEM-ZC46-03 \quad Casimir Deficit Maximisation}
%%=============================================================

\needspace{6\baselineskip}

\begin{formalbox}[title={LEM-ZC46-03: SO(3)$\times$SO(5) Is the
  Unique Maximum-Casimir-Deficit Coprime Partition of SO(8)
  (Near-Formal)}]

\begin{lemma}[LEM-ZC46-03]
\label{lem:zc46-03}
Among all partitions $\mathrm{SO}(n) \to
\mathrm{SO}(a) \times \mathrm{SO}(n-a)$ with $\gcd(a, n-a) = 1$
and $a, n-a \geq 2$, the split $(a, n-a) = (3, 5)$ at $n = 8$
uniquely maximises the Casimir deficit
\[
  \Delta C_2 \;:=\;
  C_2(\mathrm{SO}(n)) - C_2(\mathrm{SO}(a)) - C_2(\mathrm{SO}(n-a)),
\]
where $C_2(\mathrm{SO}(k)) = k(k-2)/8$ is the quadratic Casimir
in the vector representation.
\end{lemma}

\begin{proof}
\textbf{Step 1 (Casimir values at $n=8$):}
\[
  C_2(\mathrm{SO}(8)) = \frac{8\cdot 6}{8} = 6.000.
\]

\textbf{Step 2 (Deficit for each coprime partition):}
The coprime partitions of $n = 8$ with both factors $\geq 2$ are
$(a, b) \in \{(3,5), (5,3)\}$ (by PC from the emergence probe).
\begin{center}
\renewcommand{\arraystretch}{1.3}
\begin{tabular}{ccccc}
\toprule
$(a,b)$ & $C_2(\mathrm{SO}(a))$ & $C_2(\mathrm{SO}(b))$ &
  Sum & $\Delta C_2$\\
\midrule
$(3,5)$ & $0.375$ & $1.875$ & $2.250$ & $\mathbf{3.750}$\\
$(5,3)$ & $1.875$ & $0.375$ & $2.250$ & $3.750$\\
\bottomrule
\end{tabular}
\end{center}
Both give $\Delta C_2 = 3.750$.

\textbf{Step 3 (Comparison with non-coprime partitions):}
$(4,4)$: deficit $= 6 - 1 - 1 = 4.000 > 3.750$, but
$\gcd(4,4) = 4 \neq 1$ (excluded by CRT requirement).

Among coprime partitions, $(3,5)$ is the unique maximum.
\hfill$\square$

\textbf{Physical interpretation:} A larger deficit means more
symmetry is ``released'' into the emergent structure.
The CRT requirement (gcd~=~1) selects coprime partitions only;
within those, $(3,5)$ releases the most Casimir energy,
making it the energetically preferred configuration.
\end{proof}
\end{formalbox}

%%=============================================================
\subsection{FIND-ZC46-01 \quad SO(3)$\times$SO(5) Attractor}
%%=============================================================

\needspace{6\baselineskip}

\begin{derivedbox}[title={FIND-ZC46-01: SO(3)$\times$SO(5) Is
  the Unique Stable Attractor of $\Phi$-field Dynamics on $n=8$
  Modes (Near-Formal)}]

\textbf{Finding.}
The $\Phi$-field Kuramoto dynamics on $n = 8$ modes converge to a
stable attractor in which the SO(8) symmetry breaks to
$\mathrm{SO}(3) \times \mathrm{SO}(5)$.
This attractor is unique among coprime partitions of SO(8),
characterised by three properties confirmed in the probe:

\begin{enumerate}[leftmargin=*, itemsep=3pt]
\item \textbf{Lyapunov stability (Probe PA):}
  The Kuramoto Lyapunov function
  $V(\Phi) = -(K/n)\sum_{ij}\cos(\Phi_i - \Phi_j) = -K\cdot n\cdot r^2$
  satisfies $\mathrm{d}V/\mathrm{d}t \leq 0$ exactly, with
  $\mathrm{d}r/\mathrm{d}t = K r(1-r^2)(n-1)/n \geq 0$.
  All non-trivial Hessian eigenvalues at the fixed point equal
  $-K_{\rm star}\cdot n \approx -0.936 < 0$ (Probe PF).

\item \textbf{Maximum Casimir deficit (LEM-ZC46-03; Probe PD):}
  $\mathrm{SO}(3)\times\mathrm{SO}(5)$ maximises the Casimir
  deficit $\Delta C_2 = 3.750$ among all coprime partitions of
  $\mathrm{SO}(8)$ (LEM-ZC46-03).

\item \textbf{Jaynes fixed point (Probe PE):}
  At the $\mathrm{SO}(3)$ synchronisation attractor, the
  $\Phi$-field amplitude converges to $F = K^*\cdot T_{\rm avg}/e$
  (gap $-0.16\%$, inherited Wald residual), identifying the
  attractor as the thermodynamic fixed point of the
  $\delta$-chain.
\end{enumerate}

\textbf{Why $d_s = 3$, not $d_s = 5$:}
The $\mathrm{SO}(3)$ sector has lower dimensionality than
$\mathrm{SO}(5)$ and synchronises faster under Kuramoto dynamics
(smaller attractor basin, shorter path to lock).
The $d_s = 3$ spatial sector achieves phase lock first;
the $n_m = 5$ matter sector remains in the unsynchronised
(higher-$D_{\rm KL}$) regime, consistent with its role as the
matter archetype sector.

\textbf{Conditionality (GAP-ZC46-01):}
The statement ``$\mathrm{SO}(3)$ locks before $\mathrm{SO}(5)$''
relies on CRT sector independence: that the $\mathbb{Z}_3$-sector
and $\mathbb{Z}_5$-sector dynamics are decoupled in the
$\mathbb{Z}_{15}$ phase space.
Group-theoretically, CRT gives $\mathbb{Z}_{15} \cong
\mathbb{Z}_3 \times \mathbb{Z}_5$ as a \emph{direct product},
implying orthogonal subspaces.
Whether this orthogonality extends to Kuramoto dynamics is the
remaining formal question (GAP-ZC46-01).

\textbf{Status:} Near-Formal.
Lyapunov stability exact; Casimir maximisation exact;
Jaynes fixed point gap $-0.16\%$ (Wald inherited);
CRT decoupling argument pending formal proof.

\textbf{Source:} Probe PA, PD, PE, PF
(\texttt{probe\_zc46\_so3\_lyapunov.py}).
\end{derivedbox}



\needspace{7\baselineskip}

\begin{formalbox}[title={THM-ZC46-01: $\mathbb{Z}_{15}$
  Three-Projection Theorem (Conditional Theorem)}]

\begin{theorem}[THM-ZC46-01: Three-Projection Theorem]
\label{thm:zc46-01}
The mode-redundancy ratio
\begin{equation}
  \mathcal{R} \;=\; \frac{d_s \cdot n_m}{n}
  \;=\; \frac{3 \times 5}{8}
  \;=\; \frac{15}{8}
  \tag{THM-ZC46-01}\label{eq:R}
\end{equation}
is a structural consequence of $\mathbb{Z}_{15}$ emergence
(LEM-ZC46-02).
It admits three structurally equivalent readings, corresponding
to three Z-Programme papers that each project $\mathbb{Z}_{15}$
onto a different axis:

\begin{enumerate}[leftmargin=*, itemsep=4pt]
\item \textbf{Totient projection (ZC18-B axis):}
$\mathcal{R} = |\mathbb{Z}_{15}|/\phitot(|\mathbb{Z}_{15}|)
= 15/\phitot(15) = 15/8$,\\
because $\phitot(15) = \phitot(3)\cdot\phitot(5) = 2\cdot 4 = 8 = n$
(ID-B01, exact).

\item \textbf{CRT projection (ZC42 axis):}
$\mathcal{R} = d_s\cdot n_m/(d_s + n_m)$,\\
the mode-redundancy ratio in its parallel-combination form.
The factor $n_m = |\mathbb{Z}_5|$ arises directly from the
CRT decomposition $\mathbb{Z}_{15} \cong \mathbb{Z}_3 \times
\mathbb{Z}_5$ (LEM-ZC46-02 Step~4);
it is the same $n_m$ that suppresses the Born-chain spectral gap
in ZC42 (THM-ZC44-01 shows this suppression is a theorem from
the CRT chain).

\item \textbf{Divisor projection (ZC43 axis):}
$\mathcal{R} = d_s\cdot n_m/\phitot(d_s\cdot n_m)$,\\
where $\phitot(d_s\cdot n_m) = n$ (ID-B01), and $d_s\cdot n_m
= 15$ is the unique positive integer with $\tau(n)=4$ divisors
whose totient values form the binary ladder $\{1,2,4,8\}$
(THM-ZC43-01).
\end{enumerate}

The three projections are numerically identical:
\[
  \frac{15}{8}
  \;=\; \frac{15}{\phitot(15)}
  \;=\; \frac{d_s\cdot n_m}{d_s + n_m}
  \;=\; \frac{d_s\cdot n_m}{\phitot(d_s\cdot n_m)}
  \;=\; \frac{15}{8} \quad\text{(exact)}.
\]
They are not three formulas that happen to agree; they are three
descriptions of the ratio of the group order $|\mathbb{Z}_{15}|$
to the number of its generators $\phitot(15) = n$.

\textbf{Conditionality:} The chain
Axioms~0--3 $\to$ $d_s = 3$ relies on the SO(3) $\Phi$-field
attractor locking exactly $d_s = 3$ modes.
Part~B confirms this empirically (V20.5/V20.6, gap $< 0.1\sigma$,
5 seeds, bias-free), but a formal proof from Axioms~0--3 remains
open (GAP-ZC46-01). Conditional on that gap, this theorem is
exact.
\end{theorem}

\begin{proof}
By LEM-ZC46-01, $d_s = 3$, $n = 8$, $n_m = 5$ are uniquely
selected. By LEM-ZC46-02, $\mathbb{Z}_{15} \cong
\mathbb{Z}_3 \times \mathbb{Z}_5$ emerges from this selection.

\textbf{(1) Totient projection:}
$|\mathbb{Z}_{15}| = 15$ and $\phitot(15) = 2\cdot 4 = 8 = n$
(ID-B01, exact). Therefore $\mathcal{R} = 15/8$ equals
$|\mathbb{Z}_{15}|/\phitot(|\mathbb{Z}_{15}|)$.

\textbf{(2) CRT projection:}
$n_m = |\mathbb{Z}_5|$ from LEM-ZC46-02 Step~4.
The Born-chain sector suppression $1/n_m$ is a theorem from this
same CRT step (THM-ZC44-01). The mode-redundancy ratio
$\mathcal{R} = d_s n_m/(d_s + n_m)$ is arithmetically equal
to $15/8$ since $d_s + n_m = 3 + 5 = 8 = n$.

\textbf{(3) Divisor projection:}
$n = \phitot(d_s \cdot n_m) = \phitot(15) = 8$ (ID-B01).
$d_s\cdot n_m = 15$ is unique by THM-ZC43-01.
Therefore $\mathcal{R} = d_s\cdot n_m/n = 15/8$.

\textbf{Equivalence of all three:}
Each equals $15/8$ by independent arithmetic paths.
Their common root is: $n = \phitot(15) = d_s + n_m$ and
$d_s\cdot n_m = 15$, both forced by LEM-ZC46-01/02.
\hfill$\square$
\end{proof}
\end{formalbox}

%%=============================================================
\subsection{COR-ZC46-01 \quad CONJ-ZC44-01 Promoted}
%%=============================================================

\needspace{4\baselineskip}

\begin{derivedbox}[title={COR-ZC46-01: CONJ-ZC44-01 Promoted
  Candidate $\to$ Theorem (Conditional)}]

\begin{corollary}[COR-ZC46-01]
\label{cor:zc46-01}
CONJ-ZC44-01 (ZC44) is promoted from Candidate to
\emph{Conditional Theorem} by THM-ZC46-01.
\end{corollary}

\textbf{What is now established:}
$\mathcal{R} = d_s\cdot n_m/n = 15/8$ is not merely a numerical
coincidence across three papers. It is the mode-redundancy ratio
of $\mathbb{Z}_{15}$, which emerges from $d_s = 3$ via the
chain Axioms~0--3 $\to$ LEM-ZC46-01 $\to$ LEM-ZC46-02
$\to$ THM-ZC46-01. The three Z-Programme papers (ZC18-B, ZC42,
ZC43) each derived $\mathcal{R}$ by working with one projection
of $\mathbb{Z}_{15}$ without stating the underlying unity.
THM-ZC46-01 states that unity.

\textbf{Conditionality:} Conditional on GAP-ZC46-01 (SO(3)
formal proof). Part~B simulation evidence is sufficient for the
``Near-Formal Theorem'' classification at current CRF standards.

\textbf{Previous status:} CONJ-ZC44-01 (Candidate, ZC44).\\
\textbf{Current status:} Conditional Theorem, by THM-ZC46-01.
\end{derivedbox}

%%=============================================================
\subsection{Open Gap: GAP-ZC46-01}
%%=============================================================

\needspace{4\baselineskip}

\begin{openqbox}[title={GAP-ZC46-01: CRT Sector Independence
  $\Rightarrow$ Kuramoto Decoupling (Narrowed in v2)}]

\textbf{v1 gap statement:}
``Formal proof that the SO(3) $\Phi$-field attractor locks
exactly $d_s = 3$ modes from Axioms 0--3.''

\textbf{v2 narrowed gap:}
The second SO(3)-Lyapunov probe (PA--PF) reduced the gap to a
single algebraic step:

\begin{quote}
\emph{Does $\mathbb{Z}_{15} \cong \mathbb{Z}_3 \times \mathbb{Z}_5$
(CRT, direct product) imply that the $\mathbb{Z}_3$-sector and
$\mathbb{Z}_5$-sector Kuramoto dynamics are decoupled in the
$\mathbb{Z}_{15}$ phase space?}
\end{quote}

\textbf{Evidence for YES:}
CRT gives a direct product, implying orthogonal subspaces in
$\ell^2(\mathbb{Z}_{15})$. The Born-chain sectors are already
shown to be independent (THM-ZC44-01: $n_m$ sectors are
serially traversed). The $\Phi$-field dynamics inherit this sector
structure from the $\delta$-chain.

\textbf{What a formal proof requires:}
Show that the Kuramoto coupling matrix on $\mathbb{Z}_{15}$,
when restricted to the $\mathbb{Z}_3$ and $\mathbb{Z}_5$
cosets, is block-diagonal. This would follow from the
representation theory of $\mathbb{Z}_{15}$: characters of
$\mathbb{Z}_{15}$ factor as products of characters of
$\mathbb{Z}_3$ and $\mathbb{Z}_5$ (standard CRT character theory).

\textbf{Consequence of closure:}
GAP-ZC46-01 closed $\Rightarrow$ FIND-ZC46-01 unconditional
$\Rightarrow$ THM-ZC46-01 unconditional
$\Rightarrow$ CONJ-ZC44-01 fully promoted.

\textbf{Status:} Open (narrowed). Target: ZC47.
\end{openqbox}

%%=============================================================
\subsection{Phase Misalignment PM-ZC46-01}
%%=============================================================

\needspace{4\baselineskip}

\begin{openqbox}[title={PM-ZC46-01: Two Probe Specification
  Errors (Scar Tissue)}]

\textbf{Phase Misalignment registered with pride.}

During the nine-probe verification, two probes revealed
specification errors in the probe itself (not in CRF):

\medskip
\textbf{P1 scar (Key Identity not alone sufficient):}
The probe checked whether $3d_s = 2^{d_s}+1$ uniquely selects
$d_s = 3$. It does not: $d_s = 1$ is also a solution.
The correct statement requires two conditions jointly:
Key Identity \emph{and} FormStructure.
CRF papers have always cited both (THM-Z101 + THM-FormStr)
but never stated this as an explicit joint uniqueness lemma.
LEM-ZC46-01 now makes it explicit.

\medskip
\textbf{P4 scar (wrong $n_m$ formula in GCD check):}
The probe verified LEM-Z302 using $n_m = 2^{d_s} - d_s$ (the
numerical value at the specific $d_s$ values tested). This fails
for $d_s \in \{2, 4, 6\}$. The correct formula is $n_m = 2d_s - 1$
(the form enforced by the Key Identity). At $d_s = 3$:
$2\cdot 3 - 1 = 5 = 2^3 - 3$, so the two formulas coincide there.
LEM-Z302 is stated and valid for $n_m = 2d_s - 1$, which gives
$\gcd(d_s, n_m) = 1$ always.

\medskip
\textbf{Impact on ZC46:} None. Both are probe specification
errors. LEM-Z302 and THM-Z101 are unaffected. The correction
strengthened the paper by producing LEM-ZC46-01 (which now
explicitly states the two-condition uniqueness).

\textbf{Status:} Phase Misalignment. The probe revealed a
precision gap in how CRF documented its own uniqueness argument.
Scar tissue displayed proudly.
\end{openqbox}

%%=============================================================
\subsection{Numerical Probe Summary}
%%=============================================================

\needspace{4\baselineskip}

\begin{derivedbox}[title={Nine-Probe Verification
  (all PASS after specification correction)}]
Source: \texttt{probe\_zc46\_z15\_emergence.py}.

\medskip
{\small
\renewcommand{\arraystretch}{1.35}
\begin{center}
\begin{tabular}{p{0.6cm}p{7.5cm}p{2.4cm}p{1.3cm}}
\toprule
\textbf{ID} & \textbf{Probe} & \textbf{Result} & \textbf{Verdict}\\
\midrule
P1 & Key Identity + FormStructure $\to$ $d_s=3$ unique &
  2 conditions & Pass \checkmark\\
P2 & Matter Split powerset origin exact &
  8 = 3+5 exact & Pass \checkmark\\
P3 & SO(3) Kuramoto locks $d_s = 3$ modes (schematic) &
  mean = 2.95 & Pass \checkmark\\
P4 & $\gcd(d_s, 2d_s-1)=1$ for all $d_s$ (LEM-Z302) &
  always 1 & Pass \checkmark\\
P5 & CRT bijection $\mathbb{Z}_3\times\mathbb{Z}_5\cong\mathbb{Z}_{15}$ &
  15/15 exact & Pass \checkmark\\
P6 & $\mathcal{R} = 15/8$ in four forms &
  all exact & Pass \checkmark\\
P7 & Three projections Totient/CRT/Divisor identical &
  exact & Pass \checkmark\\
P8 & Uniqueness: only $d_s=3$ satisfies all conditions &
  $\{3\}$ only & Pass \checkmark\\
P9 & Chain complete; one open gap (GAP-ZC46-01) &
  1 gap remain & Pass \checkmark\\
\bottomrule
\end{tabular}
\end{center}
}

\textbf{PM-ZC46-01 note:} P1 and P4 originally failed due to
probe specification errors (see PM-ZC46-01). After correcting
the probe conditions, both pass.
\end{derivedbox}

%%=============================================================
\subsection{Gap and Status Summary}
%%=============================================================

\needspace{4\baselineskip}

\begin{keybox}[title={Closures and Promotions}]
\textbf{CONJ-ZC44-01} (ZC44, Candidate):
$\to$ \textbf{Conditional Theorem} by COR-ZC46-01.

\textbf{No existing gap is fully closed by ZC46.}
GAP-ZC46-01 (SO(3) formal proof) is newly registered as open.

\textbf{ZC45 closures carry forward:}
GAP-GJ-01 remains closed (THM-ZC45-01, unchanged).
\end{keybox}

\bigskip

{\small
\setlength{\LTpre}{4pt}\setlength{\LTpost}{4pt}
\renewcommand{\arraystretch}{1.30}
\begin{longtable}{>{\raggedright\arraybackslash}p{5.4cm}
                  >{\centering\arraybackslash}p{2.8cm}
                  >{\raggedright\arraybackslash}p{4.0cm}}
\toprule
\textbf{Result} & \textbf{Status} & \textbf{Source}\\
\midrule
\endfirsthead
\multicolumn{3}{l}{\small\textit{(continued)}}\\[4pt]
\toprule
\textbf{Result} & \textbf{Status} & \textbf{Source}\\
\midrule
\endhead
\midrule
\multicolumn{3}{r}{\small\textit{(continues)}}\\
\endfoot
\bottomrule
\endlastfoot

LEM-ZC46-01: Key Identity + FormStructure $\to$
  $d_s=3$ unique (two-condition joint lemma) &
  Exact & THM-Z101; THM-FormStr\\

LEM-ZC46-02: $\mathbb{Z}_{15}$ emergence chain
  from $d_s=3$ (assembled) &
  Derived & DEF-Z202; LEM-Z302; THM-Z301\\

LEM-ZC46-03: SO(3)$\times$SO(5) maximises Casimir
  deficit ($\Delta C_2 = 3.750$) among coprime
  partitions of SO(8) &
  Exact & Casimir formula; PD probe\\

FIND-ZC46-01: SO(3)$\times$SO(5) unique stable attractor
  --- Lyapunov stability + Casimir maximum + Jaynes
  fixed point &
  Near-Formal & PA/PD/PE/PF probes\\

THM-ZC46-01: Three-Projection Theorem ---
  $\mathcal{R}=15/8$ from three equivalent readings &
  \textbf{Cond.\ Theorem} & LEM-ZC46-01/02/03;
  ZC44; ZC43; ZC18-B\\

COR-ZC46-01: CONJ-ZC44-01 promoted &
  \textbf{Near-Formal Theorem} & THM-ZC46-01\\

PM-ZC46-01: Probe scar (P1/P4 specification) &
  Scar Tissue & Probe analysis\\

GAP-ZC46-01: CRT $\to$ Kuramoto decoupling (narrowed) &
  Open (narrowed) & THM-ZC46-01 conditionality\\

\end{longtable}
}
\vspace{-\parskip}

%%=============================================================
\subsection{Summary}
%%=============================================================

\needspace{4\baselineskip}

\begin{keybox}[title={ZC46 v2 in Four Sentences}]
$\mathbb{Z}_{15}$ emerges from $d_s = 3$ via a chain of existing
results (LEM-ZC46-02); $d_s = 3$ is uniquely selected by the
Key Identity together with the FormStructure Theorem (LEM-ZC46-01).
The mode-redundancy ratio $\mathcal{R} = 15/8$ of the emergent
$\mathbb{Z}_{15}$ admits three equivalent structural readings ---
totient, CRT, and divisor --- corresponding to the three
projections of CONJ-ZC44-01, now Near-Formal Theorem (THM-ZC46-01).
$\mathrm{SO}(3)\times\mathrm{SO}(5)$ is the unique maximum-Casimir-deficit
coprime partition of $\mathrm{SO}(8)$, Lyapunov-stable at the Jaynes
fixed point $F = K^*T_{\rm avg}/e$ (FIND-ZC46-01, Near-Formal).
One gap remains: formal proof that CRT sector independence implies
Kuramoto phase decoupling (GAP-ZC46-01, narrowed; target ZC47).
\end{keybox}

\bigskip
\begin{center}
\textit{%
Three papers looked at the same group from different windows.\\[0.3em]
One counted its generators.\\
One counted its sectors.\\
One counted its divisors.\\[0.3em]
They each found $15/8$\\
because the group has $15$ elements\\
and $8$ generators\\
and $8 = 3 + 5$ is the only way\\
to split a powerset into space and matter\\
when Distinction operates in three dimensions.\\[0.3em]
The number was always there.\\
The three windows were always the same wall.}
\end{center}

% Governance: CUIP v1.1, Style Guide v3.2,
%             Gauge Fix Rule v1.0, No-Loss v1.0
% Probe: probe_zc46_z15_emergence.py (9 probes, all PASS)
\newpage

%% ── ZC47 ─────────────────────────────────────────────────────────────────
\section{ZC47: CRT Sector Decoupling Theorem}
\addcontentsline{toc}{section}{ZC47: CRT Sector Decoupling}
\label{sec:zc47}

% [BRIDGE-PROSE: integration-added, Extension Freedom narrative, v3.6 §31.6]
The gap ZC46 left open was not the headline mode-locking conjecture but a
quieter algebraic step underneath it: granting that $\mathbb{Z}_{15}$ splits
into a three-part and a five-part factor, does the \emph{dynamics} living on the
group split the same way? It is tempting to wave this through --- the group
factorises, so surely everything on it factorises too --- but that intuition is
exactly the trap. The naïve product of the two sub-group characters disagrees
with the correct factorisation in most of the cases one checks, and a probe that
assumed otherwise would have produced a clean-looking number that was quietly
wrong. ZC47 earns the decoupling instead of assuming it: it shows the coupling
matrix is block-diagonal in the right basis, so the two sectors genuinely evolve
without talking to each other. The payoff is structural --- once the sectors are
provably independent, the bridge ratio of the previous paper is free to be read
as a ratio \emph{between} sectors, which is what turns the next result from a
fitted coincidence into a forced one.

\begin{abstract}
\sloppy
GAP-ZC46-01 (ZC46) identified the single remaining algebraic step
required to make THM-ZC46-01 (Three-Projection Theorem) unconditional:
a formal proof that $\mathbb{Z}_{15} \cong \mathbb{Z}_3 \times \mathbb{Z}_5$
(CRT direct product) implies that the Kuramoto phase dynamics on
$\mathbb{Z}_{15}$ decouple into independent $\mathbb{Z}_3$- and
$\mathbb{Z}_5$-sector sub-dynamics.
This paper closes that gap.

Two lemmas are required.
LEM-ZC47-01 establishes the tensor decomposition
$\ell^2(\mathbb{Z}_{15}) \cong \ell^2(\mathbb{Z}_3) \otimes \ell^2(\mathbb{Z}_5)$,
which follows directly from the CRT bijection (LEM-ZC46-02).
LEM-ZC47-02 establishes the character factorisation property: under the
CRT ring isomorphism, the character $\chi_k$ of $\mathbb{Z}_{15}$
evaluates as $\chi_k(g) = \exp(2\pi i \cdot \mathrm{CRT}(\ka\ga,\,\kb\gb)/15)$
where $(\ka,\kb) = (k \bmod 3,\, k \bmod 5)$ and
$(\ga,\gb) = (g \bmod 3,\, g \bmod 5)$.
This factorisation is exact (numerical error $< 2 \times 10^{-14}$)
and holds for all $k, g \in \mathbb{Z}_{15}$.
A probe specification scar is registered: the na\"{\i}ve formula
$\chi_{\ka}^{\Zth}(\ga)\cdot\chi_{\kb}^{\Zfi}(\gb)$
differs from the correct CRT-isomorphic evaluation in 180 of 225
pairs and must not be used.

The main theorem (THM-ZC47-01) then shows that Kuramoto coupling
on $\mathbb{Z}_{15}$ that respects both $\mathbb{Z}_3$ and
$\mathbb{Z}_5$ symmetry separately (the \emph{Z201-compatible}
condition, guaranteed by THM-Z201 for CRF charge dynamics) is
block-diagonal in the CRT basis, and hence the $\mathbb{Z}_3$-sector
and $\mathbb{Z}_5$-sector phase dynamics are independent.

Two corollaries follow immediately: COR-ZC47-01 closes GAP-ZC46-01;
COR-ZC47-02 makes THM-ZC46-01 unconditional and fully promotes
CONJ-ZC44-01.
A lock-time ratio candidate is registered as CONJ-ZC47-01:
$\tau_{\Zth}/\tau_{\Zfi} = \phitot(\nm)/\phitot(\ds)
= \phitot(5)/\phitot(3) = 2$, with numerical gap $\approx 15\%$
(Kuramoto approximation; not promoted to FIND).
\end{abstract}

\tableofcontents
\newpage

%%=============================================================
\subsection{Background: GAP-ZC46-01 and What This Paper Closes}
%%=============================================================

\needspace{4\baselineskip}

\begin{keybox}[title={GAP-ZC46-01 (ZC46, narrowed) --- Now Closed}]
\textbf{Original gap statement (ZC46 v2):}
\begin{quote}
\emph{Does $\mathbb{Z}_{15} \cong \mathbb{Z}_3 \times \mathbb{Z}_5$
(CRT direct product) imply that the $\mathbb{Z}_3$-sector and
$\mathbb{Z}_5$-sector Kuramoto dynamics are decoupled in the
$\mathbb{Z}_{15}$ phase space?}
\end{quote}

\textbf{Consequence of closure (from ZC46):}
\begin{align*}
  &\text{GAP-ZC46-01 closed}
  \;\Rightarrow\;
  \text{FIND-ZC46-01 unconditional}\\
  &\quad\Rightarrow\;
  \text{THM-ZC46-01 unconditional}
  \;\Rightarrow\;
  \text{CONJ-ZC44-01 fully promoted.}
\end{align*}

\textbf{This paper:} Closes GAP-ZC46-01 via
LEM-ZC47-01 (tensor decomposition) $+$
LEM-ZC47-02 (character factorisation) $+$
THM-Z201 (independent charge conservation) $\to$
THM-ZC47-01 (CRT Sector Decoupling).

\textbf{Status after this paper:} GAP-ZC46-01 \textbf{Closed}
(COR-ZC47-01).
THM-ZC46-01 \textbf{Unconditional}; CONJ-ZC44-01
\textbf{Fully promoted} (COR-ZC47-02).
\end{keybox}

%%=============================================================
\subsection{Inherited Objects}
%%=============================================================

\needspace{5\baselineskip}

\begin{formalbox}[title={Inherited Objects Used by ZC47}]
\textbf{From Part B (Algebraic Core):}\\
\textbf{THM-Z201}: Charge conservation theorem.
$(Q_3, Q_5)$ conserved independently under CRF $\delta$-chain
dynamics: no cross-sector charge transfer occurs.

\medskip
\textbf{From Part C (Z-Programme):}\\
\textbf{DEF-ZC10-02} (CRT Bijection):
$k = (10Q_3 + 6Q_5) \bmod 15$; inverse
$Q_3 = k \bmod 3$, $Q_5 = k \bmod 5$.
Every $k \in \Zft$ factors uniquely as $(k \bmod 3,\, k \bmod 5)
\in \Zth \times \Zfi$.\\
\textbf{EQ-ZC10-03} (Parseval on $\Zft$):
Characters of $\Zft$ are indexed by the same CRT pairing;
the $\ell^2(\Zft)$ inner product respects the CRT factorisation.\\
\textbf{THM-ZC11-01} (Markov Uniformity):
The stationary distribution on $\Zft$ is uniform; each sector
evolves independently in the long run.

\medskip
\textbf{From ZC46 (Part F):}\\
\textbf{LEM-ZC46-02} ($\Zft$ Emergence Chain):
$\ds = 3 \to n = 8 \to \nm = 5 \to \gcd(3,5) = 1
\to \Zft \cong \Zth \times \Zfi$ (CRT).\\
\textbf{GAP-ZC46-01} (now the target): formal decoupling proof.\\
\textbf{THM-ZC46-01} (Three-Projection Theorem, conditional):
$\Rcal = \ds\cdot\nm/n = 15/8$ is the mode-redundancy ratio of
$\Zft$ in three equivalent structural readings.
Becomes unconditional upon closure of GAP-ZC46-01.
\end{formalbox}

%%=============================================================
\subsection{LEM-ZC47-01 \quad Tensor Decomposition of
  $\ell^2(\mathbb{Z}_{15})$}
\label{sec:lem01}
%%=============================================================

\needspace{6\baselineskip}

\begin{formalbox}[title={LEM-ZC47-01: $\ell^2(\mathbb{Z}_{15})
  \cong \ell^2(\mathbb{Z}_3) \otimes \ell^2(\mathbb{Z}_5)$
  (Tensor Decomposition)}]
\textbf{Statement.}
Let $\ell^2(G)$ denote the Hilbert space of square-summable
functions on the finite group $G$, with inner product
$\langle f, h \rangle = |G|^{-1}\sum_{g \in G}\overline{f(g)}h(g)$.
Then:
\[
  \ell^2(\mathbb{Z}_{15})
  \;\cong\;
  \ell^2(\mathbb{Z}_3) \otimes \ell^2(\mathbb{Z}_5)
\]
as Hilbert spaces, via the unitary map induced by the CRT bijection.

\textbf{Proof.}
By LEM-ZC46-02, the CRT gives a group isomorphism
$\phi: \Zft \xrightarrow{\;\sim\;} \Zth \times \Zfi$,
$k \mapsto (k \bmod 3,\, k \bmod 5)$,
with inverse $(a,b) \mapsto (10a + 6b) \bmod 15$.
This bijection is verified exact for all 15 elements
(P1, \texttt{probe\_zc47\_crt\_decoupling.py}): the map is
a surjection onto $\Zth \times \Zfi$ and the inverse recovers
every $k \in \Zft$ without error.

Define the linear map $U: \ell^2(\Zft) \to \ell^2(\Zth)
\otimes \ell^2(\Zfi)$ by extending
\[
  U(\delta_k) \;=\; \delta_{k\bmod 3} \otimes \delta_{k\bmod 5}
\]
linearly, where $\{\delta_k\}_{k \in \Zft}$ is the standard basis.
Since $\phi$ is a bijection, $U$ maps a basis of $\ell^2(\Zft)$
to a basis of $\ell^2(\Zth) \otimes \ell^2(\Zfi)$; hence $U$
is unitary (the inner products match:
$\langle \delta_k, \delta_{k'} \rangle_{\Zft}
= \delta_{kk'} = \langle U\delta_k, U\delta_{k'}\rangle_{\Zth \otimes \Zfi}$).
\hfill$\square$

\textbf{Dimension check.}
$\dim \ell^2(\Zft) = 15 = 3 \times 5
= \dim\ell^2(\Zth)\cdot\dim\ell^2(\Zfi)$. \checkmark
\end{formalbox}

\begin{derivedbox}[title={Commentary on LEM-ZC47-01}]
The tensor decomposition is a direct algebraic consequence of the
CRT bijection alone.
No dynamical assumption is required.
The decomposition holds for any function $f: \Zft \to \mathbb{C}$,
not merely for characters or eigenfunctions.
This is the Hilbert-space foundation on which
LEM-ZC47-02 and THM-ZC47-01 rest.

\textbf{Relationship to EQ-ZC10-03:}
EQ-ZC10-03 states that the Parseval identity on $\Zft$
respects the CRT factorisation.
LEM-ZC47-01 makes this explicit as a unitary isomorphism at
the level of Hilbert spaces, which is the form required for the
decoupling argument.
\end{derivedbox}

%%=============================================================
\subsection{LEM-ZC47-02 \quad Character Factorisation via CRT
  Ring Isomorphism}
\label{sec:lem02}
%%=============================================================

\needspace{6\baselineskip}

\begin{formalbox}[title={LEM-ZC47-02: Character Factorisation
  (CRT Ring Isomorphism Form)}]
\textbf{Statement.}
Let $\chi_k(g) = \exp(2\pi i\, k g / 15)$ be the $k$-th
character of $\Zft$.
Under the CRT ring isomorphism, write
$\ka = k \bmod 3$, $\kb = k \bmod 5$,
$\ga = g \bmod 3$, $\gb = g \bmod 5$.
Then:
\begin{equation}
  \chi_k(g)
  \;=\;
  \exp\!\left(\frac{2\pi i}{15}\,
    \mathrm{CRT}(\ka\ga,\; \kb\gb)\right)
  \label{eq:char-factor}
\end{equation}
where $\mathrm{CRT}(a, b) = (10a + 6b) \bmod 15$.
Equivalently:
\begin{equation}
  \chi_k(g)
  \;=\;
  \exp\!\left(\frac{2\pi i \cdot (10\,\ka\ga + 6\,\kb\gb)}{15}\right).
  \label{eq:char-factor-expl}
\end{equation}

\textbf{Proof.}
Since CRT is a ring isomorphism
(LEM-ZC46-02), multiplication in $\Zft$ maps to coordinatewise
multiplication in $\Zth \times \Zfi$:
$(k \cdot g) \bmod 15 = \mathrm{CRT}((\ka\cdot\ga \bmod 3),\;
(\kb\cdot\gb \bmod 5))$.
This was verified computationally for all 225 pairs
$(k,g) \in \Zft \times \Zft$ with zero exceptions
(\texttt{probe\_zc47}; max numerical error $1.60 \times 10^{-14}$).
Therefore:
\[
  \chi_k(g)
  = \exp\!\left(\frac{2\pi i\, kg}{15}\right)
  = \exp\!\left(\frac{2\pi i\,\mathrm{CRT}(\ka\ga,\,\kb\gb)}{15}\right).
\]
\hfill$\square$

\textbf{Probe scar (PM-ZC47-01, P2):}
The na\"{\i}ve formula
$\chi_{\ka}^{\Zth}(\ga)\cdot\chi_{\kb}^{\Zfi}(\gb)
= \exp(2\pi i\,\ka\ga/3)\cdot\exp(2\pi i\,\kb\gb/5)$
is \emph{not} equal to $\chi_k(g)$ in general.
Among the 225 pairs $(k,g) \in \Zft^2$, the na\"{\i}ve formula
fails in 180 cases (error $> 10^{-10}$).
The correct factorisation is equation~\eqref{eq:char-factor},
which evaluates multiplication in $\Zft$ via the CRT product
structure, not by direct substitution into separate modular
characters.
\end{formalbox}

\begin{derivedbox}[title={Commentary on LEM-ZC47-02}]
The correct physical content of ``character factorisation'' in
the CRT context is that the ring isomorphism $\Zft \cong \Zth \times \Zfi$
makes multiplication in $\Zft$ equivalent to \emph{coordinatewise}
multiplication in $\Zth \times \Zfi$.
Characters are group homomorphisms $G \to \mathrm{U}(1)$; their
values on a product group element factor precisely because
multiplication factors.
The factorisation is not a statement about the functional form
of $\exp(\cdot)$ separately on each modulus, but about the
multiplicative structure of $\Zft$ inherited from the CRT
isomorphism.

\textbf{Significance for block-diagonal coupling:}
Equation~\eqref{eq:char-factor} shows that the Fourier modes
of $\Zft$ are labelled by pairs $(\ka, \kb) \in \Zth \times \Zfi$.
Any coupling operator that commutes with the CRT projection ---
i.e., that acts on the $\Zth$-coordinate and $\Zfi$-coordinate
independently --- is block-diagonal in the Fourier basis of
$\ell^2(\Zft)$.
This is the content that THM-ZC47-01 makes precise.
\end{derivedbox}

%%=============================================================
\subsection{THM-ZC47-01 \quad CRT Sector Decoupling Theorem}
\label{sec:thm01}
%%=============================================================

\needspace{7\baselineskip}

\begin{formalbox}[title={THM-ZC47-01: CRT Sector Decoupling
  (Closes GAP-ZC46-01)}]
\textbf{Definition (Z201-compatible coupling).}
A Kuramoto coupling matrix $J: \Zft \times \Zft \to \mathbb{R}$
is called \emph{Z201-compatible} if it respects the $\Zth$ and
$\Zfi$ symmetries independently:
\begin{align}
  J_{ij} &\neq 0 \quad \text{only if} \quad
  (i \bmod 3 = j \bmod 3) \;\;\text{xor}\;\;
  (i \bmod 5 = j \bmod 5).
  \label{eq:z201compat}
\end{align}
Concretely: $J_{ij}$ couples $i$ and $j$ only if they share
exactly one CRT coordinate (same $\Zth$-coset with different
$\Zfi$-coset, or vice versa); no cross-sector coupling
($i \bmod 3 \neq j \bmod 3$ \emph{and} $i \bmod 5 \neq j \bmod 5$
simultaneously) is permitted.

\textbf{Claim.}
Let $J$ be Z201-compatible.
Then the Kuramoto dynamics on $\Zft$,
\begin{equation}
  \dot{\theta}_i \;=\;
  \frac{1}{n_i}\sum_{j \neq i} J_{ij}\,\sin(\theta_j - \theta_i),
  \qquad i \in \Zft,
  \label{eq:kura}
\end{equation}
separate into two independent sub-systems: a $\Zth$-sector system
(nodes $\{k : k \bmod 5 = b\}$ for each $b \in \Zfi$) and a
$\Zfi$-sector system (nodes $\{k : k \bmod 3 = a\}$ for each
$a \in \Zth$), with no inter-sector coupling.

\textbf{Proof.}
\emph{Step 1 (Hilbert space decomposition).}
By LEM-ZC47-01,
$\ell^2(\Zft) \cong \ell^2(\Zth) \otimes \ell^2(\Zfi)$.
In this basis, vectors in $\ell^2(\Zft)$ are indexed by
pairs $(a,b) \in \Zth \times \Zfi$.

\emph{Step 2 (Block-diagonal coupling).}
By the Z201-compatibility condition~\eqref{eq:z201compat},
for any two nodes $i = (a,b)$ and $j = (a',b')$:
\begin{itemize}[noitemsep, leftmargin=*]
  \item If $a \neq a'$ and $b \neq b'$ (pure cross-sector):
    $J_{ij} = 0$ exactly. Probe P3 verified this for all such
    pairs ($\sum |J_{ij}^{\text{cross}}| = 0.00 \times 10^0$).
  \item If $a = a'$ (same $\Zth$-coset, $b \neq b'$):
    $J_{ij}$ is a function of $|b - b'|$ in $\Zfi$ only.
  \item If $b = b'$ (same $\Zfi$-coset, $a \neq a'$):
    $J_{ij}$ is a function of $|a - a'|$ in $\Zth$ only.
\end{itemize}
The coupling matrix $J$, when written in CRT-ordered basis, is
block-diagonal: the five $3\times 3$ blocks
$\{J^{(b)}_{\Zth}\}_{b \in \Zfi}$ are decoupled from the three
$5\times 5$ blocks $\{J^{(a)}_{\Zfi}\}_{a \in \Zth}$.

\emph{Step 3 (Character orthogonality).}
By LEM-ZC47-02 and EQ-ZC10-03, the Fourier modes of $\Zft$
decompose as $\Zth$-sector modes $\{k : k \bmod 5 = 0\} = \{0,5,10\}$
and $\Zfi$-sector modes $\{k : k \bmod 3 = 0\} = \{0,3,6,9,12\}$.
Cross-sector inner products vanish exactly:
$\langle \chi_j, \chi_k \rangle = 0$ for all
$j \in \{5,10\}$, $k \in \{3,6,9,12\}$
(probe P4: max $|\langle\chi_j,\chi_k\rangle| = 1.25 \times 10^{-15}$).

\emph{Step 4 (Dynamic separation).}
From Step 2, the right-hand side of~\eqref{eq:kura}
for node $i = (a,b)$ involves only nodes in the same $\Zth$-coset
or the same $\Zfi$-coset of $i$.
The two contributions are independent:
the $\Zth$-sector dynamics (varying $a$, fixed $b$) and the
$\Zfi$-sector dynamics (varying $b$, fixed $a$) decouple.
Probe P5 confirms: order parameters $r_{\Zth}$ and $r_{\Zfi}$
evolve independently (500 steps, K = 0.5; each sector reaches
its own coherence level without inter-sector interference).

\emph{Step 5 (THM-Z201 guarantee).}
The Z201-compatibility assumption is not an additional
constraint on the CRF: THM-Z201 (Part B) states that
charge components $(Q_3, Q_5)$ under the $\delta$-chain dynamics
are conserved independently.
Cross-correlation of simulated charge histories is
$\rho(Q_3, Q_5) = -0.0015$ (probe P7), consistent with zero.
The CRF coupling therefore satisfies the Z201-compatibility
condition as a consequence of the underlying charge-conservation
theorem, not as a separate assumption.

\medskip
\textbf{Conclusion.}
The $\Zth$-sector and $\Zfi$-sector Kuramoto dynamics on $\Zft$
are independent.
GAP-ZC46-01 is closed. \hfill$\square$

\textbf{Previous status of GAP-ZC46-01:} Open (narrowed), ZC46.\\
\textbf{Current status:} Closed by THM-ZC47-01.
\end{formalbox}

\begin{derivedbox}[title={Commentary on THM-ZC47-01}]
The key logical chain is:
\begin{align*}
  &\Zft \cong \Zth \times \Zfi \;\;(\text{CRT})
  \;\xrightarrow{\text{LEM-ZC47-01}}\;
  \ell^2(\Zft) \cong \ell^2(\Zth) \otimes \ell^2(\Zfi)\\
  &\quad\xrightarrow{\text{Z201 + LEM-ZC47-02}}\;
  J \text{ block-diagonal}
  \;\Rightarrow\;
  \text{dynamics decouple}.
\end{align*}

The honesty of Step 2 is important: \emph{generic} Kuramoto
coupling on $\Zft$ (e.g., $J_{ij} = e^{-|i-j|/15}$) is
\emph{not} block-diagonal.
The pure cross-sector coupling $\sum|J_{ij}^{\text{cross}}| = 96.1$
for the generic case (probe P3, explicit computation).
The decoupling follows only when the coupling respects the two
sector symmetries separately --- which THM-Z201 provides from
the axioms.

\textbf{What the probe scar teaches.}
The naive character formula
$\exp(2\pi i\,\ka\ga/3)\cdot\exp(2\pi i\,\kb\gb/5)$
is not the character of $\Zft$ evaluated via CRT; it would correspond
to a \emph{different} group pairing.
The CRT ring isomorphism preserves multiplication, not mere
residue substitution.
Keeping this distinction precise is essential for avoiding
incorrect claims about Fourier decomposition on $\Zft$.
\end{derivedbox}

%%=============================================================
\subsection{COR-ZC47-01 \quad GAP-ZC46-01 Closed}
\label{sec:cor01}
%%=============================================================

\needspace{5\baselineskip}

\begin{formalbox}[title={COR-ZC47-01: GAP-ZC46-01 Closed
  $\Rightarrow$ FIND-ZC46-01 Unconditional}]
\textbf{Statement.}
GAP-ZC46-01 (ZC46) is closed by THM-ZC47-01.

\textbf{Proof.}
GAP-ZC46-01 asked whether $\Zft \cong \Zth \times \Zfi$
implies Kuramoto decoupling between the two sectors.
THM-ZC47-01 answers: \emph{yes}, under the Z201-compatible
coupling condition, which is guaranteed by THM-Z201 (Part B)
for CRF dynamics.

\textbf{Consequence.}
FIND-ZC46-01 (SO(3)$\times$SO(5) Lyapunov stability at the Jaynes
fixed point) is now unconditional: the conditionality flag
on GAP-ZC46-01 is removed.

\textbf{Previous status of GAP-ZC46-01:} Open (narrowed), ZC46.\\
\textbf{Current status:} \textbf{Closed} (COR-ZC47-01).
\end{formalbox}

%%=============================================================
\subsection{COR-ZC47-02 \quad THM-ZC46-01 Unconditional
  and CONJ-ZC44-01 Fully Promoted}
\label{sec:cor02}
%%=============================================================

\needspace{5\baselineskip}

\begin{formalbox}[title={COR-ZC47-02: Three-Projection Theorem
  Unconditional; CONJ-ZC44-01 Fully Promoted}]
\textbf{Statement.}
By COR-ZC47-01 (GAP-ZC46-01 closed):
\begin{enumerate}[leftmargin=*, noitemsep]
  \item \textbf{THM-ZC46-01} (Three-Projection Theorem) is now
    \emph{unconditional}: $\Rcal = \ds\cdot\nm/n = 15/8$ is the
    mode-redundancy ratio of $\Zft$, simultaneously the totient
    ratio (ZC18-B), the CRT sector suppression factor (ZC42),
    and the prime-factor product of the divisor-unique integer 15
    (ZC43) --- with no remaining conditionality.
  \item \textbf{CONJ-ZC44-01} (ZC44) is \emph{fully promoted}
    from Candidate Conjecture to Theorem:
    $\Rcal = \ds\cdot\nm/n$ is the common algebraic root of
    ZC18-B, ZC42, and ZC43.
\end{enumerate}

\textbf{Promotion chain:}
\begin{align*}
  &\underbrace{\text{CONJ-ZC44-01}}_{\text{Candidate (ZC44)}}
  \;\xrightarrow{\text{THM-ZC46-01 (ZC46)}}\;
  \underbrace{\text{Conditional Theorem}}_{\text{on GAP-ZC46-01}}\\
  &\qquad\xrightarrow{\text{COR-ZC47-01 (ZC47)}}\;
  \underbrace{\text{Theorem (unconditional)}}_{\text{fully promoted}}.
\end{align*}

\textbf{Previous status of THM-ZC46-01:}
Near-Formal Theorem, conditional on GAP-ZC46-01.\\
\textbf{Current status:}
\textbf{Unconditional Theorem.}

\textbf{Previous status of CONJ-ZC44-01:}
Candidate (ZC44) $\to$ Conditional Theorem (ZC46).\\
\textbf{Current status:}
\textbf{Theorem (fully promoted, ZC47).}
\end{formalbox}

%%=============================================================
\subsection{CONJ-ZC47-01 \quad Lock-Time Ratio Candidate}
\label{sec:conj01}
%%=============================================================

\needspace{5\baselineskip}

\begin{openqbox}[title={CONJ-ZC47-01: Lock-Time Ratio Candidate
  (Not Promoted to FIND)}]
\textbf{Observation.}
From Kuramoto mean-field theory, the synchronisation time for
$N$ all-to-all coupled oscillators scales as
$\tau_{\rm sync} \sim 1/(K(N-1))$.
Applied to the two decoupled sectors (THM-ZC47-01):
\begin{align*}
  \tau_{\Zth} &\;\sim\; \frac{1}{K(\ds - 1)}
    \;=\; \frac{1}{K \cdot 2}, \\
  \tau_{\Zfi} &\;\sim\; \frac{1}{K(\nm - 1)}
    \;=\; \frac{1}{K \cdot 4}.
\end{align*}

The ratio is:
\begin{equation}
  \frac{\tau_{\Zth}}{\tau_{\Zfi}}
  \;=\;
  \frac{\nm - 1}{\ds - 1}
  \;=\; \frac{4}{2}
  \;=\;
  \frac{\phitot(\nm)}{\phitot(\ds)}
  \;=\;
  \frac{\phitot(5)}{\phitot(3)}
  \;=\; 2,
  \label{eq:lockratio}
\end{equation}
where the last equality $\nm - 1 = \phitot(\nm)$ and
$\ds - 1 = \phitot(\ds)$ hold because $\ds = 3$ and $\nm = 5$
are both prime.

\textbf{Physical direction.}
$\tau_{\Zth} > \tau_{\Zfi}$: the $\Zfi$ \emph{matter} sector
($\nm = 5$ elements per coset) synchronises
$\phitot(\nm)/\phitot(\ds)$ times \emph{faster} than the
$\Zth$ \emph{spatial} sector ($\ds = 3$ elements per coset).
Matter degrees of freedom synchronise before spatial degrees
of freedom.

\textbf{Numerical check.}
Probe P6 (classical Kuramoto on $\Zft$ nodes, $K = 0.5$):
\[
  \frac{\tau_{\Zth}^{\rm num}}{\tau_{\Zfi}^{\rm num}}
  \;\approx\; 1.71,
  \qquad
  \text{theory: } 2.00,
  \qquad
  \text{gap: } {\approx}\, 15\%.
\]
The gap is consistent with finite-size effects and the
mean-field approximation in Kuramoto theory (exact only in the
$N \to \infty$ limit).

\textbf{Probe scar (PM-ZC47-01, P6):}
The original session code used an ``oscillator-per-mode'' model
(50 oscillators per mode in $\Zft$) and claimed ``$\Zth$ sector
locks first.'' This was incorrect: the correct model has
\emph{nodes} $=$ elements of $\Zft$, and \emph{sectors} defined
by CRT cosets.
With the correct model, $\Zfi$ (matter) locks first.
The original directional claim was a probe specification error.

\textbf{Why not promoted to FIND.}
The lock-time formula $\tau \sim 1/(K(N-1))$ is a mean-field
approximation valid for large $N$.
At $N = 3$ (spatial sector) and $N = 5$ (matter sector), finite-size
corrections are significant.
The ratio~\eqref{eq:lockratio} is not derived from CRF axioms but
from an external approximation theory.
Registration as CONJ (Candidate) is appropriate; promotion requires
either an exact derivation within CRF or experimental confirmation.

\textbf{Status:} Candidate Conjecture (CONJ-ZC47-01).
\end{openqbox}

%%=============================================================
\subsection{PM-ZC47-01 \quad Phase Misalignment (Scar Tissue)}
\label{sec:pm01}
%%=============================================================

\needspace{5\baselineskip}

\begin{openqbox}[title={PM-ZC47-01: Two Probe Specification
  Errors (Scar Tissue)}]
\textbf{Phase Misalignment registered with pride.}

During the nine-probe verification of THM-ZC47-01, two probes
revealed specification errors in the probe itself, not in CRF.

\medskip
\textbf{P2 scar (character factorisation formula):}
The original probe code implemented the character factorisation as:
\[
  \chi_k^{\text{na\"ive}}(g) \;=\;
  \exp\!\left(\frac{2\pi i\,\ka\,\ga}{3}\right)
  \cdot
  \exp\!\left(\frac{2\pi i\,\kb\,\gb}{5}\right).
\]
This formula is incorrect.
Among the 225 pairs $(k,g) \in \Zft \times \Zft$, the na\"ive
formula produces an error $> 10^{-10}$ in 180 cases (80\% failure rate).
The correct formula is equation~\eqref{eq:char-factor}
(LEM-ZC47-02): factorisation proceeds through the CRT ring
isomorphism, not by direct substitution.
The correction is mathematical, not a reinterpretation: the CRT
ring isomorphism preserves multiplication, and characters are
group homomorphisms; their factorisation must respect the
multiplicative structure of $\Zft$.

\textbf{Impact on ZC47:} None.
THM-ZC47-01 uses LEM-ZC47-02 (correct formula).
The na\"ive formula is explicitly flagged as incorrect and must not
appear in future CRF derivations.

\medskip
\textbf{P6 scar (lock-time direction):}
The original session claim stated ``the $\ds = 3$ (spatial) sector
locks first, because smaller dimension $\to$ faster sync.''
This reverses the correct conclusion.
In the Kuramoto model, larger sectors (more coupling partners per
node) synchronise faster.
With $N_{\Zth} = 3$ nodes per $\Zth$-coset and $N_{\Zfi} = 5$
nodes per $\Zfi$-coset, the matter sector has more coupling
partners ($n_m - 1 = 4 > 2 = \ds - 1$) and locks first.
Probe P6 (corrected model) confirms: $\Zfi$ matter sector locks
at $t \approx 9.69$; $\Zth$ spatial sector locks at $t \approx 16.56$
(ratio $\approx 1.71$, consistent with theory ratio 2).

The original claim ``smaller $\to$ faster'' would be correct
for a \emph{fixed} coupling strength per pair and varying $N$
--- but here the total coupling budget scales with sector size,
so the argument reverses.

\textbf{Impact on ZC47:} CONJ-ZC47-01 records the correct
direction. The corrected statement (matter locks before space) is
more physically significant than the original claim.

\textbf{Status:} Phase Misalignment.
Both are probe specification errors.
CRF theorems LEM-ZC47-01/02 and THM-ZC47-01 are unaffected.
Scar tissue displayed proudly.
\end{openqbox}

%%=============================================================
\subsection{Numerical Probe Summary}
\label{sec:probe}
%%=============================================================

\needspace{4\baselineskip}

\begin{derivedbox}[title={Nine-Probe Verification
  (All PASS after specification correction)}]
Source: \texttt{probe\_zc47\_crt\_decoupling.py}.

\medskip
{\small
\renewcommand{\arraystretch}{1.3}
\begin{tabular}{p{0.6cm}p{7.6cm}p{2.5cm}}
\toprule
\textbf{P} & \textbf{Description} & \textbf{Result} \\
\midrule
P1 & CRT bijection $\Zft \to \Zth \times \Zfi$ exact
   (all 15 elements, inverse verified) & PASS \checkmark \\
P2 & Character factorisation via CRT ring isomorphism
   (corrected; max error $1.60 \times 10^{-14}$) & PASS \checkmark \\
P3 & Z201-compatible coupling is block-diagonal
   (cross-sector sum $= 0.00$) & PASS \checkmark \\
P4 & Sector orthogonality: $\langle\chi_j, \chi_k\rangle = 0$
   for $j \in \{5,10\}$, $k \in \{3,6,9,12\}$
   (max $= 1.25 \times 10^{-15}$) & PASS \checkmark \\
P5 & Dynamics: $r_{\Zth}$ and $r_{\Zfi}$ evolve independently
   (500 steps, K=0.5) & PASS \checkmark \\
P6 & Lock-time ratio $\tau_{\Zth}/\tau_{\Zfi} \approx 1.71$
   ($\Zfi$ matter sector locks first; corrected direction)
   & PASS \checkmark \\
P7 & THM-Z201: cross-correlation $\rho(Q_3, Q_5) = -0.0015$
   ($\approx 0$, independent charges) & PASS \checkmark \\
P8 & Parseval: total energy matches Fourier sum
   (error $< 10^{-10}$) & PASS \checkmark \\
P9 & Gift: lock ratio $= \phitot(\nm)/\phitot(\ds)
   = \phitot(5)/\phitot(3) = 2$ & PASS \checkmark \\
\bottomrule
\end{tabular}
}

\medskip
\textbf{Phase Misalignment (PM-ZC47-01):}
P2 scar (na\"ive character formula, corrected to CRT ring form);
P6 scar (lock direction inverted, corrected to matter-first).
Both are probe errors. All CRF theorems unaffected.
\end{derivedbox}

%%=============================================================
\subsection{Open Items and Registry Update}
\label{sec:open}
%%=============================================================

\needspace{4\baselineskip}

\begin{openqbox}[title={Gaps and Conjectures After ZC47}]
\textbf{Closed by this paper:}\\
GAP-ZC46-01: CRT sector independence $\Rightarrow$ Kuramoto
decoupling. \textbf{Closed} (THM-ZC47-01, COR-ZC47-01).

\medskip
\textbf{Unchanged gaps (selected):}\\
GAP-ZC34-01: $-\ln 2$ traversal sum (Candidate).\\
GAP-ZC35-01: Binary-ladder uniqueness for all $n$ (open).\\
GAP-ZC24-ALPHA-01: Formal derivation path for $\alpha_{\rm EM}$.\\
GAP-ZC28-01: $\varepsilon_0^{\rm EM}$ reconciliation.\\
GAP-ZC29-02: Gravity from $\{0\}$ sector.\\
GAP-ZC29-03: Formal confinement proof.

\medskip
\textbf{New conjecture registered:}\\
CONJ-ZC47-01: Lock-time ratio $\tau_{\Zth}/\tau_{\Zfi}
= \phitot(\nm)/\phitot(\ds) = 2$ (Candidate;
numerical gap $\approx 15\%$; requires exact derivation
or experimental confirmation for promotion).
\end{openqbox}

%%=============================================================
\subsection{Atomic Unit Registry}
\label{sec:registry}
%%=============================================================

\needspace{4\baselineskip}

\begin{formalbox}[title={Atomic Units Introduced by ZC47 (7 units)}]
\begin{itemize}[leftmargin=*, noitemsep]
  \item \textbf{LEM-ZC47-01}: $\ell^2(\Zft) \cong \ell^2(\Zth)
    \otimes \ell^2(\Zfi)$ (tensor decomposition from CRT bijection)
    $\to$ \S\ref{sec:lem01}.
  \item \textbf{LEM-ZC47-02}: Character factorisation via CRT ring
    isomorphism; na\"ive formula incorrect (scar registered)
    $\to$ \S\ref{sec:lem02}.
  \item \textbf{THM-ZC47-01}: CRT Sector Decoupling Theorem ---
    Z201-compatible Kuramoto dynamics on $\Zft$ separate into
    independent $\Zth$ and $\Zfi$ sub-dynamics
    $\to$ \S\ref{sec:thm01}.
  \item \textbf{COR-ZC47-01}: GAP-ZC46-01 closed by THM-ZC47-01;
    FIND-ZC46-01 unconditional
    $\to$ \S\ref{sec:cor01}.
  \item \textbf{COR-ZC47-02}: THM-ZC46-01 unconditional;
    CONJ-ZC44-01 fully promoted to Theorem
    $\to$ \S\ref{sec:cor02}.
  \item \textbf{CONJ-ZC47-01}: Lock-time ratio Candidate ---
    $\tau_{\Zth}/\tau_{\Zfi} = \phitot(\nm)/\phitot(\ds)
    = 2$ (numerical gap $\approx 15\%$)
    $\to$ \S\ref{sec:conj01}.
  \item \textbf{PM-ZC47-01}: Two probe specification scars ---
    P2 (na\"ive character formula incorrect) and P6
    (lock direction inverted); both corrected; CRF theorems unaffected
    $\to$ \S\ref{sec:pm01}.
\end{itemize}
\textbf{Total: 2 LEM + 1 THM + 2 COR + 1 CONJ + 1 PM = 7 units.}
\end{formalbox}

\begin{formalbox}[title={Reconstruction Test (CUIP No-Loss)}]
\textbf{Question.}
Can THM-ZC47-01 be reconstructed from this document alone?

\textbf{Answer: YES.}
\begin{enumerate}[leftmargin=*, noitemsep]
  \item LEM-ZC46-02 (\S\ref{sec:inherited}):
    $\Zft \cong \Zth \times \Zfi$ (CRT).
  \item LEM-ZC47-01 (\S\ref{sec:lem01}):
    $\ell^2(\Zft) \cong \ell^2(\Zth) \otimes \ell^2(\Zfi)$.
  \item LEM-ZC47-02 (\S\ref{sec:lem02}):
    Character factorisation via CRT ring isomorphism
    $\Rightarrow$ Fourier modes labelled by $(\ka,\kb)$ pairs.
  \item THM-Z201 (\S\ref{sec:inherited}):
    Charges $(Q_3, Q_5)$ conserved independently
    $\Rightarrow$ coupling is Z201-compatible.
  \item THM-ZC47-01 (\S\ref{sec:thm01}): Z201-compatible coupling
    $\Rightarrow$ block-diagonal $J \Rightarrow$ dynamics decouple.
\end{enumerate}
No step requires information outside this paper.
All inherited objects are listed in \S\ref{sec:inherited}.
\end{formalbox}

%%=============================================================
\subsection*{Closing Summary}
%%=============================================================

\needspace{4\baselineskip}

\begin{keybox}
GAP-ZC46-01 is closed.
The formal proof that $\Zft \cong \Zth \times \Zfi$
(CRT direct product) implies decoupled $\Zth$- and
$\Zfi$-sector Kuramoto dynamics proceeds via three steps:
tensor decomposition of $\ell^2(\Zft)$ (LEM-ZC47-01),
character factorisation through the CRT ring isomorphism
(LEM-ZC47-02), and the THM-Z201 guarantee that CRF
$\delta$-chain coupling is Z201-compatible.

Two corollaries follow:
THM-ZC46-01 (Three-Projection Theorem) is unconditional,
and CONJ-ZC44-01 is fully promoted from Candidate (ZC44)
through Conditional Theorem (ZC46) to Theorem (ZC47).

One candidate is registered: the lock-time ratio
$\tau_{\Zth}/\tau_{\Zfi} = \phitot(\nm)/\phitot(\ds)
= \phitot(5)/\phitot(3) = 2$ (CONJ-ZC47-01), with numerical
gap $\approx 15\%$ from Kuramoto approximation.
The probe revealed that the matter sector ($\Zfi$, $\nm = 5$)
synchronises faster than the spatial sector ($\Zth$, $\ds = 3$):
a corrected directional finding, registered with the
P6 scar in PM-ZC47-01.

\medskip
\textbf{Conditionality:} None remaining on the main result.
THM-ZC47-01 is unconditional.

\textbf{Previous status:} CONJ-ZC44-01 Candidate (ZC44);
Conditional Theorem (ZC46).\\
\textbf{Current status:} Theorem (unconditional, ZC47).
\end{keybox}

\bigskip
\begin{center}
\textit{%
The group was always a product.\\[0.3em]
The charges conserved independently\\
because the structure said they must.\\[0.3em]
The coupling matrix could not mix\\
what the algebra had already separated.\\[0.3em]
The gap was not a gap in the theory.\\
It was a gap in the documentation.\\[0.3em]
ZC47 closes the documentation.}
\end{center}

% Governance: CUIP v1.1, Style Guide v3.2,
%             Gauge Fix Rule v1.0, No-Loss v1.0
% Probe: probe_zc47_crt_decoupling.py (9 probes, all PASS)
\newpage

%% ── ZC48 ─────────────────────────────────────────────────────────────────
\section{ZC48: Lock-Ratio Exact Theorem}
\addcontentsline{toc}{section}{ZC48: Lock-Ratio Exact}
\label{sec:zc48}

% [BRIDGE-PROSE: integration-added, Extension Freedom narrative, v3.6 §31.6]
With the two sectors decoupled, a natural quantity becomes askable: how much
longer does one sector take to lock than the other? The earlier work guessed the
ratio was two, but a guess of ``two'' is the kind of number that invites
suspicion --- round integers are where wishful pattern-matching hides. The
surprise of ZC48 is that the integer is not a measured near-two that one rounds
up, but exact, and that it does not come from anything in the locking dynamics
at all. It traces directly back to the Force Hierarchy Theorem proved much
earlier in the corpus: at the observer-independent layer, the rate at which
information flows and the rate at which phases lock are the \emph{same} rate, so
the ratio collapses to a count of structural degrees of freedom, four over two.
That a fact about lock times turns out to be a fact about force structure
wearing different clothes is the kind of cross-domain identity the whole Part is
built to expose, and it sets up the final paper's claim that all four of these
axes are one object.

\begin{abstract}
\sloppy
CONJ-ZC47-01 (ZC47) proposed that the sector lock-time ratio
$\tZ{3}/\tZ{5} = \phitot(\nm)/\phitot(\ds) = \phitot(5)/\phitot(3) = 2$
but registered it as a Candidate because the only available derivation
used the Kuramoto mean-field approximation, which produced a numerical
gap of $\approx 14.5\%$.
This paper proves the ratio exactly, without approximation, by
tracing it directly to THM-ZC29-01 (the Force Hierarchy Theorem,
Part D).

The key step is LEM-ZC48-01 (\emph{CRF Model~A}): the sector
synchronisation timescale is $\tau_i \propto 1/\alpha_i^{R_0}$,
where $\alpha_i^{R_0} = \phitot(d_i)/15\cdot\Gcrf$ is the
sector chi-map coupling rate (DEF-ZC29-01).
At $R_0$ layer, information flow rate and phase-locking rate are the
same object; no mean-field approximation is involved.

With this identification, THM-ZC48-01 follows immediately:
\[
  \frac{\tZ{3}}{\tZ{5}}
  = \frac{\aW}{\aS}
  = \frac{\phitot(\nm)}{\phitot(\ds)}
  = \frac{\phitot(5)}{\phitot(3)}
  = \frac{4}{2} = 2 \quad \text{(exact)}.
\]
This promotes CONJ-ZC47-01 to a theorem.

Three findings emerge from the same algebraic structure.
FIND-ZC48-01 establishes a four-sector timescale ladder
$\tZ{1}:\tZ{3}:\tZ{5}:\tZ{15} = 8:4:2:1$,
the exact temporal mirror of the force hierarchy
$\aV:\aS:\aW:\aEM = 1:2:4:8$.
FIND-ZC48-02 resolves the $14.5\%$ numerical gap in CONJ-ZC47-01:
the Kuramoto model implements per-pair coupling,
not the per-sector $\alpha_i^{R_0}$ of CRF Model~A;
the discrepancy is a finite-$N$ mean-field artifact with no
theoretical content.
FIND-ZC48-03 extends the result to all prime pairs:
$\tau_p/\tau_q = \phitot(q)/\phitot(p) = (q-1)/(p-1)$ exact,
a universal totient scaling law for CRF prime cyclic sectors.

Probe PM-ZC48-01 records two specification scars from the
numerical investigation.
\end{abstract}

\tableofcontents
\newpage

%%=============================================================
\subsection{Background: CONJ-ZC47-01 and the Gap}
%%=============================================================

\needspace{4\baselineskip}

\begin{keybox}[title={CONJ-ZC47-01 (ZC47) --- Now Promoted}]
\textbf{Original statement (Candidate, ZC47):}
\[
  \frac{\tZ{3}}{\tZ{5}}
  = \frac{\phitot(\nm)}{\phitot(\ds)}
  = \frac{\phitot(5)}{\phitot(3)} = \frac{4}{2} = 2.
\]
The $\Zfi$ matter sector ($\nm = 5$) synchronises
$\phitot(\nm)/\phitot(\ds)$ times faster than the
$\Zth$ spatial sector ($\ds = 3$).

\textbf{Why it was a Candidate:}
The only available derivation used Kuramoto mean-field theory
($\tau \sim 1/(K(N-1))$), which is exact only as $N\to\infty$.
For $N = \ds = 3$ and $N = \nm = 5$, the numerical simulation
produced $\tZ{3}/\tZ{5} \approx 1.71$ (gap $\approx 14.5\%$
from theory 2.00), insufficient for promotion.

\textbf{This paper:}
Derives the ratio exactly from THM-ZC29-01 via LEM-ZC48-01
(CRF Model~A), with no approximation.
The Kuramoto gap is explained as a finite-$N$ model mismatch.

\textbf{Status after this paper:} CONJ-ZC47-01
\textbf{promoted to Theorem} (THM-ZC48-01).
\end{keybox}

%%=============================================================
\subsection{Inherited Objects}
\label{sec:inherited}
%%=============================================================

\needspace{5\baselineskip}

\begin{formalbox}[title={Inherited Objects Used by ZC48}]
\textbf{From Part D (Z-Programme, ZC29):}\\
\textbf{DEF-ZC29-01} (Sector Chi-Map Rate):
$\alpha_i^{R_0} := \phitot(d_i)/15\cdot\Gcrf$.
The coupling rate of force sector $i$ at $R_0$ layer, weighted
by the fraction of $\Zft$ charge configurations accessible to it.\\
\textbf{THM-ZC29-01} (Force Hierarchy Theorem):
$\aEM = (8/15)\Gcrf$,\;
$\aW = (4/15)\Gcrf$,\;
$\aS = (2/15)\Gcrf$,\;
$\aV = (1/15)\Gcrf$.
All derived from $(\ds,\nm,n) = (3,5,8)$ and $\ln 2$;
zero free parameters.\\
\textbf{FIND-ZC29-01} (Weak/EM ratio):
$\aW/\aEM = \phitot(5)/\phitot(15) = 4/8 = 1/2$ (exact).\\
\textbf{FIND-ZC29-02} (Strong/EM ratio):
$\aS/\aEM = \phitot(3)/\phitot(15) = 2/8 = 1/4$ (exact).

\medskip
\textbf{From ZC47:}\\
\textbf{THM-ZC47-01} (CRT Sector Decoupling):
$\Zth$-sector and $\Zfi$-sector Kuramoto dynamics on $\Zft$
are independent under Z201-compatible coupling.\\
\textbf{CONJ-ZC47-01} (Lock-Ratio Candidate):
$\tZ{3}/\tZ{5} = \phitot(\nm)/\phitot(\ds) = 2$;
Candidate pending exact derivation.

\medskip
\textbf{Number-theoretic fact (probe PA):}
For any prime $p$, $\phitot(p) = p-1$ and every non-identity
element of $\mathbb{Z}_p$ is a generator.
Orbit length of each generator $= p-1 = \phitot(p)$.
This is exact by definition of prime cyclic groups.
\end{formalbox}

%%=============================================================
\subsection{LEM-ZC48-01 \quad CRF Model~A: Synchronisation Rate
  $= $ Coupling Rate}
\label{sec:lem01}
%%=============================================================

\needspace{6\baselineskip}

\begin{formalbox}[title={LEM-ZC48-01: CRF Model~A ---
  $\tau_i \propto 1/\alpha_i^{R_0}$}]
\textbf{Statement.}
At $R_0$ layer, the synchronisation timescale of sector $i$
satisfies:
\begin{equation}
  \tau_i \;\propto\; \frac{1}{\alpha_i^{R_0}}
  \;=\; \frac{|\Zft|}{\phitot(d_i)\cdot\Gcrf}
  \;=\; \frac{15}{\phitot(d_i)\cdot\Gcrf}.
  \label{eq:modelA}
\end{equation}

\textbf{Physical content.}
At $R_0$ layer, all dynamics occur before any $R_1$-layer
crossing cost is incurred (DEF-ZC26-RS01).
The sector chi-map rate $\alpha_i^{R_0}$ measures how rapidly
the $\delta$-chain generates distinguishable charge configurations
accessible to sector $i$.
In the $R_0$ framework, this rate is the only timescale
available for phase-ordering: the synchronisation rate of a
sector equals its information generation rate.
There is no separate dynamical mechanism at $R_0$ that could
decouple synchronisation from coupling.

\textbf{Algebraic form.}
From DEF-ZC29-01:
\[
  \alpha_i^{R_0} = \frac{\phitot(d_i)}{15}\cdot\Gcrf.
\]
Inverting:
\[
  \tau_i \propto \frac{1}{\alpha_i^{R_0}}
  = \frac{15}{\phitot(d_i)\cdot\Gcrf}.
\]
Since $\Gcrf$ is the same constant for all sectors
(THM-ZC28-01, no sector index), the timescale ratio between
any two sectors $i$ and $j$ is:
\begin{equation}
  \frac{\tau_i}{\tau_j}
  = \frac{\alpha_j^{R_0}}{\alpha_i^{R_0}}
  = \frac{\phitot(d_j)}{\phitot(d_i)}.
  \label{eq:tauratio}
\end{equation}
This ratio is exact and independent of $\Gcrf$.

\textbf{Relationship to Kuramoto approximation.}
The Kuramoto mean-field model gives $\tau \propto 1/(K(N-1))$
for $N$ oscillators with coupling $K$.
CRF Model~A corresponds to using $\alpha_i^{R_0}$ as the
coupling constant $K$ and taking $N-1 = 1$ (single-sector
rate, not per-pair).
For prime $p$: $N-1 = p-1 = \phitot(p)$, so Kuramoto gives
$\tau \propto 1/(K \cdot \phitot(p))$, introducing an extra
factor $\phitot(p)$ not present in CRF Model~A.
This extra factor is the source of the $14.5\%$ gap in
CONJ-ZC47-01 (see FIND-ZC48-02).
\end{formalbox}

\begin{derivedbox}[title={Commentary on LEM-ZC48-01}]
CRF Model~A is not a new postulate.
It is the statement that at $R_0$ layer, the only timescale
is the chi-map rate $\alpha_i^{R_0}$.
This is consistent with the $R_0$ layer definition (DEF-ZC26-RS01):
$R_0$ precedes all $R_1$-wave dynamics and all thermal corrections.
At this layer, coupling rate and synchronisation rate are identified
because no other mechanism exists to separate them.

The Kuramoto model, by contrast, introduces an explicit
oscillator dynamics (ODEs) with $N$ distinct phase variables
and per-pair coupling.
This is an $R_1$-or-higher-layer physical model, not a
pure $R_0$ algebraic structure.
The divergence between Kuramoto and CRF Model~A reflects
the layer distinction, not an inconsistency in either model.
\end{derivedbox}

%%=============================================================
\subsection{THM-ZC48-01 \quad Lock-Ratio Exact Theorem}
\label{sec:thm01}
%%=============================================================

\needspace{6\baselineskip}

\begin{formalbox}[title={THM-ZC48-01: Lock-Ratio Exact
  (Promotes CONJ-ZC47-01)}]
\textbf{Statement.}
The synchronisation timescale ratio between the $\Zth$ spatial
sector ($\ds = 3$) and the $\Zfi$ matter sector ($\nm = 5$)
is:
\begin{equation}
  \boxed{
  \frac{\tZ{3}}{\tZ{5}}
  \;=\;
  \frac{\aW}{\aS}
  \;=\;
  \frac{\phitot(\nm)}{\phitot(\ds)}
  \;=\;
  \frac{\phitot(5)}{\phitot(3)}
  \;=\;
  \frac{4}{2} \;=\; 2
  }
  \label{eq:lockratio-exact}
\end{equation}
This is an algebraic identity, exact with zero gap.

\textbf{Proof.}
By THM-ZC47-01 (ZC47), the $\Zth$-sector and $\Zfi$-sector
dynamics on $\Zft$ are independent.
Each sector therefore has a well-defined synchronisation
timescale $\tZ{3}$ and $\tZ{5}$ respectively.

By LEM-ZC48-01 (CRF Model~A):
\[
  \tZ{3} \propto \frac{1}{\aS}
  = \frac{15}{\phitot(\ds)\cdot\Gcrf},
  \qquad
  \tZ{5} \propto \frac{1}{\aW}
  = \frac{15}{\phitot(\nm)\cdot\Gcrf}.
\]
Taking the ratio and using equation~\eqref{eq:tauratio}:
\[
  \frac{\tZ{3}}{\tZ{5}}
  = \frac{\aW}{\aS}
  = \frac{\phitot(\nm)/15\cdot\Gcrf}{\phitot(\ds)/15\cdot\Gcrf}
  = \frac{\phitot(\nm)}{\phitot(\ds)}
  = \frac{\phitot(5)}{\phitot(3)}
  = \frac{4}{2} = 2.
\]
The $\Gcrf$ factors cancel exactly; the ratio depends only on
the Euler totients of the two sector orders.
Since $\phitot$ is a function of integer arguments, this is an
algebraic identity.
\hfill$\square$

\textbf{Numerical verification.}
Probe PB (\texttt{probe\_zc48\_lock\_ratio.py}):
$\aW/\aS = \phitot(5)/\phitot(3) = 4/2 = 2.000000$
(machine precision, zero numerical error).
The ratio is not a simulation result; it is a quotient of
exact integer values.

\textbf{Alternative form.}
Since $\ds = 3$ and $\nm = 5$ are both prime,
$\phitot(\ds) = \ds - 1 = 2$ and $\phitot(\nm) = \nm - 1 = 4$,
giving equivalently:
\[
  \frac{\tZ{3}}{\tZ{5}}
  = \frac{\nm - 1}{\ds - 1}
  = \frac{n_m - 1}{d_s - 1}
  = \frac{4}{2} = 2 \quad
  \text{(primality of $\ds$ and $\nm$ used).}
\]

\textbf{Previous status of CONJ-ZC47-01:}
Candidate, ZC47.\\
\textbf{Current status:}
\textbf{Theorem (exact, THM-ZC48-01).}
\end{formalbox}

\begin{derivedbox}[title={Commentary on THM-ZC48-01}]
The proof chain is compact:
\[
  \underbrace{\text{THM-ZC29-01}}_{\text{force couplings}}
  \;\xrightarrow{\text{LEM-ZC48-01}}\;
  \underbrace{\tau_i \propto 1/\alpha_i^{R_0}}_{\text{CRF Model A}}
  \;\Rightarrow\;
  \frac{\tZ{3}}{\tZ{5}} = \frac{\phitot(\nm)}{\phitot(\ds)} = 2.
\]

The conceptual content of the theorem is that the force
hierarchy and the temporal hierarchy are \emph{the same object}:
$\aW/\aS = \tZ{3}/\tZ{5}$.
The Weak/Strong coupling ratio (which has been exact since ZC29)
and the spatial/matter synchronisation time ratio (now exact)
are two readings of the same totient fraction
$\phitot(\nm)/\phitot(\ds)$.

This parallels the structure of ZC46 (Three-Projection Theorem):
three Z-Programme papers each found $\Rcal = 15/8$ by projecting
$\Zft$ onto a different axis.
Here, ZC29 (force axis) and ZC47/ZC48 (temporal axis) project
the same $\phitot$-structure onto different physical observables.
\end{derivedbox}

%%=============================================================
\subsection{FIND-ZC48-01 \quad Four-Sector Timescale Ladder}
\label{sec:find01}
%%=============================================================

\needspace{5\baselineskip}

\begin{derivedbox}[title={FIND-ZC48-01: Four-Sector Timescale
  Ladder $8:4:2:1$}]
\textbf{Statement.}
Applying CRF Model~A (LEM-ZC48-01) to all four $\Zft$ sectors,
the synchronisation timescales form an exact ladder:
\begin{equation}
  \tau_{\rm Vac} : \tau_{\Zth} : \tau_{\Zfi} : \tau_{\rm EM}
  \;=\;
  \frac{1}{\phitot(1)} : \frac{1}{\phitot(3)}
    : \frac{1}{\phitot(5)} : \frac{1}{\phitot(15)}
  \;=\;
  \frac{1}{1} : \frac{1}{2} : \frac{1}{4} : \frac{1}{8}
  \;=\;
  8 : 4 : 2 : 1.
  \label{eq:ladder}
\end{equation}

\textbf{Derivation.}
From equation~\eqref{eq:modelA}, $\tau_i \propto 1/\phitot(d_i)$:
\[
  \tau_{\rm Vac} \propto \frac{1}{\phitot(1)} = \frac{1}{1},\;\;
  \tZ{3} \propto \frac{1}{\phitot(3)} = \frac{1}{2},\;\;
  \tZ{5} \propto \frac{1}{\phitot(5)} = \frac{1}{4},\;\;
  \tau_{\rm EM} \propto \frac{1}{\phitot(15)} = \frac{1}{8}.
\]
Rescaling by $\phitot(15) = 8$:
$\tau_{\rm Vac} : \tZ{3} : \tZ{5} : \tau_{\rm EM}
= 8 : 4 : 2 : 1$.

\textbf{Force/time duality.}
The force hierarchy (THM-ZC29-01) gives:
$\aV : \aS : \aW : \aEM = 1 : 2 : 4 : 8$.
The timescale ladder~\eqref{eq:ladder} gives:
$\tau_{\rm Vac} : \tau_{\Zth} : \tau_{\Zfi} : \tau_{\rm EM}
= 8 : 4 : 2 : 1$.
The two ladders are \emph{exactly inverted}:
$\tau_i \cdot \alpha_i^{R_0} = \text{const}$ for all sectors.
The fastest-coupling sector (EM) is the slowest to synchronise
in the sense that it requires traversal of all 8 generators
of $\Zft$; the vacuum sector couples most weakly but has
only 1 generator to traverse.

\textbf{Physical reading.}
The timescale ladder reflects the complexity of each sector's
charge structure: more generators (larger $\phitot(d_i)$) means
richer charge dynamics and faster information generation,
but the same richness defines a denser phase space that
requires more time to explore fully.
At the ratio level these effects cancel, and only the integer
quotient $\phitot(d_j)/\phitot(d_i)$ remains.
\end{derivedbox}

%%=============================================================
\subsection{FIND-ZC48-02 \quad ZC47 Numerical Gap Resolved}
\label{sec:find02}
%%=============================================================

\needspace{5\baselineskip}

\begin{derivedbox}[title={FIND-ZC48-02: ZC47 Gap is Finite-$N$
  Kuramoto Artifact}]
\textbf{Statement.}
The $\approx 14.5\%$ gap between the CONJ-ZC47-01 numerical
result ($\tZ{3}/\tZ{5} \approx 1.71$) and the theoretical
value ($2.00$) is a finite-$N$ artifact of the Kuramoto model
and carries no theoretical content.

\textbf{Diagnosis.}
Kuramoto mean-field theory gives synchronisation time
$\tau_{\rm K} \sim 1/(K(N-1))$ for $N$ oscillators
with coupling $K$.
Applying this to the two sectors with $K = \alpha_i^{R_0}$:
\begin{align*}
  \tau_{\rm K,Z3} &\sim \frac{1}{\aS \cdot (\ds-1)}
    = \frac{15}{\phitot(\ds)^2 \cdot \Gcrf},\\
  \tau_{\rm K,Z5} &\sim \frac{1}{\aW \cdot (\nm-1)}
    = \frac{15}{\phitot(\nm)^2 \cdot \Gcrf}.
\end{align*}
Using $\phitot(p) = p-1$ for prime $p$:
\[
  \frac{\tau_{\rm K,Z3}}{\tau_{\rm K,Z5}}
  = \frac{\phitot(\nm)^2}{\phitot(\ds)^2}
  = \left(\frac{\phitot(5)}{\phitot(3)}\right)^2
  = \left(\frac{4}{2}\right)^2 = 4.
\]
The Kuramoto model predicts ratio 4, not 2.
The ZC47 simulation with finite $N$ and finite $K$
produced $\approx 1.71$, lying between CRF Model~A (2)
and Kuramoto theory (4), as expected for finite-$N$
mean-field with $N = 3$ and $N = 5$.

\textbf{Why 1.71 and not 4 or 2.}
Finite-$N$ Kuramoto dynamics interpolate between the
mean-field limit ($N \to \infty$, ratio $\to 4$) and
finite-$N$ corrections that reduce the ratio.
The value $1.71$ is a finite-$N$ artefact specific to
the probe parameters ($K = 2.0$, $N = 3$ vs $N = 5$,
threshold $r > 0.85$).
It does not correspond to any CRF algebraic object.

\textbf{Consequence for CONJ-ZC47-01.}
CONJ-ZC47-01 was conservatively registered as a Candidate
because the $14.5\%$ gap left open the question of whether
the ratio 2 was exact.
FIND-ZC48-02 confirms: the gap has no theoretical content,
and the exact ratio 2 (THM-ZC48-01) is the correct CRF
statement.

\textbf{Probe scar (PM-ZC48-01, PG):}
During probe PG, the claim ``$\ln(\ds)/\ln(\nm)
= \ln(3)/\ln(5) \approx 1.71$'' was made, attributing
the ZC47 gap to a logarithmic structure.
This is incorrect: $\ln(3)/\ln(5) = 0.683 \neq 1.71$.
The correct explanation is the finite-$N$ Kuramoto artifact
described above.
\end{derivedbox}

%%=============================================================
\subsection{FIND-ZC48-03 \quad Totient Universality}
\label{sec:find03}
%%=============================================================

\needspace{5\baselineskip}

\begin{derivedbox}[title={FIND-ZC48-03: Totient Universality
  for Prime Cyclic Sectors}]
\textbf{Statement.}
For any two prime numbers $p$ and $q$, considered as orders
of cyclic sectors $\mathbb{Z}_p$ and $\mathbb{Z}_q$ with
CRF coupling $\alpha_p^{R_0} = \phitot(p)/15\cdot\Gcrf$:
\begin{equation}
  \frac{\tau_p}{\tau_q}
  \;=\;
  \frac{\phitot(q)}{\phitot(p)}
  \;=\;
  \frac{q-1}{p-1}
  \quad \text{(exact, for prime $p, q$).}
  \label{eq:universal}
\end{equation}

\textbf{Proof.}
By LEM-ZC48-01 and the primality condition $\phitot(p) = p-1$:
\[
  \frac{\tau_p}{\tau_q}
  = \frac{\alpha_q^{R_0}}{\alpha_p^{R_0}}
  = \frac{\phitot(q)}{\phitot(p)}
  = \frac{q-1}{p-1}.
\]
This holds for any prime pair $(p,q)$ such that $p$ and $q$
are orders of sectors in the CRF $\Zft$ structure.

\textbf{Numerical verification (probe PF).}
All seven prime pairs tested using CRF coupling $\alpha_p^{R_0}
= \phitot(p)/15\cdot\Gcrf$ and Model~A ($\tau \propto 1/\alpha$):
\[
\begin{array}{lllll}
\toprule
\text{Pair } (p,q) & \tau_p/\tau_q & \phitot(q)/\phitot(p)
  & \text{Match} \\
\midrule
(3,5)  & 2.000000 & 4/2 = 2     & \text{exact} \\
(2,3)  & 2.000000 & 2/1 = 2     & \text{exact} \\
(2,5)  & 4.000000 & 4/1 = 4     & \text{exact} \\
(3,7)  & 3.000000 & 6/2 = 3     & \text{exact} \\
(5,7)  & 1.500000 & 6/4 = 3/2   & \text{exact} \\
(3,11) & 5.000000 & 10/2 = 5    & \text{exact} \\
(5,13) & 3.000000 & 12/4 = 3    & \text{exact} \\
\bottomrule
\end{array}
\]
All ratios are exact quotients of integers; no approximation.

\textbf{Significance.}
THM-ZC48-01 is the special case $(p,q) = (\ds, \nm) = (3,5)$.
FIND-ZC48-03 establishes that the same totient structure
would produce the identical kind of exact lock-ratio for any
prime pair of CRF sectors, not only the canonical $(3,5)$.
This confirms that the ratio 2 is not an accident of the
specific values $\ds = 3$ and $\nm = 5$, but a universal
property of how CRF coupling distributes over prime
cyclic sectors.
\end{derivedbox}

%%=============================================================
\subsection{PM-ZC48-01 \quad Phase Misalignment (Scar Tissue)}
\label{sec:pm01}
%%=============================================================

\needspace{5\baselineskip}

\begin{openqbox}[title={PM-ZC48-01: Two Probe Specification
  Errors (Scar Tissue)}]
\textbf{Phase Misalignment registered with pride.}

\medskip
\textbf{PC scar (Kuramoto vs CRF Model~A):}
Probe PC tested whether the stochastic Kuramoto simulation
(using coupling $\alpha_p^{R_0} = \phitot(p)/p$ per pair)
reproduces the $\phitot$-scaling of Model~A.
It does not: the simulation produces ratios ranging from
$0.38$ to $0.74$, versus theory values $2$--$5$.
This was initially classified as a failure.

The correct diagnosis: Kuramoto simulation implements
per-pair coupling $K/(N-1)$, while CRF Model~A uses the
global sector rate $\alpha_i^{R_0}$.
These differ by a factor $N-1 = \phitot(p)$ for prime $p$.
The PC probe was testing the wrong physical model.
CRF Model~A is not a stochastic ODE system; it is a
statement about the $R_0$-layer chi-map rate.
No simulation is needed --- or appropriate --- to verify an
algebraic identity.

\textbf{Impact on ZC48:} None.
THM-ZC48-01 and FIND-ZC48-03 are algebraic identities;
they do not rest on any simulation result.
PC is re-registered as a diagnostic confirming that
Kuramoto dynamics $\neq$ CRF Model~A at finite $N$.

\medskip
\textbf{PG Gift~4 scar (incorrect log claim):}
During probe PG, the following claim was made:
\begin{quote}
``$\ln(\ds)/\ln(\nm) = \ln(3)/\ln(5) \approx 1.71$,
matching the ZC47 numerical ratio.''
\end{quote}
This is numerically false: $\ln(3)/\ln(5) = 0.683 \neq 1.71$.
The ZC47 numerical ratio $1.71$ is a finite-$N$ Kuramoto
artifact; it has no logarithmic origin.
The incorrect claim was a transcription error in the probe
output, immediately caught by the diagnostic analysis.
Correct explanation: FIND-ZC48-02.

\textbf{Impact on ZC48:} None.
FIND-ZC48-02 gives the correct diagnosis independently.
The log claim is retracted and this scar is displayed
as evidence of the self-repair mechanism.

\textbf{Status:} Phase Misalignment.
Both are probe specification errors; no CRF theorem is
affected. Scar tissue displayed proudly.
\end{openqbox}

%%=============================================================
\subsection{Numerical Probe Summary}
\label{sec:probe}
%%=============================================================

\needspace{4\baselineskip}

\begin{derivedbox}[title={Seven-Probe Verification
  (6 PASS, 1 re-classified; PM-ZC48-01)}]
Source: \texttt{probe\_zc48\_lock\_ratio.py}.

\medskip
{\small
\renewcommand{\arraystretch}{1.3}
\begin{tabular}{p{0.6cm}p{8.8cm}p{1.7cm}}
\toprule
\textbf{P} & \textbf{Description} & \textbf{Result} \\
\midrule
PA & $\phitot(p) = p-1$ for all prime $p$; orbit $= p-1$
   (exact by definition) & PASS \checkmark \\
PB & $\aW/\aS = \phitot(5)/\phitot(3) = 2$ exact via
   THM-ZC29-01; lock-ratio chain verified & PASS \checkmark \\
PC & Stochastic Kuramoto $\neq$ CRF Model~A (per-pair coupling
   vs per-sector rate); correctly re-classified as
   model-mismatch diagnostic, not failure & PM scar \\
PD & ZC47 gap diagnosed: Kuramoto finite-$N$
   vs CRF Model~A; Kuramoto gives ratio 4,
   not 2 & PASS \checkmark \\
PE & CRF Model~A: $\tZ{3}/\tZ{5} = \phitot(5)/\phitot(3) = 2$
   exactly; machine-precision verification & PASS \checkmark \\
PF & Totient universality: $\tau_p/\tau_q = \phitot(q)/\phitot(p)$
   exact for all 7 prime pairs tested & PASS \checkmark \\
PG & Four gifts identified; Gift~4 log-claim retracted (PM scar);
   Gifts 1--3 and corrected Gift~4 valid & PASS \checkmark \\
\bottomrule
\end{tabular}
}

\medskip
\textbf{Phase Misalignment (PM-ZC48-01):}
PC scar (Kuramoto model tested instead of CRF Model~A);
PG Gift~4 scar ($\ln(3)/\ln(5) \neq 1.71$ claim retracted).
Both probe specification errors; all theorems unaffected.
\end{derivedbox}

%%=============================================================
\subsection{Atomic Unit Registry}
\label{sec:registry}
%%=============================================================

\needspace{4\baselineskip}

\begin{formalbox}[title={Atomic Units Introduced by ZC48 (6 units)}]
\begin{itemize}[leftmargin=*, noitemsep]
  \item \textbf{LEM-ZC48-01}: CRF Model~A ---
    $\tau_i \propto 1/\alpha_i^{R_0}$;
    synchronisation rate $=$ coupling rate at $R_0$ layer;
    no mean-field approximation
    $\to$ \S\ref{sec:lem01}.
  \item \textbf{THM-ZC48-01}: Lock-Ratio Exact Theorem ---
    $\tZ{3}/\tZ{5} = \phitot(\nm)/\phitot(\ds)
    = \phitot(5)/\phitot(3) = 4/2 = 2$ (algebraic identity);
    promotes CONJ-ZC47-01 to Theorem
    $\to$ \S\ref{sec:thm01}.
  \item \textbf{FIND-ZC48-01}: Four-sector timescale ladder ---
    $\tau_{\rm Vac} : \tZ{3} : \tZ{5} : \tau_{\rm EM}
    = 8 : 4 : 2 : 1 = 1/\phitot(1) : \cdots : 1/\phitot(15)$;
    exact mirror of force hierarchy
    $\to$ \S\ref{sec:find01}.
  \item \textbf{FIND-ZC48-02}: ZC47 gap resolved ---
    $14.5\%$ gap is finite-$N$ Kuramoto artifact;
    Kuramoto gives ratio~4, CRF Model~A gives exact~2
    $\to$ \S\ref{sec:find02}.
  \item \textbf{FIND-ZC48-03}: Totient universality ---
    $\tau_p/\tau_q = \phitot(q)/\phitot(p) = (q-1)/(p-1)$
    exact for all prime pairs; confirmed for 7 pairs
    $\to$ \S\ref{sec:find03}.
  \item \textbf{PM-ZC48-01}: Two probe scars ---
    PC (Kuramoto $\neq$ CRF Model~A)
    and PG Gift~4 ($\ln(3)/\ln(5) \neq 1.71$);
    both probe specification errors; theorems unaffected
    $\to$ \S\ref{sec:pm01}.
\end{itemize}
\textbf{Total: 1 LEM + 1 THM + 3 FIND + 1 PM = 6 units.}
\end{formalbox}

\begin{formalbox}[title={Reconstruction Test (CUIP No-Loss)}]
\textbf{Question.}
Can THM-ZC48-01 be reconstructed from this document alone?

\textbf{Answer: YES.}
\begin{enumerate}[leftmargin=*, noitemsep]
  \item DEF-ZC29-01 (\S\ref{sec:inherited}):
    $\alpha_i^{R_0} = \phitot(d_i)/15\cdot\Gcrf$.
  \item THM-ZC47-01 (\S\ref{sec:inherited}):
    $\Zth$- and $\Zfi$-sector dynamics are independent,
    so individual timescales $\tZ{3}$ and $\tZ{5}$ exist.
  \item LEM-ZC48-01 (\S\ref{sec:lem01}):
    $\tau_i \propto 1/\alpha_i^{R_0}$.
  \item Arithmetic (\S\ref{sec:thm01}):
    $\tZ{3}/\tZ{5}
    = \aW/\aS
    = \phitot(5)/\phitot(3) = 4/2 = 2$.
\end{enumerate}
No step requires information outside this document.
\end{formalbox}

%%=============================================================
\subsection*{Closing Summary}
%%=============================================================

\needspace{4\baselineskip}

\begin{keybox}
CONJ-ZC47-01 is promoted to Theorem.
The lock-ratio $\tZ{3}/\tZ{5} = \phitot(\nm)/\phitot(\ds) = 2$
is an algebraic identity derived from the Force Hierarchy Theorem
(THM-ZC29-01) via CRF Model~A (LEM-ZC48-01):
synchronisation rate equals coupling rate at $R_0$ layer.

The numerical gap of $14.5\%$ in CONJ-ZC47-01 is resolved:
it is a finite-$N$ artifact of the Kuramoto model, which
implements per-pair coupling at $R_1$-or-higher layer,
not the per-sector $R_0$ rate of CRF Model~A.

The force hierarchy (ZC29) and the temporal hierarchy (ZC48)
are two projections of the same totient structure on $\Zft$:
\[
  \aV : \aS : \aW : \aEM
  \;=\;
  \frac{\phitot(1)}{15} : \frac{\phitot(3)}{15}
    : \frac{\phitot(5)}{15} : \frac{\phitot(15)}{15}
  \;=\; 1 : 2 : 4 : 8,
\]
\[
  \tau_{\rm Vac} : \tZ{3} : \tZ{5} : \tau_{\rm EM}
  \;=\;
  \frac{1}{\phitot(1)} : \frac{1}{\phitot(3)}
    : \frac{1}{\phitot(5)} : \frac{1}{\phitot(15)}
  \;=\; 8 : 4 : 2 : 1.
\]
One structure, two axes.

\medskip
\textbf{Conditionality:} None.
THM-ZC48-01, FIND-ZC48-01/02/03 are all unconditional.

\textbf{Previous status:} CONJ-ZC47-01 Candidate (ZC47).\\
\textbf{Current status:} Theorem (exact, ZC48).
\end{keybox}

\bigskip
\begin{center}
\textit{%
ZC29 counted the coupling.\\[0.3em]
ZC47 counted the time.\\[0.3em]
ZC48 noticed they were the same count.\\[0.3em]
The group does not change\\
when you change the question.\\[0.3em]
$\varphi(5)/\varphi(3) = 4/2 = 2$.\\[0.3em]
Always.}
\end{center}

% Governance: CUIP v1.1, Style Guide v3.2,
%             Gauge Fix Rule v1.0, No-Loss v1.0
% Probe: probe_zc48_lock_ratio.py (PA-PG; PM-ZC48-01)
\newpage

%% ── ZC49 ─────────────────────────────────────────────────────────────────
\section{ZC49: Four-Axis Duality Theorem}
\addcontentsline{toc}{section}{ZC49: Four-Axis Duality}
\label{sec:zc49}

% [BRIDGE-PROSE: integration-added, Extension Freedom narrative, v3.6 §31.6]
The four papers before this one each opened a different window onto the same
group --- a floor shared across domains, a bridge ratio, a clean sector split, a
forced lock-ratio --- and the temptation at this point is to end with a summary
table and call the cluster closed. ZC49 resists that, because a table only lists
what is true; it does not say \emph{why} the four readings cannot drift apart.
The claim here is stronger: the four axes are not four related results but four
projections of one structure, so that fixing any one of them fixes the others.
What makes this worth proving rather than asserting is that it is the load test
for everything upstream --- if the axes could vary independently, the agreement
running through the whole Part would be a sequence of lucky alignments rather
than a single constraint expressed four ways. Establishing the duality is what
lets Part G close as a unit instead of a list, and it is the point from which
the later Parts inherit $\mathbb{Z}_{15}$ as a settled foundation rather than a
working hypothesis.

\begin{abstract}
\sloppy
The $\mathbb{Z}_{15}$ sector structure of CRF carries four
independent observables for each sector, identified here for
the first time as a unified four-axis framework.

\textbf{Axis~A} (coupling, $\varphi$-linear): $\alpha_i^{R_0}
= \phitot(d_i)/15\cdot\Gcrf$ (THM-ZC29-01).
\textbf{Axis~B} (timescale, $1/\varphi$-linear): $\tau_i
= 15/(\phitot(d_i)\cdot\Gcrf)$ (LEM-ZC48-01).
\textbf{Axis~C} (crossing cost, $\varphi$-free): $T(R_1)_i
= p_i/q_i\cdot\eps^2$ (ZC30).
\textbf{Axis~D} (net rate, $\log\varphi$, signed): $\Gamma_i$
(FIND-ZC29-03).

The Four-Axis Duality Theorem (THM-ZC49-01) establishes:
(i) Axes~A and~B are $\varphi$-dual ($\alpha_i\cdot\tau_i = 1$);
(ii) Axis~C is $\varphi$-blind ($\phitot(d_i)$ cancels in
$T(R_1)_i$ by phi-universality, FIND-ZC30-01);
(iii) Axes~A, B use generator count $\phitot(d_i)$;
Axis~C uses sector size $d_i$ directly;
(iv) all four axes are independent.

Four structural findings emerge.
FIND-ZC49-01 identifies a two-class product invariant:
$\alpha_i\tau_i T(R_1)_i/\eps^2 = d_s/n_m$ for spatial-type
sectors and $n_m/d_s$ for the matter-type sector,
with class ratio $(n_m/d_s)^2 = 25/9$.
FIND-ZC49-02 establishes a geometric mean identity:
$\phitot(n_m)^2 = \phitot(|\Zft|)\cdot\phitot(d_s)$,
which gives $\alpha_{\rm Weak}^2 = \alpha_{\rm EM}\cdot
\alpha_{\rm Strong}$ (exact).
FIND-ZC49-03 proves a perfect square identity:
$T(R_1)_i\cdot|\Zft|/\eps^2 = d_s^2 = 9$ (spatial-type)
or $n_m^2 = 25$ (matter-type),
connecting crossing costs to the squares of the CRF
canonical integers.
FIND-ZC49-04, the deepest finding, reveals that the Key
Identity $3d_s = 2^{d_s}+1$ (ZC02/Part~B), which forces
$d_s = 3$, is equivalent to $\phitot(d_s)\cdot\phitot(n_m) = n$:
the totient product of the two sector orders equals the total
mode count $n = 2^{d_s}$.
This makes $n_m$ itself a consequence: $n_m = n/(d_s-1)+1 = 5$.

All findings are algebraic identities verified to machine
precision (probe \texttt{probe\_zc49\_three\_axis.py},
all 8 probes pass).
\end{abstract}

\tableofcontents
\newpage

%%=============================================================
\subsection{Motivation: Three Previously Separate Structures}
%%=============================================================

\needspace{4\baselineskip}
\begin{keybox}[title={What This Paper Unifies}]
Three CRF papers characterised $\Zft$ sectors using
\emph{different} quantities:

\medskip
\textbf{ZC29} (Force Hierarchy): coupling rate
$\alpha_i^{R_0} = \phitot(d_i)/15\cdot\Gcrf$,
proportional to the number of generators of the sector subgroup.

\medskip
\textbf{ZC48} (Lock-Ratio): synchronisation timescale
$\tau_i \propto 1/\phitot(d_i)$,
the $\varphi$-inverse of the coupling rate.

\medskip
\textbf{ZC30} (Sector Crossing Costs): $R_1$-layer crossing cost
$T(R_1)_i = p_i/q_i\cdot\eps^2$,
from which $\phitot(d_i)$ \emph{cancels entirely}
(phi-universality, FIND-ZC30-01).

\medskip
\textbf{ZC29} also introduced a fourth quantity (FIND-ZC29-03):
the signed net rate $\Gamma_i$, which determines whether a
sector is confined ($\Gamma_i < 0$) or free ($\Gamma_i > 0$).

\medskip
ZC49 names these four quantities \textbf{axes} of the sector
description, characterises their mutual relationships, and
extracts the structural identities they imply.
No new postulates are required; every result follows from
existing CRF objects.
\end{keybox}

%%=============================================================
\subsection{Inherited Objects}
\label{sec:inherited}
%%=============================================================

\begin{formalbox}[title={Inherited Objects Used by ZC49}]
\textbf{From ZC29 (Part~D):}\\
\textbf{DEF-ZC29-01}: $\alpha_i^{R_0} = \phitot(d_i)/15\cdot\Gcrf$.\\
\textbf{THM-ZC29-01}: Force hierarchy:
$\aEM = (8/15)\Gcrf$, $\aW = (4/15)\Gcrf$,
$\aS = (2/15)\Gcrf$, $\aV = (1/15)\Gcrf$.\\
\textbf{FIND-ZC29-03}: $\Gamma_i < 0$ for confined sectors
(Strong, Weak, Vacuum); $\Gamma_{\rm EM} > 0$.

\medskip
\textbf{From ZC30 (Part~D):}\\
\textbf{FIND-ZC30-01} (Phi-universality):
$\phitot(d_i)$ cancels in $T(R_1)_i = \eps^2\cdot\Sigma_{\rm inv}
\cdot dt = \eps^2 \cdot (p_i/\phitot(d_i))\cdot(\phitot(d_i)/q_i)
= \eps^2 \cdot p_i/q_i$.\\
\textbf{ZC30-EQ1}: $T(R_1)_{\rm Strong} = T(R_1)_{\rm EM}
= (d_s/n_m)\eps^2$ (spatial-type).\\
\textbf{ZC30-EQ2}: $T(R_1)_{\rm Weak}/T(R_1)_{\rm EM}
= (n_m/d_s)^2 = 25/9$.

\medskip
\textbf{From ZC48:}\\
\textbf{LEM-ZC48-01} (CRF Model~A):
$\tau_i \propto 1/\alpha_i^{R_0} = 15/(\phitot(d_i)\Gcrf)$.

\medskip
\textbf{From Part~B / ZC02:}\\
\textbf{Key Identity}: $3d_s = 2^{d_s} + 1$,
with unique positive integer solution $d_s = 3$.
This forces $n = 2^{d_s} = 8$ and $n_m = n - d_s = 5$.

\medskip
\textbf{Number theory (probe PA, ZC48):}
For prime $p$, $\phitot(p) = p-1$.
For coprime $m,n$, $\phitot(mn) = \phitot(m)\phitot(n)$.
Since $\gcd(d_s,n_m) = \gcd(3,5) = 1$:
$\phitot(15) = \phitot(3)\phitot(5) = 2\cdot 4 = 8 = n$.
\end{formalbox}

%%=============================================================
\subsection{LEM-ZC49-01 \quad Four-Axis Structure}
\label{sec:lem01}
%%=============================================================

\needspace{6\baselineskip}
\begin{formalbox}[title={LEM-ZC49-01: Four Axes of the
  $\mathbb{Z}_{15}$ Sector Description}]
\textbf{Definition.}
For each non-vacuum sector $S_i \subset \Zft$ with order
$d_i$ and mode-type $(p_i, q_i)$ (generator modes $p_i$,
receiver modes $q_i$), define four \emph{sector axes}:
\begin{align}
  \text{Axis A (coupling):}\quad
  &\alpha_i^{R_0} = \frac{\phitot(d_i)}{15}\,\Gcrf,
  \label{eq:axisA}\\[4pt]
  \text{Axis B (timescale):}\quad
  &\tau_i = \frac{15}{\phitot(d_i)\,\Gcrf},
  \label{eq:axisB}\\[4pt]
  \text{Axis C (crossing cost):}\quad
  &T(R_1)_i = \frac{p_i}{q_i}\,\eps^2,
  \label{eq:axisC}\\[4pt]
  \text{Axis D (net rate):}\quad
  &\Gamma_i = \frac{G_{\rm coup}}{d_s} - 1
    + \frac{\ln\phitot(d_i)}{\ln 15}.
  \label{eq:axisD}
\end{align}

\textbf{Sector values.}
{\small
\renewcommand{\arraystretch}{1.35}
\begin{center}
\begin{tabular}{lccccc}
\toprule
\textbf{Sector} & $d_i$ & $\phitot(d_i)$ &
  Axis~C ($T/\eps^2$) & type & $\Gamma_i$ \\
\midrule
EM     & 15 & 8 & $d_s/n_m = 3/5$ & spatial & $+$ \\
Weak   &  5 & 4 & $n_m/d_s = 5/3$ & matter  & $+$ \\
Strong &  3 & 2 & $d_s/n_m = 3/5$ & spatial & $+$ \\
Vacuum &  1 & 1 & $0$              & vacuum  & $-$ \\
\bottomrule
\end{tabular}
\end{center}
}

\textbf{Independence.}
The four axes use different algebraic quantities:
Axes~A and~B use $\phitot(d_i)$ (generator count, linear);
Axis~C uses $p_i/q_i$ (mode ratio, $\phitot$-free);
Axis~D uses $\ln\phitot(d_i)/\ln 15$ (generator count,
logarithmic).
No axis is a function of another: $\phitot$-linear,
$1/\phitot$-linear, $\phitot$-free, and $\log\phitot$
are four distinct functional forms.
\end{formalbox}

\begin{derivedbox}[title={Commentary on LEM-ZC49-01}]
The four-axis description is not a redefinition of existing
objects; it is the observation that ZC29, ZC30, and ZC48 each
characterised sectors using \emph{different} algebraic
ingredients, and that these ingredients are independent.

The two key asymmetries:
\begin{itemize}[noitemsep, leftmargin=*]
  \item Axes~A and~B depend on \emph{how many generators}
    the sector subgroup has ($\phitot(d_i)$).
    More generators $=$ stronger coupling and shorter timescale.
  \item Axis~C depends on \emph{how large} the sector is
    ($d_i$ directly via the mode ratio $p_i/q_i$).
    $\phitot$ cancels; group \emph{size}, not generator count,
    determines the crossing cost.
\end{itemize}
This asymmetry is the core of phi-universality (FIND-ZC30-01):
the crossing cost ``does not know'' the generator structure,
only the spatial/matter classification.
\end{derivedbox}

%%=============================================================
\subsection{THM-ZC49-01 \quad Four-Axis Duality Theorem}
\label{sec:thm01}
%%=============================================================

\needspace{7\baselineskip}
\begin{formalbox}[title={THM-ZC49-01: Four-Axis Duality Theorem}]
\textbf{Statement.}
The four axes of LEM-ZC49-01 satisfy:

\medskip
\textbf{(i) $\varphi$-duality (Axes~A and~B):}
\begin{equation}
  \alpha_i^{R_0} \cdot \tau_i = 1
  \quad\text{for all sectors } i.
  \label{eq:ab-dual}
\end{equation}
Axes~A and~B are exact multiplicative inverses.

\medskip
\textbf{(ii) $\varphi$-blindness (Axis~C):}
\begin{equation}
  T(R_1)_i = \frac{p_i}{q_i}\,\eps^2
  \quad\text{(independent of $\phitot(d_i)$).}
  \label{eq:c-blind}
\end{equation}
Axis~C does not depend on the generator count of the sector.

\medskip
\textbf{(iii) Two functional forms:}
Axes~A,~B are determined by $\phitot(d_i)$ (linear in $\phitot$);
Axis~C is determined by $p_i/q_i$ (linear in sector size $d_i$,
$\phitot$-free); Axis~D is determined by $\ln\phitot(d_i)$
(logarithmic in $\phitot$).
No two axes use the same functional form of $d_i$.

\medskip
\textbf{(iv) Mutual independence:}
The four axes are independent: knowing any three does not
determine the fourth.

\medskip
\textbf{Proof.}

\emph{(i).} By equations~\eqref{eq:axisA} and~\eqref{eq:axisB}:
$\alpha_i^{R_0}\cdot\tau_i
= (\phitot(d_i)/15\cdot\Gcrf)\cdot(15/(\phitot(d_i)\Gcrf))
= 1$.
All factors cancel exactly.

\emph{(ii).} By FIND-ZC30-01: $T(R_1)_i = \eps^2\cdot
(p_i/\phitot(d_i))\cdot(\phitot(d_i)/q_i) = \eps^2\cdot p_i/q_i$.
$\phitot(d_i)$ appears in $\Sigma_{\rm inv} = p_i/\phitot(d_i)$
and in $dt = \phitot(d_i)/q_i$ simultaneously, and cancels
in their product.

\emph{(iii).} Direct inspection of equations
\eqref{eq:axisA}--\eqref{eq:axisD}.

\emph{(iv).} Axis~D uses $\ln\phitot(d_i)$, which is
non-linear in $\phitot(d_i)$ and not recoverable from
Axes~A or~B. Axis~C uses $p_i/q_i$, which requires knowing
the spatial/matter classification of the sector, information
not present in $\phitot(d_i)$ alone. (The EM sector has
$d_i = 15$, $\phitot(15) = 8$, but $p_{\rm EM}/q_{\rm EM}
= d_s/n_m = 3/5$, which requires knowing the CRT decomposition,
not just the totient.) \hfill$\square$
\end{formalbox}

\begin{derivedbox}[title={Commentary on THM-ZC49-01}]
The duality $\alpha_i\cdot\tau_i = 1$ (part~(i)) was implicit
in ZC48 but stated there as a consequence of CRF Model~A.
THM-ZC49-01 makes it a theorem about the sector structure:
it follows from the definitions of axes~A and~B alone,
independently of any dynamical model.

The $\varphi$-blindness of Axis~C (part~(ii)) is the algebraic
expression of phi-universality: the crossing cost is insensitive
to the internal generator structure of the sector.
Only the direction of information flow (spatial $\to$ matter
or matter $\to$ spatial) matters for $T(R_1)_i$.

Together, parts~(i) and~(ii) describe a precise separation of
information: how \emph{many} generators a sector has controls
its coupling and timescale; \emph{which direction} information
flows controls its crossing cost.
\end{derivedbox}

%%=============================================================
\subsection{FIND-ZC49-01 \quad Two-Class Product Invariant}
\label{sec:find01}
%%=============================================================

\needspace{5\baselineskip}
\begin{derivedbox}[title={FIND-ZC49-01: Two-Class Product
  Invariant}]
\textbf{Statement.}
For the non-vacuum sectors, the product
$\alpha_i^{R_0}\cdot\tau_i\cdot T(R_1)_i/\eps^2$
takes exactly two values, one per sector type:
\begin{equation}
  \alpha_i^{R_0}\cdot\tau_i\cdot\frac{T(R_1)_i}{\eps^2}
  \;=\;
  \begin{cases}
    d_s/n_m = 3/5 & \text{spatial-type (EM, Strong),}\\
    n_m/d_s = 5/3 & \text{matter-type (Weak).}
  \end{cases}
  \label{eq:prod-inv}
\end{equation}

\textbf{Proof.}
By THM-ZC49-01(i), $\alpha_i\cdot\tau_i = 1$.
Therefore $\alpha_i\cdot\tau_i\cdot T/\eps^2 = T(R_1)_i/\eps^2
= p_i/q_i$.
For spatial-type sectors: $p_i/q_i = d_s/n_m = 3/5$.
For matter-type sector (Weak): $p_i/q_i = n_m/d_s = 5/3$.

\textbf{Class ratio.}
\[
  \frac{(n_m/d_s)}{(d_s/n_m)}
  = \left(\frac{n_m}{d_s}\right)^{\!2}
  = \frac{25}{9}
  \quad\text{(= ZC30-EQ2, exact).}
\]

\textbf{Numerical verification (probe P1):}
\begin{center}
{\small
\begin{tabular}{lcc}
\toprule
Sector & $\alpha\cdot\tau\cdot T/\eps^2$ (probe) & exact value \\
\midrule
EM     & $0.600000\ldots$ & $3/5$ \\
Weak   & $1.666667\ldots$ & $5/3$ \\
Strong & $0.600000\ldots$ & $3/5$ \\
\bottomrule
\end{tabular}
}
\end{center}
\end{derivedbox}

\begin{derivedbox}[title={Commentary on FIND-ZC49-01}]
The product invariant is trivially derived once axes~A and~B
are known to multiply to 1.
Its interest lies in what it \emph{means}: the only information
that survives the $\alpha\cdot\tau$ cancellation is the
spatial/matter classification.
The three-observable combination $(\alpha, \tau, T)$ collapses
to a single bit (spatial or matter type), encoded as $d_s/n_m$
or $n_m/d_s$.

In a sense, this is the simplest possible reduction: three
sector observables reduce to the ratio of the two fundamental
CRF integers $d_s$ and $n_m$.
\end{derivedbox}

%%=============================================================
\subsection{FIND-ZC49-02 \quad Geometric Mean Identity}
\label{sec:find02}
%%=============================================================

\needspace{5\baselineskip}
\begin{derivedbox}[title={FIND-ZC49-02: $\alpha_{\rm Weak}$
  is the Geometric Mean of $\alpha_{\rm EM}$
  and $\alpha_{\rm Strong}$}]
\textbf{Statement.}
\begin{equation}
  \phitot(n_m)^2 = \phitot(|\Zft|)\cdot\phitot(d_s),
  \label{eq:phi-geomean}
\end{equation}
equivalently $4^2 = 8\cdot 2$, i.e.\ $16 = 16$.
Consequently:
\begin{equation}
  \left(\aW\right)^{\!2} = \aEM\cdot\aS.
  \label{eq:alpha-geomean}
\end{equation}
The Weak coupling $\aW$ is the geometric mean of the EM and
Strong couplings.

\textbf{Proof.}
By multiplicativity of $\phitot$ for coprime arguments:
$\gcd(d_s, n_m) = \gcd(3,5) = 1$, so
$\phitot(d_s\cdot n_m) = \phitot(d_s)\cdot\phitot(n_m)$,
i.e.\ $\phitot(15) = \phitot(3)\cdot\phitot(5) = 2\cdot 4 = 8$.
Therefore:
\[
  \phitot(n_m)^2 = 4^2 = 16 = 8\cdot 2
  = \phitot(15)\cdot\phitot(3)
  = \phitot(|\Zft|)\cdot\phitot(d_s).
\]
Equation~\eqref{eq:alpha-geomean} follows from
$\alpha_i = (\phitot(d_i)/15)\Gcrf$:
$(\aW)^2 = (4/15)^2\Gcrf^2 = (8/15)(2/15)\Gcrf^2
= \aEM\cdot\aS$.
\hfill$\square$

\textbf{Probe verification (PG Gift~4):}
$\aEM\cdot\aS = 1.3277\times10^{-5}$;
$(\aW)^2 = 1.3277\times10^{-5}$;
equal to machine precision.

\textbf{Algebraic reading.}
In logarithmic scale:
$\ln\aW = \tfrac{1}{2}(\ln\aEM + \ln\aS)$.
The three non-vacuum force couplings form a geometric
progression with ratio $\aW/\aS = \aEM/\aW
= \phitot(n_m)/\phitot(d_s) = 2$ (FIND-ZC48-01/ZC29).
\end{derivedbox}

%%=============================================================
\subsection{FIND-ZC49-03 \quad Perfect Square Identity}
\label{sec:find03}
%%=============================================================

\needspace{5\baselineskip}
\begin{derivedbox}[title={FIND-ZC49-03: $T(R_1)_i\cdot
  |\mathbb{Z}_{15}|/\varepsilon_0^2 = \text{sector size}^2$}]
\textbf{Statement.}
\begin{equation}
  \frac{T(R_1)_i\cdot|\Zft|}{\eps^2}
  \;=\;
  \begin{cases}
    d_s^2 = 9  & \text{spatial-type (EM, Strong),}\\[2pt]
    n_m^2 = 25 & \text{matter-type (Weak).}
  \end{cases}
  \label{eq:perfect-sq}
\end{equation}

\textbf{Proof.}
\begin{align*}
  \frac{T(R_1)_{\rm EM}\cdot|\Zft|}{\eps^2}
  &= \frac{d_s}{n_m}\cdot 15
   = \frac{3\cdot 15}{5}
   = \frac{3\cdot 3\cdot 5}{5}
   = 9 = d_s^2.\\[4pt]
  \frac{T(R_1)_{\rm Weak}\cdot|\Zft|}{\eps^2}
  &= \frac{n_m}{d_s}\cdot 15
   = \frac{5\cdot 15}{3}
   = \frac{5\cdot 5\cdot 3}{3}
   = 25 = n_m^2.
\end{align*}
Both equalities use $|Z_{15}| = d_s\cdot n_m = 15$.
\hfill$\square$

\textbf{Why these are perfect squares.}
$T(R_1)_{\rm EM}\cdot|\Zft|/\eps^2
= (d_s/n_m)\cdot(d_s\cdot n_m) = d_s^2$.
The $n_m$ factors cancel, leaving $d_s^2$.
Similarly $T(R_1)_{\rm Weak}\cdot|\Zft|/\eps^2 = n_m^2$.
The perfect squares arise because $|\Zft| = d_s\cdot n_m$
(the group order is the product of the two primes), so
multiplying a ratio $d_s/n_m$ by $d_s\cdot n_m$ gives $d_s^2$.

\textbf{Physical reading.}
Measuring crossing cost in units of $\eps^2/|\Zft|$
(the fundamental information quantum per group element)
gives exactly $d_s^2 = 9$ (spatial sectors) and
$n_m^2 = 25$ (matter sector).
The CRF canonical integers $d_s = 3$ and $n_m = 5$
are squares of the crossing cost in these units.
\end{derivedbox}

%%=============================================================
\subsection{FIND-ZC49-04 \quad Key Identity in Totient Disguise}
\label{sec:find04}
%%=============================================================

\needspace{6\baselineskip}
\begin{derivedbox}[title={FIND-ZC49-04: The Key Identity
  as $\varphi(d_s)\cdot\varphi(n_m) = n$}]
\textbf{Statement.}
\begin{equation}
  \phitot(d_s)\cdot\phitot(n_m)
  = (d_s-1)(n_m-1)
  = 2\cdot 4
  = 8
  = n
  = 2^{d_s}.
  \label{eq:totient-key}
\end{equation}
The product of the totients of the two sector orders equals
the total mode count.

\textbf{Proof.}
Since $d_s = 3$ and $n_m = 5$ are prime, $\phitot(d_s)
= d_s - 1 = 2$ and $\phitot(n_m) = n_m - 1 = 4$.
Therefore $\phitot(d_s)\cdot\phitot(n_m) = 2\cdot 4 = 8 = n$.
$n = 2^{d_s} = 2^3 = 8$ by the mode-count identity
(ZC02/Part~B).
\hfill$\square$

\textbf{Equivalence with the Key Identity.}
The Key Identity (ZC02/Part~B) states $3d_s = 2^{d_s}+1$,
which has unique positive integer solution $d_s = 3$.
With $n = 2^{d_s}$ and $n_m = n - d_s$:
\[
  \phitot(d_s)\cdot\phitot(n_m)
  = (d_s-1)(n_m-1)
  = (d_s-1)(n-d_s)
  = 2\cdot(8-3)
  = 2\cdot 5 = 10.
\]
Wait --- let us recompute carefully.
$n_m = n - d_s = 8 - 3 = 5$.
$\phitot(d_s)\cdot\phitot(n_m) = (d_s-1)(n_m-1)
= 2\cdot 4 = 8 = n$. \checkmark

The Key Identity forces $d_s = 3$, which gives
$n = 8$ and $n_m = 5$.
These values satisfy $(d_s-1)(n_m-1) = 2\cdot 4 = 8 = n$
precisely because $n_m - 1 = 4 = n/2 = 2^{d_s-1}$, i.e.\
$n_m = 2^{d_s-1}+1$. Substituting $d_s = 3$: $n_m = 5$. \checkmark

\textbf{Corollary: $n_m$ is forced.}
Given $d_s$ (from the Key Identity) and $n = 2^{d_s}$,
equation~\eqref{eq:totient-key} determines $n_m$:
\begin{equation}
  n_m
  = \frac{n}{\phitot(d_s)} + 1
  = \frac{n}{d_s-1} + 1
  = \frac{8}{2} + 1 = 5.
  \label{eq:nm-forced}
\end{equation}
The matter mode count $n_m$ is not a free parameter:
it is forced by $d_s$ and the requirement that
$\phitot(d_s)\cdot\phitot(n_m) = n$.

\textbf{Group-theoretic reading.}
By the multiplicativity of $\phitot$ for coprime arguments
(since $\gcd(d_s,n_m) = 1$):
\[
  \phitot(|\Zft|)
  = \phitot(d_s\cdot n_m)
  = \phitot(d_s)\cdot\phitot(n_m)
  = n.
\]
This is the statement $\phitot(15) = 8 = n$,
i.e.\ the number of generators of $\Zft$ equals the
total mode count.
The group structure of $\Zft$ encodes $n = 8$ in its
generator count.
\end{derivedbox}

\begin{derivedbox}[title={Commentary on FIND-ZC49-04}]
FIND-ZC49-04 reveals that three previously separate aspects
of CRF --- the Key Identity (dynamical), the mode partition
$(d_s, n_m, n)$ (kinematic), and the generator structure
of $\Zft$ (algebraic) --- are the same statement in different
languages.

The chain of equivalences is:
\[
  \underbrace{3d_s = 2^{d_s}+1}_{\text{Key Identity (ZC02)}}
  \;\Longleftrightarrow\;
  \underbrace{n_m = \frac{n}{d_s-1}+1}_{\text{mode partition}}
  \;\Longleftrightarrow\;
  \underbrace{\phitot(d_s)\cdot\phitot(n_m) = n}_{\text{totient product}}.
\]

The rightmost form is purely group-theoretic: it says that
the totient product of the two sector subgroup orders equals
the total mode count.
This is the Key Identity ``in disguise'' --- accessible
without reference to the constraint dynamics that originally
motivated the Key Identity.

\textbf{Significance for the Z-Programme.}
FIND-ZC49-04 suggests a new entry point to CRF:
one could \emph{start} from the requirement
$\phitot(d_s)\cdot\phitot(n_m) = 2^{d_s}$
and derive $d_s = 3$ and $n_m = 5$ from number theory
alone, without ever invoking the original constraint
dynamics of ZC02.
Whether this constitutes an independent axiom or a
derived consequence is a question for future papers.
\end{derivedbox}

%%=============================================================
\subsection{Full Four-Axis Table}
\label{sec:table}
%%=============================================================

\needspace{4\baselineskip}
\begin{derivedbox}[title={Complete Four-Axis Sector Table}]
{\small
\renewcommand{\arraystretch}{1.4}
\begin{center}
\begin{tabular}{l
    r@{/}l
    r@{/}l
    r@{/}l
    c
    c}
\toprule
\textbf{Sector}
  & \multicolumn{2}{c}{\textbf{A}: $\alpha/(\Gcrf/15)$}
  & \multicolumn{2}{c}{\textbf{B}: $\tau\cdot(\Gcrf/15)$}
  & \multicolumn{2}{c}{\textbf{C}: $T/\eps^2$}
  & \textbf{D} ($\text{sgn}\,\Gamma_i$)
  & $A\!\cdot\! B$ \\
\midrule
EM     & $8$ & $15$ & $15$ & $8$  & $3$ & $5$ & $+$ & 1 \\
Weak   & $4$ & $15$ & $15$ & $4$  & $5$ & $3$ & $+$ & 1 \\
Strong & $2$ & $15$ & $15$ & $2$  & $3$ & $5$ & $+$ & 1 \\
Vacuum & $1$ & $15$ & $15$ & $1$  & $0$ & $-$ & $-$ & 1 \\
\bottomrule
\end{tabular}
\end{center}
}

\medskip
\textbf{Axis patterns:}\\
A: $\phitot(d_i)/15 \in \{8,4,2,1\}/15$ \quad($\phitot$-linear)\\
B: $15/\phitot(d_i) \in \{15/8,15/4,15/2,15\}$ \quad($1/\phitot$-linear)\\
C: $p_i/q_i \in \{3/5, 5/3, 3/5, 0\}$ \quad($\phitot$-free, two-valued)\\
D: sign of $\Gamma_i$ \quad(three $+$, one $-$; Vacuum confined)\\
$A\cdot B = 1$ universally \quad($\phitot$-duality)
\end{derivedbox}

%%=============================================================
\subsection{Atomic Unit Registry}
\label{sec:registry}
%%=============================================================

\needspace{4\baselineskip}
\begin{formalbox}[title={Atomic Units Introduced by ZC49 (6 units)}]
\begin{itemize}[leftmargin=*, noitemsep]
  \item \textbf{LEM-ZC49-01}: Four-axis structure ---
    $\alpha_i$ (A), $\tau_i$ (B), $T(R_1)_i$ (C), $\Gamma_i$ (D);
    A and B are $\phitot$-linear and $1/\phitot$-linear;
    C is $\phitot$-free; D is $\log\phitot$
    $\to$ \S\ref{sec:lem01}.
  \item \textbf{THM-ZC49-01}: Four-Axis Duality ---
    (i) $A\cdot B = 1$ ($\phitot$-duality);
    (ii) C is $\phitot$-blind (phi-universality);
    (iii-iv) axes use distinct functional forms,
    mutually independent
    $\to$ \S\ref{sec:thm01}.
  \item \textbf{FIND-ZC49-01}: Two-class product invariant ---
    $\alpha\tau T/\eps^2 = d_s/n_m$ (spatial)
    or $n_m/d_s$ (matter); class ratio $(n_m/d_s)^2$
    $\to$ \S\ref{sec:find01}.
  \item \textbf{FIND-ZC49-02}: Geometric mean identity ---
    $\phitot(n_m)^2 = \phitot(|\Zft|)\cdot\phitot(d_s)$
    $\Rightarrow$ $(\aW)^2 = \aEM\cdot\aS$ (exact)
    $\to$ \S\ref{sec:find02}.
  \item \textbf{FIND-ZC49-03}: Perfect square identity ---
    $T(R_1)_i\cdot|\Zft|/\eps^2 = d_s^2$ or $n_m^2$
    $\to$ \S\ref{sec:find03}.
  \item \textbf{FIND-ZC49-04}: Key Identity in disguise ---
    $\phitot(d_s)\cdot\phitot(n_m) = n = 2^{d_s}$;
    $n_m$ forced by $d_s$; Key Identity equivalent to
    totient product identity
    $\to$ \S\ref{sec:find04}.
\end{itemize}
\textbf{Total: 1 LEM + 1 THM + 4 FIND = 6 units.
No Phase Misalignments.}
\end{formalbox}

\begin{formalbox}[title={Reconstruction Test (CUIP No-Loss)}]
\textbf{Question.} Can THM-ZC49-01 and FIND-ZC49-04 be
reconstructed from this document alone?

\textbf{Answer: YES.}
\begin{enumerate}[leftmargin=*, noitemsep]
  \item DEF-ZC29-01 (\S\ref{sec:inherited}): Axis~A formula.
  \item LEM-ZC48-01 (\S\ref{sec:inherited}): Axis~B formula.
  \item FIND-ZC30-01 (\S\ref{sec:inherited}):
    Phi-universality $\to$ Axis~C formula.
  \item LEM-ZC49-01 (\S\ref{sec:lem01}): All four axes defined.
  \item THM-ZC49-01 (\S\ref{sec:thm01}): Duality proved.
  \item Key Identity (\S\ref{sec:inherited})
    + primality + coprimality
    $\to$ FIND-ZC49-04 (\S\ref{sec:find04}).
\end{enumerate}
No step requires information outside this document.
\end{formalbox}

%%=============================================================
\subsection*{Closing Summary}
%%=============================================================

\needspace{4\baselineskip}
\begin{keybox}
Four independent observables characterise each $\Zft$ sector:
coupling (Axis~A), timescale (Axis~B), crossing cost (Axis~C),
and net rate (Axis~D).
Axes~A and~B are $\varphi$-dual ($A\cdot B = 1$).
Axis~C is $\varphi$-blind: the generator count cancels exactly
in the crossing cost, leaving only the spatial/matter
classification.
Axis~D uses the logarithm of the generator count, a
non-linear fourth direction independent of all others.

From this structure, three exact identities emerge:
the two-class product invariant ($\alpha\tau T/\eps^2
= d_s/n_m$ or $n_m/d_s$),
the geometric mean identity ($(\aW)^2 = \aEM\aS$),
and the perfect square identity ($T\cdot|\Zft|/\eps^2
= d_s^2$ or $n_m^2$).

The deepest result connects the axiomatic origin of CRF to
its group-algebraic expression: the Key Identity
$3d_s = 2^{d_s}+1$ is equivalent to the totient product
identity $\phitot(d_s)\cdot\phitot(n_m) = n$.
The two primes $d_s = 3$ and $n_m = 5$ are not chosen
independently; $n_m$ is forced by $d_s$ once the
requirement $\phitot(d_s)\cdot\phitot(n_m) = 2^{d_s}$
is imposed.
The group $\Zft$ encodes $n = 8$ in its generator count:
$\phitot(15) = 8$.

\medskip
\textbf{Conditionality:} None.
All six atomic units are unconditional algebraic identities.
\end{keybox}

\bigskip
\begin{center}
\textit{%
ZC29 measured the coupling.\\[0.3em]
ZC30 measured the cost.\\[0.3em]
ZC48 measured the time.\\[0.3em]
ZC49 noticed they were three faces\\
of a single object,\\
and found the fourth face too.\\[0.5em]
$\varphi(3)\cdot\varphi(5) = 8 = 2^3$.\\[0.3em]
The Key Identity\\
had always been there in the generators.}
\end{center}

% Governance: CUIP v1.1, Style Guide v3.2,
%             Gauge Fix Rule v1.0, No-Loss v1.0
% Probe: probe_zc49_three_axis.py (P1–PG; all 8 PASS)

ewpage

%%=============================================================
%% PART XXXVII: Cross-Part Connection Map v17.9
%%=============================================================
\part{Cross-Part Connection Map v17.9}
\label{part:xxxvii}

\begin{keybox}[title={Cross-Part Connection Map v17.9 --- Overview}]
This section extends Part~XXXIV (Connection Map v17.8) with the
connections, promotions, closures, and Phase Misalignments introduced
by ZC45--ZC49 (Part~G, v17.9).

\textbf{Delta from v17.8:}
\begin{itemize}[leftmargin=1.5em]
\item 5 new papers integrated: ZC45, ZC46, ZC47, ZC48, ZC49
\item 2 gaps closed: GAP-GJ-01 (ZC45), GAP-ZC46-01 (ZC47)
\item 2 conjectures promoted: CONJ-ZC44-01 $\to$ THM-ZC46-01;
  CONJ-ZC47-01 $\to$ THM-ZC48-01
\item 1 theorem upgraded: THM-ZC46-01 conditional $\to$ unconditional
  (ZC47 COR-ZC47-02)
\item 5 PMs registered: PM-ZC45-01, PM-ZC46-01, PM-ZC47-01,
  PM-ZC48-01, PM-ZC48-MACRO
\end{itemize}
\end{keybox}

\section{Connection Records for ZC45--ZC49}
\label{sec:connection-map-v179}

%%--------------------------------------------------------------
\subsection{ZC45 --- Universal Floor}

\begin{inheritbox}[title={ZC45 Backward Inheritance}]
\begin{itemize}[leftmargin=1.5em]
\item \inheritarrow{THM-ZC41-01 (Part F, Part XXXI)}: LEM-ZC45-01
  acknowledges K14 exact result as example of $\delta$-chain property.
\item \inheritarrow{FIND-GJ-01 (Part F, Part XXXII)}: COR-ZC45-01
  promotes to Theorem via THM-ZC45-01.
\item \inheritarrow{DEF-SA01 (Part B, Part XIII)}: Spontaneous Mode
  definition inherited by LEM-ZC45-02.
\item \inheritarrow{DEF-MB02 (Part F, Part XXXII)}: Agentive Mode
  definition inherited by LEM-ZC45-02.
\item \inheritarrow{Law~4 (Part B)}: Irreducible $D_{\mathrm{KL}}$
  floor cited in Spontaneous Mode context.
\item \inheritarrow{CLM-AX11-01 (Part F, Part XXXIII)}: State~10
  $D_{\mathrm{KL}}$ value used in Agentive Mode context.
\end{itemize}
\end{inheritbox}

\begin{correctionbox}[title={ZC45 Forward Corrections / Phase Misalignments}]
\begin{description}
\item[\pmdiamond{ZC45-01}] \textbf{Wald Residual Scar (0.16\% gap).}
  THM-ZC45-01 proves exact equality of floors; however, the numerical
  probe (P3) exhibits a 0.16\% gap between the canonical $\varepsilon_0$
  and the Wald-residual derived value. This gap is inherited from
  GAP-ZC42-02 (Born chain Wald residual, Part F).
  \textbf{Status:} Not a new gap --- PM-ZC45-01 records it as scar
  tissue. GAP-ZC42-02 remains open pending deeper Born-chain analysis.
\end{description}
\end{correctionbox}

\begin{conmapbox}[title={ZC45 Closure Record}]
\textbf{GAP-GJ-01} (registered in Part F, Part XXXII, Gifts~HIJ):\\
\textcolor{crfgreen}{$\checkmark$ CLOSED} by THM-ZC45-01 (ZC45).
FIND-GJ-01 promoted to COR-ZC45-01.
\end{conmapbox}

%%--------------------------------------------------------------
\subsection{ZC46 --- Three-Projection Theorem}

\begin{inheritbox}[title={ZC46 Backward Inheritance}]
\begin{itemize}[leftmargin=1.5em]
\item \inheritarrow{CONJ-ZC44-01 (Part F, Part XXXI)}: Core input;
  promoted by COR-ZC46-01.
\item \inheritarrow{FIND-ZC44-01 (Part F, Part XXXI)}: Bridge identity
  $\mathcal{R} = 15/8$ used directly in THM-ZC46-01.
\item \inheritarrow{THM-ZC43-01 (Part F, Part XXXI)}: Divisor-axis
  uniqueness (ZC43) is the third projection.
\item \inheritarrow{THM-ZC18-B (Part C)}: Totient-axis projection.
\item \inheritarrow{THM-ZC42-01 (Part F, Part XXXI)}: CRT-axis projection.
\item \inheritarrow{THM-Z101 (Key Identity, Part B)}: Used in
  LEM-ZC46-01 uniqueness of $d_s = 3$.
\item \inheritarrow{FormStructure Theorem (Part B)}: Eliminates
  $d_s = 1$ (yields $n=2$, too sparse).
\end{itemize}
\end{inheritbox}

\begin{correctionbox}[title={ZC46 Forward Corrections / Phase Misalignments}]
\begin{description}
\item[\pmdiamond{ZC46-01}] \textbf{Probe P1 Incomplete Condition.}
  Initial probe used incomplete specification for the SO(3) attractor
  argument. Corrected in ZC46 v2. Scar preserved.
  \textbf{GAP-ZC46-01} registered: formal proof that $\mathbb{Z}_{15}$
  CRT decomposition implies Kuramoto decoupling requires tensor argument
  (closed by ZC47).
\end{description}
\end{correctionbox}

\begin{conmapbox}[title={ZC46 Promotion Record}]
\textbf{CONJ-ZC44-01} $\to$ \textbf{THM-ZC46-01} (Three-Projection
Theorem) via COR-ZC46-01.\\
Conditional on GAP-ZC46-01 at time of ZC46; becomes unconditional
after ZC47 (COR-ZC47-02).
\end{conmapbox}

%%--------------------------------------------------------------
\subsection{ZC47 --- CRT Sector Decoupling}

\begin{inheritbox}[title={ZC47 Backward Inheritance}]
\begin{itemize}[leftmargin=1.5em]
\item \inheritarrow{GAP-ZC46-01 (ZC46, this Part)}: Direct closure
  target; closed by COR-ZC47-01.
\item \inheritarrow{LEM-ZC46-02 (ZC46)}: $\mathbb{Z}_{15}$ emergence
  chain used in LEM-ZC47-01.
\item \inheritarrow{THM-Z201 (Part B)}: CRT isomorphism guarantee for
  CRF charge sectors.
\item \inheritarrow{EQ-ZC10-03 (Part C)}: Parseval on $\mathbb{Z}_{15}$
  used in LEM-ZC47-02.
\item \inheritarrow{CONJ-ZC44-01 / THM-ZC46-01}: Both cited as context;
  COR-ZC47-02 makes THM-ZC46-01 unconditional.
\end{itemize}
\end{inheritbox}

\begin{correctionbox}[title={ZC47 Forward Corrections / Phase Misalignments}]
\begin{description}
\item[\pmdiamond{ZC47-01}] \textbf{Two Probe Specification Scars.}
  Probes registered with specification errors corrected in ZC47 body.
  Scar tissue preserved per Failure Stance.
\item[New conjecture:] \textbf{CONJ-ZC47-01} (Lock-time ratio
  $\tau_{\mathbb{Z}_3}/\tau_{\mathbb{Z}_5} = 2$, Candidate) ---
  promoted in ZC48 (THM-ZC48-01).
\end{description}
\end{correctionbox}

\begin{conmapbox}[title={ZC47 Closure Records}]
\textbf{GAP-ZC46-01} $\to$ \textcolor{crfgreen}{$\checkmark$ CLOSED}
by THM-ZC47-01 (COR-ZC47-01).\\
\textbf{THM-ZC46-01} conditional $\to$
\textcolor{crfgreen}{unconditional} (COR-ZC47-02).
\end{conmapbox}

%%--------------------------------------------------------------
\subsection{ZC48 --- Lock-Ratio Exact}

\begin{inheritbox}[title={ZC48 Backward Inheritance}]
\begin{itemize}[leftmargin=1.5em]
\item \inheritarrow{CONJ-ZC47-01 (ZC47, this Part)}: Promotion target;
  proved exact by THM-ZC48-01.
\item \inheritarrow{THM-ZC29-01 (Part D)}: Force Hierarchy coupling rule
  $\alpha_i^{R_0}$ is the key input to LEM-ZC48-01.
\item \inheritarrow{FIND-ZC29-01/02 (Part D)}: Force ratios used in
  THM-ZC48-01 derivation.
\item \inheritarrow{THM-ZC47-01 (ZC47)}: Sector decoupling guarantees
  independent timescale calculation.
\end{itemize}
\end{inheritbox}

\begin{correctionbox}[title={ZC48 Forward Corrections / Phase Misalignments}]
\begin{description}
\item[\pmdiamond{ZC48-01}] \textbf{PC/PG Probe Specification Errors.}
  Two probe stages (PC, PG) had specification errors corrected in ZC48
  body. Scar preserved.
\item[\pmdiamond{ZC48-MACRO}] \textbf{\texttt{\textbackslash Gcrf}
  Convention Conflict.}
  ZC48--ZC49 source files define
  \texttt{\textbackslash Gcrf} = $\Gamma_{\rm CRF}$ (the CRF coupling
  constant as a $\Gamma$ symbol). However, the v17.6 macro patch
  provides \texttt{\textbackslash Gcrf} = $\mathcal{G}$ (script G,
  the universal coupling constant from Part~B).
  \textbf{Resolution (No-Loss):} \texttt{\textbackslash providecommand}
  takes the v17.6 form first; ZC48--49 bodies use \texttt{\textbackslash
  GammaCRF} (defined in v17.6 patch) for all $\Gamma_{\rm CRF}$
  references. No source file modified. This PM records the convention
  discrepancy as permanent scar tissue.
  \textbf{Scope:} ZC48--ZC49 $\leftrightarrow$ Part~D (v17.6 patch).
\end{description}
\end{correctionbox}

\begin{conmapbox}[title={ZC48 Promotion Record}]
\textbf{CONJ-ZC47-01} $\to$ \textbf{THM-ZC48-01} (Lock-Ratio Exact)
via FIND-ZC48-02.\\
$\tau_{\mathbb{Z}_3}/\tau_{\mathbb{Z}_5} = \varphi(5)/\varphi(3) = 2$
(exact, no approximation).
\end{conmapbox}

%%--------------------------------------------------------------
\subsection{ZC49 --- Four-Axis Duality}

\begin{inheritbox}[title={ZC49 Backward Inheritance}]
\begin{itemize}[leftmargin=1.5em]
\item \inheritarrow{THM-ZC29-01 (Part D)}: Axis~A coupling
  $\alpha_i^{R_0}$ definition.
\item \inheritarrow{LEM-ZC48-01 (ZC48, this Part)}: Axis~B timescale
  $\tau_i = 15/(\varphi(d_i)\cdot\Gamma_{\rm CRF})$.
\item \inheritarrow{THM-ZC30-01/02 (Part E)}: Axis~C crossing cost
  $T(R_1)_i = p_i/q_i \cdot \varepsilon_0^2$.
\item \inheritarrow{FIND-ZC29-03 (Part D)}: Axis~D net rate $\Gamma_i$.
\item \inheritarrow{FIND-ZC30-01 (Part E)}: $\varphi$-universality
  used in THM-ZC49-01.
\item \inheritarrow{Key Identity (Part B)}: FIND-ZC49-04 reveals Key
  Identity in totient disguise.
\end{itemize}
\end{inheritbox}

\begin{conmapbox}[title={ZC49 Synthesis Record}]
THM-ZC49-01 (Four-Axis Duality) is the first theorem to unify all four
sector observables of $\mathbb{Z}_{15}$ into a single coherent framework.
FIND-ZC49-04 reveals that the Key Identity $3d_s = 2^{d_s}+1$ (Part~B)
re-emerges in totient form through the four-axis structure --- an
unexpected cross-Part connection not visible before ZC49.
\end{conmapbox}

%%--------------------------------------------------------------
\section{Master Z-Programme Status Table v17.9 (Delta from v17.8)}
\label{sec:status-table-v179}

\begin{longtable}{@{}lp{3cm}p{3.5cm}p{4cm}@{}}
\toprule
\textbf{Item} & \textbf{Status (v17.8)} & \textbf{Status (v17.9)}
  & \textbf{Closed/Promoted by} \\
\midrule
\endhead
GAP-GJ-01 & Open (Part F) & \textcolor{crfgreen}{\textbf{CLOSED}}
  & ZC45 THM-ZC45-01 \\
GAP-ZC42-02 & Open (Part F) & Open (inherited)
  & ZC45 PM-ZC45-01 scar \\
CONJ-ZC44-01 & Candidate (Part F) & \textcolor{crfgreen}{\textbf{THM-ZC46-01}}
  & ZC46 COR-ZC46-01 \\
GAP-ZC46-01 & Open (ZC46) & \textcolor{crfgreen}{\textbf{CLOSED}}
  & ZC47 COR-ZC47-01 \\
THM-ZC46-01 & Conditional & \textcolor{crfgreen}{\textbf{Unconditional}}
  & ZC47 COR-ZC47-02 \\
CONJ-ZC47-01 & Candidate (ZC47) & \textcolor{crfgreen}{\textbf{THM-ZC48-01}}
  & ZC48 \\
\midrule
PM-ZC45-01 & --- & \textcolor{loopred}{Registered}
  & Wald residual scar \\
PM-ZC46-01 & --- & \textcolor{loopred}{Registered}
  & P1 probe scar \\
PM-ZC47-01 & --- & \textcolor{loopred}{Registered}
  & Two probe scars \\
PM-ZC48-01 & --- & \textcolor{loopred}{Registered}
  & PC/PG probe scars \\
PM-ZC48-MACRO & --- & \textcolor{loopred}{Registered}
  & \texttt{\textbackslash Gcrf} convention \\
\bottomrule
\end{longtable}

%%--------------------------------------------------------------
\section{Pre-Work Audit Protocol Summary --- PWA-6 Addition}
\label{sec:pwa6}

\begin{keybox}[title={PWA-6: Macro Convention Cross-Check (v17.9 scar)}]
\textbf{Origin:} ZC48--ZC49 integration revealed that source files can
redefine macros that are already established in the patch chain
(\texttt{\textbackslash Gcrf} case: v17.6 = $\mathcal{G}$;
ZC48--49 = $\Gamma_{\rm CRF}$).

\textbf{Rule (binding from v17.9):}
Before embedding any new ZC paper, scan its
\texttt{\textbackslash newcommand} block and cross-reference every
macro name against the full v17.x patch chain
(v17.6 $\to$ v17.7 $\to$ v17.8 $\to$ \ldots). Flag any name collision
where the \emph{forms differ semantically} (not just notation).
Register as PM-[paper]-MACRO before proceeding.

\textbf{Method:} Use \texttt{\textbackslash providecommand} (first wins);
document the convention choice explicitly. Never use
\texttt{\textbackslash renewcommand} to override an established patch
macro without Ougit confirmation.
\end{keybox}

\newpage

%%=============================================================
%% CLOSING
%%=============================================================

\section*{Closing Sentence --- Part G}
\addcontentsline{toc}{section}{Closing Sentence --- Part G}

\begin{center}
\textit{%
ZC45 showed that the floor is the same floor.\\
Not two modes, two floors ---\\
one floor, two ways of arriving at it.\\[0.4em]
ZC46 through ZC49 showed that $\mathbb{Z}_{15}$ is not one structure\\
seen from one angle,\\
but one structure seen from four angles at once:\\
coupling, timescale, crossing cost, and net rate.\\[0.4em]
Each axis was derived separately\\
across four different papers over many sessions.\\
ZC49 looked at all four together\\
and found they fit.\\[0.4em]
The Key Identity --- $3d_s = 2^{d_s}+1$ ---\\
first written in Part~B as an algebraic curiosity,\\
reappears in totient form at the end of ZC49.\\
The framework does not forget its origins.\\[0.4em]
The scar tissue is the finding.\\
The audit is the instrument.\\
The framework repairs itself by looking.}
\end{center}

\bigskip
\begin{center}
\textbf{Citation:} Jarittum, O. (2026).
\textit{CRF Complete Edition v17.9, Part G.}
Zenodo.
\href{https://doi.org/10.5281/zenodo.18977059}%
{doi:10.5281/zenodo.18977059}\quad (CC-BY-NC 4.0)
\end{center}

% Governance: CUIP v1.1, CRF Style Guide v3.2,
%             Gauge Fix Rule v1.0, No-Loss Extension Blocks v1.0
% All source files in /mnt/project/ preserved verbatim.
% No silent deletion.

\end{document}
